% Options for packages loaded elsewhere
% Options for packages loaded elsewhere
\PassOptionsToPackage{unicode}{hyperref}
\PassOptionsToPackage{hyphens}{url}
\PassOptionsToPackage{dvipsnames,svgnames,x11names}{xcolor}
%
\documentclass[
  letterpaper,
  DIV=11,
  numbers=noendperiod]{scrreprt}
\usepackage{xcolor}
\usepackage{amsmath,amssymb}
\setcounter{secnumdepth}{1}
\usepackage{iftex}
\ifPDFTeX
  \usepackage[T1]{fontenc}
  \usepackage[utf8]{inputenc}
  \usepackage{textcomp} % provide euro and other symbols
\else % if luatex or xetex
  \usepackage{unicode-math} % this also loads fontspec
  \defaultfontfeatures{Scale=MatchLowercase}
  \defaultfontfeatures[\rmfamily]{Ligatures=TeX,Scale=1}
\fi
\usepackage{lmodern}
\ifPDFTeX\else
  % xetex/luatex font selection
\fi
% Use upquote if available, for straight quotes in verbatim environments
\IfFileExists{upquote.sty}{\usepackage{upquote}}{}
\IfFileExists{microtype.sty}{% use microtype if available
  \usepackage[]{microtype}
  \UseMicrotypeSet[protrusion]{basicmath} % disable protrusion for tt fonts
}{}
\makeatletter
\@ifundefined{KOMAClassName}{% if non-KOMA class
  \IfFileExists{parskip.sty}{%
    \usepackage{parskip}
  }{% else
    \setlength{\parindent}{0pt}
    \setlength{\parskip}{6pt plus 2pt minus 1pt}}
}{% if KOMA class
  \KOMAoptions{parskip=half}}
\makeatother
% Make \paragraph and \subparagraph free-standing
\makeatletter
\ifx\paragraph\undefined\else
  \let\oldparagraph\paragraph
  \renewcommand{\paragraph}{
    \@ifstar
      \xxxParagraphStar
      \xxxParagraphNoStar
  }
  \newcommand{\xxxParagraphStar}[1]{\oldparagraph*{#1}\mbox{}}
  \newcommand{\xxxParagraphNoStar}[1]{\oldparagraph{#1}\mbox{}}
\fi
\ifx\subparagraph\undefined\else
  \let\oldsubparagraph\subparagraph
  \renewcommand{\subparagraph}{
    \@ifstar
      \xxxSubParagraphStar
      \xxxSubParagraphNoStar
  }
  \newcommand{\xxxSubParagraphStar}[1]{\oldsubparagraph*{#1}\mbox{}}
  \newcommand{\xxxSubParagraphNoStar}[1]{\oldsubparagraph{#1}\mbox{}}
\fi
\makeatother


\usepackage{longtable,booktabs,array}
\newcounter{none} % for unnumbered tables
\usepackage{calc} % for calculating minipage widths
% Correct order of tables after \paragraph or \subparagraph
\usepackage{etoolbox}
\makeatletter
\patchcmd\longtable{\par}{\if@noskipsec\mbox{}\fi\par}{}{}
\makeatother
% Allow footnotes in longtable head/foot
\IfFileExists{footnotehyper.sty}{\usepackage{footnotehyper}}{\usepackage{footnote}}
\makesavenoteenv{longtable}
\usepackage{graphicx}
\makeatletter
\newsavebox\pandoc@box
\newcommand*\pandocbounded[1]{% scales image to fit in text height/width
  \sbox\pandoc@box{#1}%
  \Gscale@div\@tempa{\textheight}{\dimexpr\ht\pandoc@box+\dp\pandoc@box\relax}%
  \Gscale@div\@tempb{\linewidth}{\wd\pandoc@box}%
  \ifdim\@tempb\p@<\@tempa\p@\let\@tempa\@tempb\fi% select the smaller of both
  \ifdim\@tempa\p@<\p@\scalebox{\@tempa}{\usebox\pandoc@box}%
  \else\usebox{\pandoc@box}%
  \fi%
}
% Set default figure placement to htbp
\def\fps@figure{htbp}
\makeatother





\setlength{\emergencystretch}{3em} % prevent overfull lines

\providecommand{\tightlist}{%
  \setlength{\itemsep}{0pt}\setlength{\parskip}{0pt}}



 


\KOMAoption{captions}{tableheading}
\makeatletter
\@ifpackageloaded{bookmark}{}{\usepackage{bookmark}}
\makeatother
\makeatletter
\@ifpackageloaded{caption}{}{\usepackage{caption}}
\AtBeginDocument{%
\ifdefined\contentsname
  \renewcommand*\contentsname{Table of contents}
\else
  \newcommand\contentsname{Table of contents}
\fi
\ifdefined\listfigurename
  \renewcommand*\listfigurename{List of Figures}
\else
  \newcommand\listfigurename{List of Figures}
\fi
\ifdefined\listtablename
  \renewcommand*\listtablename{List of Tables}
\else
  \newcommand\listtablename{List of Tables}
\fi
\ifdefined\figurename
  \renewcommand*\figurename{Figure}
\else
  \newcommand\figurename{Figure}
\fi
\ifdefined\tablename
  \renewcommand*\tablename{Table}
\else
  \newcommand\tablename{Table}
\fi
}
\@ifpackageloaded{float}{}{\usepackage{float}}
\floatstyle{ruled}
\@ifundefined{c@chapter}{\newfloat{codelisting}{h}{lop}}{\newfloat{codelisting}{h}{lop}[chapter]}
\floatname{codelisting}{Listing}
\newcommand*\listoflistings{\listof{codelisting}{List of Listings}}
\makeatother
\makeatletter
\makeatother
\makeatletter
\@ifpackageloaded{caption}{}{\usepackage{caption}}
\@ifpackageloaded{subcaption}{}{\usepackage{subcaption}}
\makeatother
\usepackage{bookmark}
\IfFileExists{xurl.sty}{\usepackage{xurl}}{} % add URL line breaks if available
\urlstyle{same}
\hypersetup{
  pdftitle={UMRS Restoration Features},
  colorlinks=true,
  linkcolor={blue},
  filecolor={Maroon},
  citecolor={Blue},
  urlcolor={Blue},
  pdfcreator={LaTeX via pandoc}}


\title{UMRS Restoration Features}
\author{}
\date{2027-09-09}
\begin{document}
\maketitle

\renewcommand*\contentsname{Table of contents}
{
\setcounter{tocdepth}{2}
\tableofcontents
}

\bookmarksetup{startatroot}

\chapter*{Preface}\label{preface}
\addcontentsline{toc}{chapter}{Preface}

\markboth{Preface}{Preface}

This restoration features book documents the key design elements of
ecosystem restoration features built on the Upper Mississippi River
System (UMRS). This data model has been developed to represent the wide
range of restoration practices employed on the UMRS throughout many
years and be adaptable into the future. As we learn lessons about
restoration, this data model has evolved along with our best practice.

\section*{Authors}\label{authors}
\addcontentsline{toc}{section}{Authors}

\markright{Authors}

\begin{itemize}
\tightlist
\item
  Michael Dougherty, Geographer, US Army Corps of Engineers, Rock Island
  District
\item
  Bobby Oliver, Geographer, US Army Corps of Engineers, Rock Island
  District
\end{itemize}

\section*{Abstract}\label{abstract}
\addcontentsline{toc}{section}{Abstract}

\markright{Abstract}

\bookmarksetup{startatroot}

\chapter{Introduction}\label{introduction}

\section{What are Restoration
Features?}\label{what-are-restoration-features}

Restoration features are what we build.

\section{UMRS Restoration Data Model}\label{umrs-restoration-data-model}

\subsection{Purpose}\label{purpose}

The restoration features are GIS feature classes designed to document
the key design elements of ecosystem restoration features built on the
Upper Mississippi River System (UMRS). This data model has been
developed to represent the wide range of restoration practices employed
on the UMRS throughout many years and be adaptable into the future. As
we learn lessons about restoration, this data model has evolved along
with our best practice.

\subsection{Data Design Standards}\label{data-design-standards}

The entities in this data model represent the physical structures
constructed to accomplish the restoration objectives of each project.
These entities are grouped into standardized GIS ``feature classes''
(groups of objects of single type having a common set of attributes).
This standard was adapted from the DOD SDSFIE-V (https://www.sdsfie.mil)
to represent the specialized features built on the UMRS. This data model
was designed to represent UMRS best practice as documented in the UMRR
Design Handbook (2012, 2025).

\subsection{Data Use}\label{data-use}

\begin{itemize}
\tightlist
\item
  Common Terms - Provides a common lexicon of ecosystem restoration
  features representing current best practice to streamline
  communication across a diverse group of stakeholders, spread across
  many organizations, from many disciplinary backgrounds.
\item
  Cartography - Defines a standardized framework for representing a wide
  variety of complex spatial features on different media types (i.e.,
  static paper maps, interactive web maps). Maps are one of the most
  effective communication tools to describe project benefits to the
  public.
\item
  Upward Reporting - Expressing these diverse features in a standard
  framework supports automated summarization to streamline routine
  upward reporting of project status and benefits.
\item
  Record Keeping - Forty years of UMRS restoration has generated a rich
  set of built features to monitor and learn from. This standardized
  framework, that can evolve with current best practice, is the best way
  to track our investments.
\item
  Whole Project Lifecycle - This standardized framework allows us to
  track projects throughout their entire lifecycle, from planning and
  feasibility, through design and construction, and through operation
  and maintenance. UMRS Restoration Data Model
\end{itemize}

\section{Implementation}\label{implementation}

The Army Corps of Engineers project process typically involves four main
phases: pre-study initiation (community request), feasibility study,
design and construction, and operation and maintenance.

\subsection{Feasibility}\label{feasibility}

\begin{itemize}
\tightlist
\item
  The Corps conducts a detailed study to assess the problem, evaluate
  potential solutions, and determine project feasibility.
\item
  Analysis includes environmental impacts, economic benefits, and
  cost-benefit analysis.
\item
  Report is submitted to higher authorities for review and approval.
\end{itemize}

\subsection{Design \& Construction}\label{design-construction}

\begin{itemize}
\tightlist
\item
  If the project is approved, detailed design plans are developed.
\item
  Necessary permits are obtained.
\item
  Construction contracts are awarded and project is built.
\end{itemize}

\subsubsection{Operation and
Maintenance}\label{operation-and-maintenance}

Once completed, the Corps manages the project, including routine
maintenance and repairs.

\section{Current Programs Usage}\label{current-programs-usage}

\subsection{UMRR}\label{umrr}

To work within a partnership among federal and state agencies and other
organizations; to construct high-performing habitat restoration,
rehabilitation projects; to produce state-of-the-art knowledge through
monitoring, research, and assessment; to engage other organizations to
accomplish the Upper Mississippi River Restoration Program's vision.

The HREP element has fostered interdisciplinary and collaborative
planning for habitat restoration, preservation, and enhancement
previously non-existent on any other river system in the United States.

\subsection{NESP}\label{nesp}

The Navigation and Ecosystem Sustainability Program (NESP) is a
long-term, dual-purpose program that integrates navigation improvements
and ecosystem restoration together to provide Upper Mississippi River
System once in a generation-type positive impact.

The purpose of the Navigation and Ecosystem Sustainability Program is to
improve efficiency and capacity of the nationally significant Upper
Mississippi River -- Illinois Waterway (UMR-IWW) navigation system while
protecting preserving and enhancing the structure, diversity, and
function of this nationally significant ecosystem.

\subsection{Notes}\label{notes}

Several ecosystem restoration programs are active on the UMRS at the
local, state, and federal levels. This data model was designed to
function as a repository of the ecosystem restoration features common
within a large river floodplain system such as the UMRS, independent of
jurisdiction or funding source.

\bookmarksetup{startatroot}

\chapter{Case Studies}\label{case-studies}

\subsection{\texorpdfstring{\textbf{An Integrated Framework in Practice:
Detailed Case Studies in USACE Upper Mississippi River
Restoration}}{An Integrated Framework in Practice: Detailed Case Studies in USACE Upper Mississippi River Restoration}}\label{an-integrated-framework-in-practice-detailed-case-studies-in-usace-upper-mississippi-river-restoration}

\subsubsection{\texorpdfstring{\textbf{Introduction: From a Toolkit of
Features to Integrated
Systems}}{Introduction: From a Toolkit of Features to Integrated Systems}}\label{introduction-from-a-toolkit-of-features-to-integrated-systems}

The ecological restoration of the Upper Mississippi River System (UMRS)
is a monumental undertaking guided by the sophisticated methodology
outlined in the U.S. Army Corps of Engineers' (USACE) \emph{2012 UMRR
EMP Environmental Design Handbook}. While the handbook provides a
detailed lexicon of individual ``restoration features,'' their true
power is only realized when they are integrated into a cohesive project
designed to address specific, large-scale ecological deficits. To
understand this methodology in its practical application, it is best to
examine it through the lens of real-world case studies. This chapter
will therefore explore the USACE restoration process through four major
Habitat Rehabilitation and Enhancement Projects (HREPs). Each section
will deconstruct a project to show how a unique combination of primary
and supporting features was used to address its specific ecological
problems, demonstrating how the abstract toolkit is transformed into
functional, resilient, and interconnected ecosystems on the ground.

\begin{center}\rule{0.5\linewidth}{0.5pt}\end{center}

\subsection{\texorpdfstring{\textbf{Case Study 1: The Pool 8 Islands
HREP - Geomorphic Restoration for Ecological
Resurgence}}{Case Study 1: The Pool 8 Islands HREP - Geomorphic Restoration for Ecological Resurgence}}\label{case-study-1-the-pool-8-islands-hrep---geomorphic-restoration-for-ecological-resurgence}

\textbf{Project Context and Goal:} The Pool 8 Islands HREP, located in
the vast Weaver Bottoms area near Brownsville, Minnesota, addressed one
of the most significant consequences of river impoundment: the creation
of a wide, shallow, lake-like environment. The 4,000-acre backwater
suffered from a wind fetch of over two miles, which generated constant
wave action. This energy re-suspended fine bottom sediments, leading to
extremely high turbidity that blocked sunlight, prevented the growth of
vital submerged aquatic vegetation (SAV), and eroded the few remaining
natural shorelines. The primary goal was to re-engineer the
geomorphology of the pool to reduce wave energy, improve water clarity,
and spur the large-scale recovery of the aquatic ecosystem.

\textbf{Integrated Feature Application in Detail:} This project is a
classic example of geomorphic restoration centered on the beneficial use
of dredged material. The foundational activity was a large-scale
\textbf{Dredge Event}, which sourced over 1.5 million cubic yards of
clean sand during routine maintenance of the 9-foot navigation channel.
This material was then transported and strategically placed to construct
a complex of 17 new \textbf{Islands}. The curvilinear, naturalistic
layout of these islands was precisely guided by a detailed
\textbf{Island Line} survey plan, which provided the exact coordinates
for construction crews to place material accurately and create the
desired footprint.

The immediate stability of these new sand landforms was critical. To
prevent their immediate erosion by the very forces they were designed to
block, their exposed, windward shores were armored with a heavy rock
\textbf{Revetment}. This involved placing thousands of tons of Class III
riprap over a durable geotextile fabric, creating a resilient shield
against wave attack that was essential for the project's long-term
survival.

Once the physical forms were stable, a massive \textbf{Flora Planting or
Seeding} effort transformed the sterile sand into living habitat. On the
highest elevations of the islands, designed to be above the 2-year flood
frequency, over 80,000 hardwood trees---including species like swamp
white oak, river birch, and pecan---were planted to initiate a
floodplain forest. On the lower, marshy benches and shorelines, the
substrate was seeded with native grasses and forbs to create a diverse
emergent \textbf{Wetland} fringe, providing both erosion control and
valuable habitat.

The project's success was profound. In the protected zones behind the
islands, water clarity improved from inches to several feet. This
triggered a massive resurgence of SAV, with coverage increasing by
hundreds of acres. This, in turn, created a rich nursery and foraging
habitat for fish and invertebrates, and a major feeding area for
migratory waterfowl. To manage the area and inform the public, numerous
\textbf{Signs} were installed. These include kiosks at public access
points explaining the project's success and ``Slow No Wake Zone'' buoys
to protect the fragile, newly established vegetation from boat wakes.

\begin{center}\rule{0.5\linewidth}{0.5pt}\end{center}

\subsection{\texorpdfstring{\textbf{Case Study 2: The Weaver Bottoms
(Pool 5) HREP - Large-Scale Drawdown and Geomorphic
Control}}{Case Study 2: The Weaver Bottoms (Pool 5) HREP - Large-Scale Drawdown and Geomorphic Control}}\label{case-study-2-the-weaver-bottoms-pool-5-hrep---large-scale-drawdown-and-geomorphic-control}

\textbf{Project Context and Goal:} The Weaver Bottoms project in Pool 5
was another massive HREP aimed at restoring a 4,000-acre backwater that
had degraded into a wide, shallow lake with almost no vegetation. Like
Pool 8, it suffered from immense wind and wave action. The project had a
dual goal: first, to reduce wave energy through physical barriers, and
second, to stimulate the germination of the dormant but viable native
seed bank that lay buried in the riverbed sediments, a process that
required physically dewatering a vast area of the Mississippi River.

\textbf{Integrated Feature Application in Detail:} This complex project
combined geomorphic construction with large-scale hydraulic
manipulation. Similar to Pool 8, a major \textbf{Dredge Event} sourced
material for the construction of several new \textbf{Islands}, including
a large, crescent-shaped island, to break up wave fetch and create
sheltered zones. However, the most innovative component was a
large-scale, partial \textbf{drawdown} of the backwater. This was a
significant engineering challenge that required modifying existing
\textbf{Water Control Structures} at the outlet of the backwater to
release water faster than it was entering. To facilitate the dewatering
of such a vast area, targeted \textbf{Dredging} was also performed to
create deeper \textbf{Flow Channels} that would act as drainage
arteries, ensuring the entire area could dewater effectively as the main
pool level was lowered.

The drawdown exposed thousands of acres of mudflats for an entire
growing season. This had two key ecological effects: first, it allowed
the wet, mucky sediment to dewater and consolidate, creating a firmer
substrate. Second, and more importantly, the combination of sunlight and
air stimulated the germination of a massive seed bank of native aquatic
plants (like wild celery and arrowhead) that had been dormant in the
sediment for years. When the area was slowly reflooded the following
year, a lush and diverse plant community emerged, transforming the
``dead lake'' back into a productive marsh. The newly constructed
islands, protected by \textbf{Revetments}, provided the long-term wave
protection necessary to ensure this newly established vegetation would
persist.

\begin{center}\rule{0.5\linewidth}{0.5pt}\end{center}

\subsection{\texorpdfstring{\textbf{Case Study 3: The Swan Lake HREP -
Active Management through Complex
Infrastructure}}{Case Study 3: The Swan Lake HREP - Active Management through Complex Infrastructure}}\label{case-study-3-the-swan-lake-hrep---active-management-through-complex-infrastructure}

\textbf{Project Context and Goal:} The Swan Lake HREP, a vast
10,000-acre floodplain complex on the Illinois River, is one of the most
ambitious examples of active habitat management. The site had become
disconnected from natural flood pulses, and its wetlands were degraded.
The primary goal was to restore a highly productive mosaic of wetland
habitats by giving resource managers the ability to control water levels
with absolute precision, completely independent of the unpredictable and
often ill-timed fluctuations of the river.

\textbf{Integrated Feature Application in Detail:} This project is a
masterclass in the integration of supporting infrastructure to achieve
an ecological goal. The project's core is a series of large
\textbf{Moist Soil Management Units (MSUs)}. This first required the
construction of miles of new \textbf{Embankments}, whose alignment was
precisely dictated by an \textbf{Embankment Centerline}. The design was
a major geotechnical challenge, requiring extensive soil analysis from
data collected at dozens of \textbf{Borehole Points} to ensure the
embankments would be stable on the soft, alluvial floodplain soils.

The ability to manage water is enabled by several key features. Large,
cast-in-place concrete \textbf{Water Control Structures}, fitted with
removable stoplogs, were built into the embankments. Because gravity
flow was insufficient, several large \textbf{Pump Stations} were
constructed, with a combined capacity of over 300 cubic feet per second.
This immense power gives managers the ability to fill the vast wetland
units from the river in the fall or, critically, to dewater them after
an unseasonal summer flood to save the growing food crop.

This operational infrastructure required an equally complex network of
utilities and access. A robust electrical grid was installed, consisting
of miles of buried armored cable (\textbf{Electrical Utility Segments})
that run along the levee crests. These segments connect the main power
grid to the \textbf{Electrical Utility Nodes}---the transformers and
motor control panels---located at each remote pump station. A water
conveyance network, composed of large-diameter HDPE pipes (\textbf{Water
Utility Segments}) and large, wheel-operated gate valves (\textbf{Water
Utility Nodes}), moves the massive volumes of water from the pumps to
the various wetland cells.

None of this would be possible to build, operate, or maintain without
physical access. The project included the construction of several miles
of \textbf{Road Segments} on the crests of the main embankments. These
roads were surfaced with a thick \textbf{Pavement Section} of compacted
aggregate, specifically designed to support the weight of heavy vehicles
like cranes and excavators needed for long-term maintenance. The entire
construction process was guided by a \textbf{Feature Stationing} system,
which provided a linear ``address'' for every component, allowing for
the precise placement of every culvert, pump, and water control
structure along the project's miles of embankments.

\begin{center}\rule{0.5\linewidth}{0.5pt}\end{center}

\subsection{\texorpdfstring{\textbf{Case Study 4: Restoring Connectivity
and Complexity - Peoria Lake, Pool 12, and Big Timber
HREPs}}{Case Study 4: Restoring Connectivity and Complexity - Peoria Lake, Pool 12, and Big Timber HREPs}}\label{case-study-4-restoring-connectivity-and-complexity---peoria-lake-pool-12-and-big-timber-hreps}

\textbf{Project Context and Goal:} Many HREPs are not about large-scale
geomorphic change but instead focus on more surgical interventions to
restore depth, flow, and habitat complexity in degraded backwaters and
side channels that have become choked with sediment and cut off from the
main river.

\textbf{Integrated Feature Application in Detail:} At the \textbf{Peoria
Lake HREP}, the foundational activity was a \textbf{Dredge Event} that
removed over 1.3 million cubic yards of sediment to excavate deep
channels in the sediment-filled lake. Because the dredged material was
fine-grained silt, it was unsuitable for island building and was instead
pumped to a designated \textbf{Dredge Placement} site---an adjacent
upland containment area where it could be safely dewatered. To bring the
newly dredged pools to life, a new \textbf{Flow Channel} was excavated
to connect them to the main river, providing a constant supply of fresh
water. Within these pools, \textbf{Fish Structures}, including rock
piles and anchored logs, were placed to provide immediate habitat
complexity for the local fishery.

In \textbf{Pool 12} of the Mississippi River, the focus was on restoring
flow to side channels cut off by historic \textbf{Riverine Flow
Structures} (rock wing dikes). Rather than removing the dikes, a more
targeted approach was taken. A \textbf{Notch}, typically 30-50 feet
wide, was cut into each dike. The design and placement of these notches
were guided by a \textbf{Riverine Flow Structure Line} survey plan to
ensure they captured the right amount of flow. This allowed a controlled
portion of the river's current to pass through, scouring the side
channel and restoring a healthy, flowing environment without
compromising the navigation function of the main channel.

Finally, at the \textbf{Big Timber HREP} in Pool 17, managers faced a
different problem: the water in the adjacent Mississippi River channel
was often too turbid and nutrient-laden to be desirable for their
managed wetland. To overcome this, the project included the installation
of a high-capacity \textbf{Well}. This feature drills deep into the
alluvial aquifer to provide a source of clean, clear, sediment-free
groundwater. The well gives managers a critical tool to fill the wetland
with high-quality water, allowing them to achieve their vegetation goals
even when the river itself is not a suitable water source.

\bookmarksetup{startatroot}

\chapter{Restoration Features
Summary}\label{restoration-features-summary}

\subsection{The Integrated Design of River Restoration
Features}\label{the-integrated-design-of-river-restoration-features}

The ``2012 UMRR EMP Environmental Design Handbook'' outlines a
comprehensive and integrated approach to river restoration, where a
diverse suite of features are designed to work together to achieve
large-scale ecological uplift in the Upper Mississippi River System. The
relationship between the individual restoration features and the
overarching design practices is not linear but rather a complex,
synergistic web. The core design practice is to view the river not as a
series of isolated problems but as a complete ecosystem. Therefore,
projects combine multiple features to mimic natural processes, restore
habitat, and ensure the long-term resilience of the restored
environment.

The restoration features can be broadly understood in two categories:
Primary Ecological Features and the Supporting Engineering and
Infrastructure Features that enable them. The design practice involves
layering these two categories to create a functional and sustainable
project.

\begin{enumerate}
\def\labelenumi{\arabic{enumi}.}
\tightlist
\item
  Primary Ecological Features: The Building Blocks of Restoration These
  features are the visible, direct interventions aimed at restoring
  ecological structure and function. Design practices for these features
  focus on creating specific habitats and restoring natural river
  processes.
\end{enumerate}

\begin{itemize}
\item
  Objective: Enhancing Aquatic Habitat \& Connectivity

  \begin{itemize}
  \item
    Features Used: Dredge Events are a cornerstone practice to restore
    depth in sediment-filled backwaters. This creates overwintering
    refugia for fish and improves flow. The design is not simply about
    digging a hole; it involves creating varied depths and complex
    shapes. To further enhance this, Fish Structures (like rock piles
    and log snags) are added to provide cover. Flow Channels and Notches
    in dikes are engineered to reintroduce vital, oxygenated water from
    the main channel into these restored backwaters, preventing
    stagnation.
  \item
    Relationship to Design: The design practice here is one of habitat
    layering. A single backwater dredging project is designed in concert
    with flow introduction and in-stream cover to create a complete,
    multi-faceted aquatic habitat, rather than just a deep pool.
  \end{itemize}
\item
  Objective: Restoring Landforms and Floodplain Dynamics

  \begin{itemize}
  \item
    Features Used: Islands, often built using material from Dredge
    Events (a key synergy), are a primary tool. Their design is
    strategic: they act as barriers to reduce wind and wave erosion in
    backwaters, creating sheltered zones for aquatic plants to grow. The
    islands themselves are designed with varied topography, using Flora
    Planting and Seeding to establish Floodplain Forests and provide
    secure nesting habitat.
  \item
    Relationship to Design: The design practice is geomorphic
    restoration. It's not just about building an island, but about using
    that island as a tool to reshape the river's flow, energy, and light
    conditions to benefit the wider ecosystem. The island's design
    (shape, elevation, vegetation) is dictated by these broader goals.
  \end{itemize}
\item
  Objective: Managing Riparian and Wetland Zones

  \begin{itemize}
  \item
    Features Used: Moist Soil Management Units (MSUs) and Wetlands are
    created or restored using Embankments to impound water. The key
    design practice is the precise management of water levels, which is
    accomplished using Water Control Structures, Pump Stations, and
    Wells. By controlling the timing and duration of flooding and
    drawdowns, managers can stimulate the growth of specific native
    plants that provide critical food for migratory waterfowl.
  \item
    Relationship to Design: This demonstrates the practice of active
    ecological management. The features are designed not as static
    elements, but as a dynamic, controllable system that allows managers
    to adapt to environmental conditions and meet specific ecological
    targets year after year.
  \end{itemize}
\end{itemize}

\begin{enumerate}
\def\labelenumi{\arabic{enumi}.}
\setcounter{enumi}{1}
\tightlist
\item
  Supporting Engineering and Infrastructure Features: The Foundation of
  Success These features are the critical engineering components that
  make the primary ecological restorations possible, stable, and
  accessible. The design practice is to integrate these robust
  engineering solutions seamlessly to ensure the project's longevity and
  functionality.
\end{enumerate}

\begin{itemize}
\item
  Geotechnical and Structural Stability: A Revetment (rock armor) is
  designed to protect a newly built Island from eroding. The design of a
  massive Embankment relies entirely on data from Borehole Points to
  understand the soil's capacity to support the structure. Riverine Flow
  Structures like wing dikes are designed based on extensive hydraulic
  modeling to ensure they steer river currents effectively without
  causing unintended erosion elsewhere.
\item
  Construction and Operational Logistics: The entire system is built
  upon a framework of civil infrastructure. Feature Stationing,
  Embankment Centerlines, and Island Lines provide the precise survey
  control needed for construction. Road Segments and Pavement Sections
  are designed to provide all-weather access for heavy machinery during
  construction and for long-term maintenance and monitoring.
\item
  Power and Utilities: The active management of an MSU is impossible
  without a reliable Pump Station, which in turn depends on a network of
  Electrical Utility Nodes and Segments. The design practice here
  ensures that this critical infrastructure is itself resilient, often
  elevated above flood levels to maintain operational capacity during
  high-water events.
\end{itemize}

\subsubsection{Conclusion}\label{conclusion}

In summary, the relationship between restoration features and design
practices is one of interdependence and synergy. The design process for
a UMRR EMP project begins with a clear ecological objective, such as
improving fish habitat or restoring a floodplain forest. To achieve
this, a suite of primary ecological features are selected and designed
as an interconnected system. Finally, this system is built upon a robust
foundation of supporting engineering and infrastructure features, which
are designed to ensure the primary features are constructed correctly,
remain stable over their design life, and can be actively managed to
maximize their ecological benefits. No single feature stands alone; its
design is always considered in the context of the entire project and the
broader river ecosystem.

\part{Restoration Features}

\chapter{Restoration Features}\label{restoration-features-1}

\section{Purpose}\label{purpose-1}

Restoration of ecosystems through a variety of features, including
wetlands, riparian areas, aquatic systems, and oyster reefs. The goal of
these restoration projects is to restore the natural function,
structure, and processes of degraded ecosystems.

\section{Notes}\label{notes-1}

Data Sources - If a project has been constructed, CAD drawings or
As-Built drawings should be used to guide feature digitizing. If a
project is in the planning stages, the schema can be used to record
potential features and track which planning alternatives they
participate in. Elevation data should be attached to geometries as
specified in this document.

These feature classes are stored in the MVR UMRR Oracle database to
ensure that authoritative data is available to all. See these reference
guides to access the oracle database for editing features.

Restoration features are a topical group of feature classes used to
represent, at various phases, engineering features constructed or
implemented to directly meet the goals and objectives of programs and
the projects they support.

\section{Restoration Features}\label{restoration-features-2}

\begin{itemize}
\tightlist
\item
  DredgingEvent (polygon) -- The extent of the area subject to dredging
  action. (area of material removal)
\item
  Embankment (polygon) -- The centerline of an embankment for
  controlling the waters of the sea, river, or other water bodies.
\item
  EmbankmentCenterline (line) -- The centerline of an embankment for
  controlling the waters of the sea, river, or other water bodies.
\item
  FishStructure (polygon) -- An engineered structure designed to improve
  fish habitat or allow improved fish movement through the river system.
\item
  FloraPlantingOrSeeding (polygon) -- A location where vegetation is
  planted or seeded.
\item
  FlowChannel (polygon) -- An excavated channel (via mechanical or
  hydraulic means) intended to allow for the flow of water across, into,
  or out of the project area or any subcomponent of a project area.
\item
  Island (polygon) -- Island structures engineered with a variety of
  techniques and materials that serve as barriers to flow, wind and
  waves.
\item
  IslandLine (line) -- Island structures engineered with a variety of
  techniques and materials that serve as barriers to flow, wind and
  waves.
\item
  MoistSoilMgmtUnits (polygon) -- Areas or units designed to enhance
  wetland habitat through the use of water containment, water supply and
  water control structures.
\item
  Notch (point) -- A notch in a riverine flow structure (typically a
  dike) used to control the flow of water into and out of the low flow
  area.
\item
  PumpStation (point) -- A location where a pump is located to pump
  water from the interior of a levee.
\item
  Revetment (polygon) -- An area faced with stone or other material
  along the shoreline designed to stabilize the bank or protect it from
  erosion.
\item
  RiverineFlowStructureLine (line) -- A natural or man-made flow (or
  sediment) control structure in a water course or water body such as a
  dike or weir. This feature should not be used to model a levee.
\item
  RiverineFlowStructure (polygon) -- A natural or man-made flow (or
  sediment) control structure in a water course or water body such as a
  dike or weir. This feature should not be used to model a levee.
\item
  WaterControlStructure (point) -- A structure designed to allow flow
  through an embankment.
\item
  Well (point) -- A shaft dug or drilled into the Earth for the purpose
  of extracting fluids from the subsurface, collecting environmental
  samples, injecting fluids into the subsurface or extracting
  contamination or other impurities from the subsurface.
\item
  Wetland (polygon) -- Geospatial feature composed of one unique wetland
  class (no complexes) used to determine USACE jurisdiction and
  regulatory impacts.
\end{itemize}

\chapter{Dredging Event}\label{dredging-event}

\section{Spatial \& Geometric
Definition}\label{spatial-geometric-definition}

\begin{itemize}
\item
  Geometry Type: Polygon (Dredge\_Event\_Area).
\item
  Spatial Delineation: Extents of authorized aquatic cut prisms within
  backwater lakes, secondary chutes, side channels, access cuts, and
  sediment deflection basins.
\item
  Vertical Bounds: Existing hydrographic bathymetric profile down to
  engineered project design grade, including allowable design overdepth
  tolerance (+0.5 to +1.0 ft).
\item
  Primary Restoration \& Engineering Objectives Overwintering Fish
  Refugia: Deepen degraded shallow backwaters to restore minimum safe
  thermal depths (≥8.0 to 10.0 ft) beneath winter river ice cover.
\item
  Riverine Velocity Dissipation: Provide low-velocity resting basins
  (\textless0.05 ft/s) protected from cold, turbulent main-stem drafts.
\item
  Aquatic Habitat Longevity: Calculate design dredge depths using
  long-term historic bathymetric sedimentation rates to guarantee a
  minimum 25- to 50-year biological lifespan before re-shoaling.
\end{itemize}

\subsection{Typical Cross-Sections, Dimensions \& Sizing
Criteria}\label{typical-cross-sections-dimensions-sizing-criteria}

\begin{itemize}
\item
  Bottom Widths: Primary fish cuts span 50 to 150 ft; deep winter holes
  span 100 to 300 ft.
\item
  Side Slopes: Typically specified at 1V:3H to 1V:5H in unconsolidated
  alluvial sands, and 1V:6H to 1V:10H in soft cohesive lacustrine silts.
\item
  Target Water Depths: Maintained at 8.0 to 12.0 ft below flat low
  navigation pool elevation.
\end{itemize}

\subsection{Construction \& Implementation
Methodology}\label{construction-implementation-methodology}

\begin{itemize}
\item
  Hydraulic Cutterhead Dredging: Slurry pumping (12 to 24 in pipelines)
  utilized when material is targeted for long-distance transport into
  confined disposal facilities (CDFs) or island containment berms.
\item
  Mechanical Clamshell / Excavator Dredging: Barge-mounted long-reach
  excavators and scows utilized when sediments are coarse sand or placed
  immediately along adjacent island cores.
\item
  Turbidity Controls: Silt curtains, hydraulic velocity dissipators, and
  continuous down-drift total suspended solids (TSS) and dissolved
  oxygen (DO) monitoring.
\end{itemize}

\subsection{Quality Control, Inspection \& Long-Term
Maintenance}\label{quality-control-inspection-long-term-maintenance}

\begin{itemize}
\item
  Pre- and Post-Dredge Hydrographic Surveys: Multi-beam and
  dual-frequency single-beam acoustic sonar runs tied to USACE River
  Gages and NAVD88 vertical datum.
\item
  Long-Term Monitoring: Bathymetric transects established at 200-ft
  intervals, surveyed at 1-, 3-, 5-, and 10-year intervals
  post-construction to quantify sediment accretion rates.
\end{itemize}

\section{Spatial Feature Attribute}\label{spatial-feature-attribute}

\subsection{Definition}\label{definition}

Backwater dredging typically consists of dredging channels with fingers
(dredged channels that extend out away from the main dredge cut). The
depth and size (length and width) of the dredge cut depends on several
site specific factors.

\subsection{Purpose}\label{purpose-2}

Sediment is dredged to create depth diversity in the backwaters can be
used to enhance aquatic areas with islands or terrestrial areas with
increase topographic diversity, which promotes the growth of mast trees.

\subsection{Applicability}\label{applicability}

Many Upper Mississippi River System (UMRS) backwaters have been degraded
by excessive amounts of sediment emanating from the basin, tributaries,
and mainstem sources. This degradation is in the form of loss of depth,
poor sediment quality, poor water quality, and sediment resuspension
that blocks light required by aquatic plants. Streambank erosion
throughout the basin is another important source of sediment that fills
the backwaters. (UMRR Design Handbook, 2012)

\subsection{Notes}\label{notes-2}

Restoration features now have a Program\_UID field to distinguish
between NESP and UMRR features. Aliases without program identification
will be created for any field that currently has UMRR in the name, so
they can be used for NESP or UMRR data.

Attribute criteria for each field includes the following:

\begin{itemize}
\tightlist
\item
  Attribute Status -- defines whether the field requires a value or if
  it is optional based on available information.
\item
  Phases -- indicates the earliest time in the project lifecycle values
  for this attribute as available.
\item
  Digitizing requirements -- indicates when the field value should be
  entered in the attribute table.
\end{itemize}

Geometry Note: Polygon nodes should have z values; elevations should be
taken from as-built survey or design documents.

\subsection{DredgingEvent}\label{dredgingevent}

(polygon) -- The extent of the area subject to dredging action. (area of
material removal)

\section{\texorpdfstring{\protect\includegraphics[width=0.75\linewidth,height=\textheight,keepaspectratio]{images/dredging.jpg}}{}}\label{section}

\section{Feature Attribute Criteria}\label{feature-attribute-criteria}

\#

Attribute

Attribute Status

Phase

Digitizing Requirements

1

Program\_UID

Required

Feasibility

At feature creation

2

PROJECT\_ID

Required

Feasibility

At feature creation

3

Type

Required

Feasibility

At feature creation

4

Subtype

Required

Feasibility

At feature creation

5

RestorFeatStatus

Optional

Design

At feature creation

6

sdsFeatureName

Optional

Design

As required and available

7

sdsFeatureDescription

Optional

Design

As required and available

8

RestorFeatVdatum

Optional

Design

As required and available

9

DredgeMethod

Optional

Design

As required and available

10

VolumeCY

Optional

Construction

As required and available

11

CutWidthFT

Optional

Construction

As required and available

12

CutLengtFT

Optional

Construction

As required and available

13

CutDepthFT

Optional

Construction

As required and available

14

sdsMetaDataID

Optional

Construction

As required and available

15

RestorFeatIDPK

Optional

Construction

As required and available

16

sdsID

Optional

Construction

As required and available

17

mediaIDFK

Optional

Construction

As required and available

\section{Parameters}\label{parameters}

\subsection{Program\_UID}\label{program_uid}

Designate program NESP or UMRR.

\subsection{Project\_ID (alias:
PROJECT\_ID)}\label{project_id-alias-project_id}

Foreign key to the unique NESP/UMRR project ID.

\subsection{Type (alias)}\label{type-alias}

A broad classification used to group NESP/UMRR restoration features into
categories of similar function.

Value

Description

Dredging

Area subject to dredging action.

\subsection{SubType (alias: Subtype)}\label{subtype-alias-subtype}

A classification used to group NESP/UMRR restoration features into
specific design feature categories.

Value

Description

Channel Dredging

Removal of sediment to create channels to redirect flow or provide
equipment access to construction areas

Deep hole dredging

Removal of sediment, using a variety of methods, to create backwater
lakes generally to increase fish habitat

Sediment basin dredging

Removal of sediment from or to create sediment basins

Geo-textile fill dredging

Removal of sediment to fill geo-textile tubes

Borrow

Removal of sediment from an area to be used for a specific purpose in
another area.

\subsection{RestorFeatStatus}\label{restorfeatstatus}

A description of the status of the individual restoration. It also
includes the party responsible for the individual feature.

Value

Description

USACE Planned

Feature is in a USACE planning stage.

USACE Built

Feature has been built by USACE. If a feature has been planned and built
by USACE, once constructed, it will always be considered USACE built
even if it has been repaired under later project stages.

USACE Enhanced

The function of an existing feature that was not built by USACE has been
enhanced by a feature that was USACE built or repaired.

USACE Repaired

An existing feature that was not built by USACE has been directly
improved or repaired.

Other Planned

A feature that is planned by another organization or government agency.

Other Built

Feature has been built by another organization or government agency and
is an integral part of the UMRR project.

Other Enhanced

The function of an existing feature, that was built by another
organization or government agency, has been enhanced by a feature that
was built or repaired by a non-USACE organization or government agency.

Other Repaired

An existing feature that has been directly improved or repaired by
another organization or government agency.

\subsection{sdsFeatureName}\label{sdsfeaturename}

The common name of the feature.

\subsection{sdsFeatureDescription}\label{sdsfeaturedescription}

A narrative describing the feature.

\subsection{RestorFeatVdatum}\label{restorfeatvdatum}

Vertical datum of elevations recorded in attribute table or used as
feature z values.

Value

Description

MSL 1912

Mean Sea Level of 1912

NGVD 1929

National Geodetic Vertical Datum of 1929

NAVD 1988

North American Vertical Datum of 1988

Other

Other, note datum in sdsFeatureDescription

Unknown

Unknown vertical datum

\subsection{DredgeMethod}\label{dredgemethod}

Method used to dredge.

Value

Description

Mechanical

Dredge method using mechanical equipment

Hydraulic

Dredge method using hydraulic equipment

High Solids

Dredge method used specifically for filling geotextile tubes

Unknown

Unknown

\subsection{VolumeCY}\label{volumecy}

Volume of material removed in cubic yards.

\subsection{CutWidthFT}\label{cutwidthft}

Width of dredge cut.

\subsection{CutLengthFT}\label{cutlengthft}

Length of dredge cut.

\subsection{CutDepthFT}\label{cutdepthft}

Average depth of dredge cut.

\subsection{sdsMetaDataID}\label{sdsmetadataid}

References attached metadata.

\subsection{RestorFeatIDPK}\label{restorfeatidpk}

References attached media.

\subsection{sdsID}\label{sdsid}

References attached media.

\subsection{mediaIDFK}\label{mediaidfk}

References attached media.

\chapter{Embankment}\label{embankment}

\section{Spatial \& Geometric
Definition}\label{spatial-geometric-definition-1}

\begin{itemize}
\item
  Geometry Type: Polygon (Embankment\_Footprint).
\item
  Spatial Delineation: Terrestrial earthen footprint bounded by the
  outer intersection lines of outer fill slopes with native subgrade
  alluvium.
\item
  Functional Roles: Isolating moist soil units, containing perimeter
  dredged placement cells, protecting wetlands from catastrophic river
  flood surges, and providing heavy maintenance vehicle access.
\end{itemize}

\subsection{Primary Restoration \& Engineering
Objectives}\label{primary-restoration-engineering-objectives}

\begin{itemize}
\item
  Hydraulic Isolation: Retain or exclude river floodwaters to facilitate
  independent moisture and pool-depth management.
\item
  Controlled Overtopping Resistance: Engineered to withstand controlled,
  non-catastrophic submergence during 10- to 25-year flood recurrence
  events.
\item
  Containment Integrity: Form stable retention dikes capable of
  impounding fine dredged hydraulic slurries without slope shear
  failures.
\end{itemize}

\subsection{Typical Cross-Sections, Dimensions \& Sizing
Criteria}\label{typical-cross-sections-dimensions-sizing-criteria-1}

\begin{itemize}
\item
  Crest Width: 10.0 to 14.0 ft for maintenance roadways; 8.0 ft for
  non-vehicular perimeter berms.
\item
  Side Slopes: Interior slopes at 1V:3H to 1V:4H; riverward exterior
  slopes at 1V:4H to 1V:6H to mitigate wind-wave runup.
\item
  Structural Core: Central compacted low-permeability cohesive clay core
  (CL/CH) with a minimum 3.0-ft bottom cutoff trench keyed into native
  impermeable strata.
\item
  Freeboard: Minimum 2.0 to 3.0 ft above maximum managed impoundment
  pool stage.
\end{itemize}

\subsection{Construction \& Implementation
Methodology}\label{construction-implementation-methodology-1}

\begin{itemize}
\item
  Subgrade Preparation: Stripping topsoil and organic muck (6 to 18 in);
  proof-rolling subgrade with heavy vibratory compactors.
\item
  Lift Placement \& Compaction: Clay and silty-clay fill spread in loose
  horizontal lifts of 8.0 in maximum, compacted to ≥95\% of Standard
  Proctor Maximum Dry Density (ASTM D698) within ±2\% of optimum
  moisture.
\item
  Slope Armoring: Class I to Class II rock riprap placed over non-woven
  geotextile on exterior wave faces; native perennial root mats
  established on interior slopes.
\end{itemize}

\subsection{Quality Control, Inspection \& Long-Term
Maintenance}\label{quality-control-inspection-long-term-maintenance-1}

\begin{itemize}
\item
  Geotechnical Monitoring: Piezometer installations along toe lines to
  monitor phreatic water lines and internal seepage gradients.
\item
  Routine Inspection: Quarterly inspections for burrowing animals
  (muskrats/beavers), desiccation cracking, mechanical rutting, and
  woody tree encroachment along levee toes.
\end{itemize}

\section{Spatial Feature Attribute}\label{spatial-feature-attribute-1}

\subsection{Definition}\label{definition-1}

An embankment is a raised earthen structure constructed to control or
contain water, provide flood protection, or create topographic diversity
within floodplain ecosystems.

\subsection{Purpose}\label{purpose-3}

Embankment systems are used to provide adequate flood protection, manage
water levels, and shape floodplain topography. They can isolate
management units, guide or restrict flows, protect infrastructure or
habitats from inundation or erosion, and create elevated or sheltered
areas that support diverse plant and wildlife communities.

\subsection{Applicability}\label{applicability-1}

Embankments are applicable wherever control of water levels or patterns
of inundation is needed in Upper Mississippi River System (UMRS)
restoration projects. Typical uses include forming the boundaries of
moist soil management units, creating berms around backwater lakes or
wetlands, protecting restoration features from unwanted flooding or
erosion, and shaping topographic relief to improve habitat diversity.
(UMRR Design Handbook, 2012)

\subsection{Notes}\label{notes-3}

Restoration features now have a Program\_UID field to distinguish
between NESP and UMRR features. Aliases without program identification
will be created for any field that currently has UMRR in the name, so
they can be used for NESP or UMRR data.

Attribute criteria for each field includes the following:

\begin{itemize}
\tightlist
\item
  Attribute Status -- defines whether the field requires a value or if
  it is optional based on available information.
\item
  Phases -- indicates the earliest time in the project lifecycle values
  for this attribute are available.
\item
  Digitizing requirements -- indicates when the field value should be
  entered in the attribute table.
\end{itemize}

Geometry Note: Elevations should be taken from as-built survey or design
documents.

\subsection{Embankment}\label{embankment-1}

(polygon) -- Embankments are raised earthen structures for the purpose
of controlling water or creating topographic diversity in floodplain
ecosystems.

\section{\texorpdfstring{\protect\includegraphics[width=0.75\linewidth,height=\textheight,keepaspectratio]{images/engriverembankment.jpg}}{}}\label{section-1}

\section{Feature Attribute Criteria}\label{feature-attribute-criteria-1}

\#

Attributes

Attribute Status

Phase

Digitizing Requirement

1

Program\_UID

Required

Feasibility

At feature creation

2

PROJECT\_ID

Required

Feasibility

At feature creation

3

Type

Required

Feasibility

At feature creation

4

Subtype

Required

Feasibility

At feature creation

5

RestorFeatStatus

Required

Design

At feature creation

6

sdsFeatureName

Optional

Design

As required and available

7

sdsFeatureDescription

Optional

Design

As required and available

8

RestorFeatVdatum

Optional

Design

As required and available

9

CrestWidthFT

Optional

Construction

As required and available

10

EnbankmentMaterial

Optional

Construction

As required and available

11

EnbankmentLenghtFT

Optional

Construction

As required and available

12

EmbankmentSideSlope

Optional

Construction

As required and available

13

sdsMetaDataID

Optional

Construction

As required and available

14

RestorFeatIDPK

Optional

Construction

As required and available

15

sdsID

Optional

Construction

As required and available

16

mediaIDFK

Optional

Construction

As required and available

\begin{center}\rule{0.5\linewidth}{0.5pt}\end{center}

\section{Parameters}\label{parameters-1}

\subsection{Program\_UID}\label{program_uid-1}

Designate program NESP or UMRR.

\subsection{Project\_ID (alias:
PROJECT\_ID)}\label{project_id-alias-project_id-1}

Foreign key to the unique NESP/UMRR project ID.

\subsection{Type (alias)}\label{type-alias-1}

A broad classification used to group NESP/UMRR restoration features into
categories of similar function.

\begin{quote}
Note: Specific coded values for Type were not provided in the source
text for Embankment.\\
In many implementations, Embankments fall under Floodplain Restoration
or Water Level Management.
\end{quote}

\subsection{SubType (alias: Subtype)}\label{subtype-alias-subtype-1}

A classification used to group NESP/UMRR restoration features into
specific embankment design feature categories.

\begin{quote}
Note: SubType coded values were not provided in the source text.\\
Example values might include Exterior Berm, Interior Berm, Levee
Segment, or Training Embankment, depending on project conventions.
\end{quote}

\subsection{RestorFeatStatus}\label{restorfeatstatus-1}

A description of the status of the individual restoration. It also
includes the party responsible for the individual feature.

Value

Description

USACE Planned

Feature is in a USACE planning stage.

USACE Built

Feature has been built by USACE. If a feature has been planned and built
by USACE, once constructed, it will always be considered USACE built
even if it has been repaired under later project stages.

USACE Enhanced

The function of an existing feature that was not built by USACE has been
enhanced by a feature that was USACE built or repaired.

USACE Repaired

An existing feature that was not built by USACE has been directly
improved or repaired.

Other Planned

A feature that is planned by another organization or government agency.

Other Built

Feature has been built by another organization or government agency and
is an integral part of the UMRR project.

Other Enhanced

The function of an existing feature, that was built by another
organization or government agency, has been enhanced by a feature that
was built or repaired by a non-USACE organization or government agency.

Other Repaired

An existing feature that has been directly improved or repaired by
another organization or government agency.

\subsection{sdsFeatureName}\label{sdsfeaturename-1}

The common name of the feature.

\subsection{sdsFeatureDescription}\label{sdsfeaturedescription-1}

A narrative describing the feature.

\subsection{RestorFeatVdatum}\label{restorfeatvdatum-1}

Vertical datum of elevations recorded in the attribute table or used as
feature z values.

Value

Description

MSL 1912

Mean Sea Level of 1912.

NGVD 1929

National Geodetic Vertical Datum of 1929.

NAVD 1988

North American Vertical Datum of 1988.

Other

Other; note datum in sdsFeatureDescription.

Unknown

Unknown vertical datum.

\subsection{CrestWidthFT}\label{crestwidthft}

Width of the embankment top or crest in feet.

\subsection{EnbankmentMaterial}\label{enbankmentmaterial}

Primary material used to construct the embankment.

Value

Description

Sand

Sand material.

Clay

Clay material.

Sand with clay core

Embankment composed primarily of sand with a central clay core.

Sand with clay cap

Embankment composed primarily of sand with a clay cap.

Earthen

Mixture of sand and clay, used when material type is not specified in
design or as-built drawings.

\subsection{EnbankmentLenghtFT}\label{enbankmentlenghtft}

Length of the embankment in feet (note: field name includes a
typographical error ``Lenght'' retained for schema consistency).

\subsection{EmbankmentSideSlope}\label{embankmentsideslope}

Side slope of the embankment, typically expressed as a
horizontal-to-vertical ratio (e.g., 3H:1V). May be stored as text or
numeric representation per data model convention.

\subsection{sdsMetaDataID}\label{sdsmetadataid-1}

References attached metadata.

\subsection{RestorFeatIDPK}\label{restorfeatidpk-1}

Primary key used to uniquely identify the restoration feature record.

\subsection{sdsID}\label{sdsid-1}

Identifier used to link the feature to supporting SDS (spatial data
system) records or related datasets.

\subsection{mediaIDFK}\label{mediaidfk-1}

Foreign key used to reference associated media (e.g., design drawings,
cross sections, as-built surveys, photographs) documenting the
embankment.

\chapter{Embankment Centerline}\label{embankment-centerline}

\section{Spatial \& Geometric
Definition}\label{spatial-geometric-definition-2}

\begin{itemize}
\item
  Geometry Type: Polyline (Embankment\_Centerline\_Line).
\item
  Spatial Delineation: Longitudinal survey axis running along the
  geometric centroid of the top crest of the embankment prism.
\item
  Data Fields: Linear stationing (0+00), top-of-crest design profile
  elevation, horizontal curve radii, design freeboard height, and crest
  width domain codes.
\end{itemize}

\subsection{Primary Restoration \& Engineering
Objectives}\label{primary-restoration-engineering-objectives-1}

\begin{itemize}
\item
  Engineering Geometry Control: Serves as the master horizontal and
  vertical control baseline for all construction layout, cross-sectional
  cut/fill grading, and post-settlement monitoring.
\item
  Profile Consistency: Establishes the uniform, non-dip crest grade to
  ensure overtopping occurs exclusively at designated, armored spillway
  notches rather than across unprotected earthen crest segments.
\end{itemize}

\subsection{Typical Cross-Sections, Dimensions \& Sizing
Criteria}\label{typical-cross-sections-dimensions-sizing-criteria-2}

\begin{itemize}
\item
  Survey Baseline: Stationing established at 100-ft linear increments
  with permanent control monuments positioned at Points of Curvature
  (PC) and Points of Tangency (PT).
\item
  Grade Tolerances: Construction grade tolerance along centerline is
  strict: −0.00 ft/+0.20 ft from contract design profile.
\end{itemize}

\subsection{Construction \& Implementation
Methodology}\label{construction-implementation-methodology-2}

\begin{itemize}
\item
  RTK-GPS \& Total Station Layout: Field staked with high-precision
  RTK-GPS tied directly into national State Plane Coordinates (NAD83)
  and NAVD88 elevation benchmarks.
\item
  Machine Control Integration: Centerline coordinate and profile strings
  loaded into 3D GPS machine-guidance earthmoving excavators and motor
  graders.
\end{itemize}

\subsection{Quality Control, Inspection \& Long-Term
Maintenance}\label{quality-control-inspection-long-term-maintenance-2}

\begin{itemize}
\tightlist
\item
  Settlement Monitoring: Permanent brass survey disks embedded in
  concrete surface monuments every 500 ft along the centerline, surveyed
  annually to track long-term primary consolidation of underlying river
  silts.
\end{itemize}

\section{Spatial Feature Attribute}\label{spatial-feature-attribute-2}

\subsection{Definition}\label{definition-2}

A linear embankment feature represents the centerline of a raised
earthen structure constructed to control or contain water, provide flood
protection, or create topographic diversity within floodplain
ecosystems.

\subsection{Purpose}\label{purpose-4}

Embankment systems are used to provide adequate flood protection, manage
water levels, and shape floodplain topography. Linear representation of
the embankment centerline supports design, hydraulic analysis, and asset
management by defining the primary alignment of the structure.
Embankments can isolate management units, guide or restrict flows,
protect infrastructure or habitats from inundation or erosion, and
create elevated or sheltered areas that support diverse plant and
wildlife communities.

\subsection{Applicability}\label{applicability-2}

Linear embankment features are applicable wherever centerline
representation of raised earthen structures is needed for planning,
design, modeling, or operations within Upper Mississippi River System
(UMRS) restoration projects. This includes levees, berms around MSMUs
and backwater lakes, training embankments, and other floodplain
earthworks that control water or create topographic variation. (UMRR
Design Handbook, 2012)

\subsection{Notes}\label{notes-4}

Restoration features now have a Program\_UID field to distinguish
between NESP and UMRR features. Aliases without program identification
will be created for any field that currently has UMRR in the name, so
they can be used for NESP or UMRR data.

Attribute criteria for each field includes the following:

\begin{itemize}
\tightlist
\item
  Attribute Status -- defines whether the field requires a value or if
  it is optional based on available information.
\item
  Phases -- indicates the earliest time in the project lifecycle values
  for this attribute are available.
\item
  Digitizing requirements -- indicates when the field value should be
  entered in the attribute table.
\end{itemize}

Geometry Note: Line should be 3D; elevations should be taken from
as-built survey or design documents.

\subsection{Embankment}\label{embankment-2}

(line) -- The centerline of a linear, raised earthen structure, for the
purpose of controlling water or creating topographic diversity in
floodplain ecosystems.

\begin{center}\rule{0.5\linewidth}{0.5pt}\end{center}

\section{Feature Attribute Criteria}\label{feature-attribute-criteria-2}

\#

Attributes

Attribute Status

Phase

Digitizing Requirement

1

Program\_UID

Required

Feasibility

At feature creation

2

PROJECT\_ID

Required

Feasibility

At feature creation

3

Type

Required

Feasibility

At feature creation

4

Subtype

Required

Feasibility

At feature creation

5

RestorFeatStatus

Required

Design

At feature creation

6

sdsFeatureName

Optional

Design

As required and available

7

sdsFeatureDescription

Optional

Design

As required and available

8

RestorFeatVdatum

Optional

Design

As required and available

9

CrestWidthFT

Optional

Construction

As required and available

10

EnbankmentMaterial

Optional

Construction

As required and available

11

EnbankmentLenghtFT

Optional

Construction

As required and available

12

EmbankmentSideSlope

Optional

Construction

As required and available

13

sdsMetaDataID

Optional

Construction

As required and available

14

RestorFeatIDPK

Optional

Construction

As required and available

15

sdsID

Optional

Construction

As required and available

16

mediaIDFK

Optional

Construction

As required and available

\begin{center}\rule{0.5\linewidth}{0.5pt}\end{center}

\section{Parameters}\label{parameters-2}

\subsection{Program\_UID}\label{program_uid-2}

Designate program NESP or UMRR.

\subsection{Project\_ID (alias:
PROJECT\_ID)}\label{project_id-alias-project_id-2}

Foreign key to the unique NESP/UMRR project ID.

\subsection{Type (alias)}\label{type-alias-2}

A broad classification used to group NESP/UMRR restoration features into
categories of similar function.

\begin{quote}
Note: Specific coded values for Type were not provided for the line
Embankment.\\
Embankments are commonly associated with Floodplain Restoration or Water
Level Management in project schemas.
\end{quote}

\subsection{SubType (alias: Subtype)}\label{subtype-alias-subtype-2}

A classification used to group NESP/UMRR restoration features into
specific embankment design feature categories.

\begin{quote}
Note: SubType coded values were not provided for the line Embankment.\\
Example values might include Exterior Berm, Interior Berm, Levee
Centerline, or Training Embankment, aligned with the polygon Embankment
subtypes.
\end{quote}

\subsection{RestorFeatStatus}\label{restorfeatstatus-2}

A description of the status of the individual restoration. It also
includes the party responsible for the individual feature.

Value

Description

USACE Planned

Feature is in a USACE planning stage.

USACE Built

Feature has been built by USACE. If a feature has been planned and built
by USACE, once constructed, it will always be considered USACE built
even if it has been repaired under later project stages.

USACE Enhanced

The function of an existing feature that was not built by USACE has been
enhanced by a feature that was USACE built or repaired.

USACE Repaired

An existing feature that was not built by USACE has been directly
improved or repaired.

Other Planned

A feature that is planned by another organization or government agency.

Other Built

Feature has been built by another organization or government agency and
is a integral part of the UMRR project.

Other Enhanced

The function of an existing feature, that was built by another
organization or government agency, has been enhanced by a feature that
was built or repaired by a non-USACE organization or government agency.

Other Repaired

An existing feature that has been directly improved or repaired by
another organization or government agency.

\subsection{sdsFeatureName}\label{sdsfeaturename-2}

The common name of the feature.

\subsection{sdsFeatureDescription}\label{sdsfeaturedescription-2}

A narrative describing the feature, including alignment, purpose, and
notable design characteristics.

\subsection{RestorFeatVdatum}\label{restorfeatvdatum-2}

Vertical datum of elevations recorded in the attribute table or used as
feature z values for the 3D line.

Value

Description

MSL 1912

Mean Sea Level of 1912.

NGVD 1929

National Geodetic Vertical Datum of 1929.

NAVD 1988

North American Vertical Datum of 1988.

Other

Other; note datum in sdsFeatureDescription.

Unknown

Unknown vertical datum.

\subsection{EnbankmentMaterial}\label{enbankmentmaterial-1}

Primary material used to construct the embankment, identified for the
centerline segment.

Value

Description

Sand

Sand material.

Clay

Clay material.

Sand with clay core

Embankment composed primarily of sand with a central clay core.

Sand with clay cap

Embankment composed primarily of sand with a clay cap.

Earthen

Some mixture of sand and clay, used when material type is not specified
in design or as-built drawings.

\subsection{CrestWidthFT}\label{crestwidthft-1}

Width of the embankment top or crest in feet (associated with the line
segment).

\subsection{EnbankmentLenghtFT}\label{enbankmentlenghtft-1}

Length of the embankment in feet for the represented line feature (note:
field name includes typographical error ``Lenght'' retained for schema
consistency).

\subsection{EmbankmentSideSlope}\label{embankmentsideslope-1}

Side slope of the embankment, expressed as a horizontal-to-vertical
ratio (e.g., 3H:1V), applicable to the represented centerline segment.

\subsection{sdsMetaDataID}\label{sdsmetadataid-2}

References attached metadata.

\subsection{RestorFeatIDPK}\label{restorfeatidpk-2}

Primary key used to uniquely identify the restoration feature record.

\subsection{sdsID}\label{sdsid-2}

Identifier used to link the feature to supporting SDS (spatial data
system) records or related datasets.

\subsection{mediaIDFK}\label{mediaidfk-2}

Foreign key used to reference associated media (e.g., design profiles,
cross-sections, levee certification documents, as-built surveys,
photographs) documenting the embankment.

\chapter{Fish Structure}\label{fish-structure}

\section{Spatial \& Geometric
Definition}\label{spatial-geometric-definition-3}

\begin{itemize}
\item
  Geometry Type: Point or Polygon (Fish\_Structure\_Location).
\item
  Spatial Delineation: Precise spatial coordinates or bounding
  footprints of submerged structural habitat units embedded in backwater
  cuts, secondary channels, and plunge basins.
\item
  Component Typology: Anchored large woody debris (LWD) complexes,
  rootwad revetments, brush bundles, boulder clusters, and wooden crib
  reefs.
\end{itemize}

\subsection{Primary Restoration \& Engineering
Objectives}\label{primary-restoration-engineering-objectives-2}

\begin{itemize}
\item
  Cover \& Velocity Shelters: Create localized zones of zero velocity
  (0.0 ft/s) alongside predatory feeding seams for gamefish species
  (largemouth bass, bluegill, black crappie, walleye).
\item
  Macroinvertebrate Substrate: Provide durable, non-silt-covered hard
  biological surfaces for net-spinning caddisflies, mayflies, and
  epilithic periphyton.
\item
  Thermal Microhabitats: Maintain localized deep cover within
  overwintering basins, preventing schooling fish from expending
  critical winter energy reserves.
\end{itemize}

\subsection{Typical Cross-Sections, Dimensions \& Sizing
Criteria}\label{typical-cross-sections-dimensions-sizing-criteria-3}

\begin{itemize}
\item
  Boulder Clusters: Groups of 3 to 7 interlocking Class IV--V boulders
  (d 50 =24 to 36 in), placed in staggered diamond formations spanning
  15×25 ft footprints.
\item
  Rootwad / Tree Deflectors: Hardwood tree trunks (≥18 in diameter) with
  intact root fans (6 to 10 ft diameter), projecting 10 to 20 ft from
  channel margins.
\item
  Ballast Requirements: Stone ballast sizing calculated to offset
  submerged wood buoyancy by a factor of safety FS≥1.75 (1.5 to 2.0 tons
  of rock per full-size tree bole).
\end{itemize}

\subsection{Construction \& Implementation
Methodology}\label{construction-implementation-methodology-3}

\begin{itemize}
\item
  Barge-Mounted Crane Placement: Lowered into aquatic cuts using cable
  chokers or mechanical clamshells.
\item
  Anchoring Systems: Hardwoods secured to riverbeds using duckbill earth
  anchors, helical steel ground screws, or heavy galvanized aircraft
  cables (3/8 in diameter) clamped through drilled stone anchors.
\end{itemize}

\subsection{Quality Control, Inspection \& Long-Term
Maintenance}\label{quality-control-inspection-long-term-maintenance-3}

\begin{itemize}
\item
  Underwater Acoustic Surveys: High-resolution multi-beam side-scan
  sonar runs conducted at 2- and 5-year intervals to verify structural
  presence, orientation, and absence of silt smothering.
\item
  Biological Electrofishing: Annual fisheries index surveys comparing
  fish density at structure sites against un-structured control reaches.
\end{itemize}

\section{Spatial Feature Attributes}\label{spatial-feature-attributes}

\subsection{Definition}\label{definition-3}

FishStructure features are engineered structures designed to improve
fish habitat or allow improved fish movement through the river system by
modifying hydraulics, providing cover, or circumventing in-channel
barriers.

\subsection{Purpose}\label{purpose-5}

Improving fish habitat and facilitating fish movement through the river
system are critical steps in maintaining the health, biodiversity, and
sustainability of the UMR ecosystem. These efforts support a wide range
of species, restore ecological functions, enhance fish populations, and
contribute to the local economy through sustainable fisheries and
recreation.

\subsection{Applicability}\label{applicability-3}

Fish structures are applicable wherever in-channel barriers, degraded
habitat, or limited connectivity constrain fish movement or habitat
quality within Upper Mississippi River System (UMRS) restoration
projects. They are typically used at dams, diversion structures,
constricted channels, and along eroding banks or simplified shorelines
to restore passage, create refuge, and diversify hydraulic and substrate
conditions. (UMRR Design Handbook, 2012)

\subsection{Notes}\label{notes-5}

Restoration features now have a Program\_UID field to distinguish
between NESP and UMRR features. Aliases without program identification
will be created for any field that currently has UMRR in the name, so
they can be used for NESP or UMRR data.

Attribute criteria for each field includes the following:

\begin{itemize}
\tightlist
\item
  Attribute Status -- defines whether the field requires a value or if
  it is optional based on available information.
\item
  Phases -- indicates the earliest time in the project lifecycle values
  for this attribute are available.
\item
  Digitizing requirements -- indicates when the field value should be
  entered in the attribute table.
\end{itemize}

Geometry Note: Points should have well top elevation as Z value if
available.

\subsection{FishStructure}\label{fishstructure}

(polygon) -- An engineered structure designed to improve fish habitat or
allow improved fish movement through the river system.

\begin{center}\rule{0.5\linewidth}{0.5pt}\end{center}

\section{Feature Attribute Criteria}\label{feature-attribute-criteria-3}

\#

Attributes

Attribute Status

Phase

Digitizing Requirement

1

Program\_UID

Required

Feasibility

At feature creation

2

PROJECT\_ID

Required

Feasibility

At feature creation

3

Type

Required

Feasibility

At feature creation

4

Subtype

Required

Feasibility

At feature creation

5

RestorFeatStatus

Required

Design

At feature creation

6

sdsFeatureName

Optional

Design

As required and available

7

sdsFeatureDescription

Optional

Design

As required and available

8

RestorFeatVdatum

Optional

Design

As required and available

9

TargetFishSpecies

Optional

Design

As required and available

10

StartStation

Optional

Construction

As required and available

11

EndStation

Optional

Construction

As required and available

12

sdsMetaDataID

Optional

Construction

As required and available

13

RestorFeatIDPK

Optional

Construction

As required and available

14

sdsID

Optional

Construction

As required and available

15

mediaIDFK

Optional

Construction

As required and available

\begin{center}\rule{0.5\linewidth}{0.5pt}\end{center}

\section{Parameters}\label{parameters-3}

\subsection{Program\_UID}\label{program_uid-3}

Designate program NESP or UMRR.

\subsection{Project\_ID (alias:
PROJECT\_ID)}\label{project_id-alias-project_id-3}

Foreign key to the unique NESP/UMRR project ID.

\subsection{Type (alias)}\label{type-alias-3}

A broad classification used to group NESP/UMRR restoration features into
categories of similar function.

Value

Description

Channel Modifications

Measures taken to train the main channel of a river and to modify
secondary channels.

\subsection{SubType (alias: Subtype)}\label{subtype-alias-subtype-3}

A classification used to group NESP/UMRR restoration features into
specific fish structure design feature categories.

Value

Description

Fish LUNKERS

Engineered overhanging-bank structure that provides fish habitat and
bank stabilization, generally about 8 ft long, 1--2 ft tall and 3 ft
deep, constructed of hardwood, concrete, or plastic lumber.

Fish bypass channel

Channel created to circumvent in-channel barriers to natural fish
movement.

Fish ladder

Engineered step or pool-and-weir structure created to circumvent
in-channel barriers to natural fish movement.

Natural boulder ramp

Short, steep ramp comprised of boulders created to circumvent in-channel
barriers to natural fish movement.

\subsection{RestorFeatStatus}\label{restorfeatstatus-3}

A description of the status of the individual restoration. It also
includes the party responsible for the individual feature.

Value

Description

USACE Planned

Feature is in a USACE planning stage.

USACE Built

Feature has been built by USACE. If a feature has been planned and built
by USACE, once constructed, it will always be considered USACE built
even if it has been repaired under later project stages.

USACE Enhanced

The function of an existing feature that was not built by USACE has been
enhanced by a feature that was USACE built or repaired.

USACE Repaired

An existing feature that was not built by USACE has been directly
improved or repaired.

Other Planned

A feature that is planned by another organization or government agency.

Other Built

Feature has been built by another organization or government agency and
is an integral part of the UMRR project.

Other Enhanced

The function of an existing feature, that was built by another
organization or government agency, has been enhanced by a feature that
was built or repaired by a non-USACE organization or government agency.

Other Repaired

An existing feature that has been directly improved or repaired by
another organization or government agency.

\subsection{sdsFeatureName}\label{sdsfeaturename-3}

The common name of the feature.

\subsection{sdsFeatureDescription}\label{sdsfeaturedescription-3}

A narrative describing the feature, including design intent, structural
details, and operational assumptions.

\subsection{RestorFeatVdatum}\label{restorfeatvdatum-3}

Vertical datum of elevations recorded in the attribute table or used as
feature z values.

Value

Description

MSL 1912

Mean Sea Level of 1912.

NGVD 1929

National Geodetic Vertical Datum of 1929.

NAVD 1988

North American Vertical Datum of 1988.

Other

Other; note datum in sdsFeatureDescription.

Unknown

Unknown vertical datum.

\subsection{TargetFishSpecies}\label{targetfishspecies}

Specific fish species or species groups the structure is intended to
benefit (e.g., pallid sturgeon, paddlefish, centrarchids). May include
one or more species names or functional groups.

\subsection{StartStation}\label{startstation}

Beginning station of the feature, typically referenced to a project or
river engineering baseline (e.g., 10+00).

\subsection{EndStation}\label{endstation}

Ending station of the feature along the same baseline.

\subsection{sdsMetaDataID}\label{sdsmetadataid-3}

References attached metadata (e.g., design reports, hydraulic modeling
files, biological objectives).

\subsection{RestorFeatIDPK}\label{restorfeatidpk-3}

Primary key used to uniquely identify the restoration feature record.

\subsection{sdsID}\label{sdsid-3}

Identifier used to link the feature to supporting SDS (spatial data
system) records or related datasets.

\subsection{mediaIDFK}\label{mediaidfk-3}

Foreign key used to reference associated media (e.g., design drawings,
plans and profiles, construction photographs, as-built surveys)
documenting the fish structure.

\chapter{Flora Planting Seeding}\label{flora-planting-seeding}

\section{Spatial \& Geometric
Definition}\label{spatial-geometric-definition-4}

\begin{itemize}
\item
  Geometry Type: Polygon (Flora\_Planting\_Area).
\item
  Spatial Delineation: Managed zones receiving revegetation treatments
  across restored island crests, moist soil units, dredge disposal
  sites, and riparian forest floodplains.
\item
  Vegetation Types: Hard-mast bottomland hardwoods, native emergent
  wetland tubers/plugs, pollinator forb mixes, and deep-rooted native
  warm-season prairie grasses.
\end{itemize}

\subsection{Primary Restoration \& Engineering
Objectives}\label{primary-restoration-engineering-objectives-3}

\begin{itemize}
\item
  Soil Bio-Stabilization: Rapid root-mesh development to anchor sandy
  and fine-grained dredged placement materials against wave-wash and
  flood shear stress.
\item
  Wildlife Mast Production: Re-establish critical hard-mast forest
  species (swamp white oak, pin oak, shellbark hickory, pecan) displaced
  by altered post-dam hydrologic regimes.
\item
  Invasive Species Suppression: Establish continuous native vegetative
  cover to prevent colonization by monotypic reed canary grass (Phalaris
  arundinacea).
\end{itemize}

\subsection{Typical Specifications, Densities \& Seed
Mixes}\label{typical-specifications-densities-seed-mixes}

\begin{itemize}
\item
  Hard-Mast Tree Stock: Root Production Method (RPM) or bare-root
  containerized saplings (1.0-gal to 3.0-gal pots, ≥0.5 in caliper).
  Planted at densities of 200 to 300 stems/acre on 12×12 ft to 15×15 ft
  grids.
\item
  Riparian Seed Mixes: Native cool- and warm-season grasses (Virginia
  wildrye, river bulrush, cordgrass, bluejoint grass) broadcast at 10 to
  20 lbs Pure Live Seed (PLS)/acre.
\item
  Emergent Wetlands: Bare-root tubers/plugs of Sagittaria latifolia and
  Scirpus fluviatilis installed along littoral margins (0.5 to 2.0 ft
  water depth) on 3-ft centers.
\end{itemize}

\subsection{Construction \& Implementation
Methodology}\label{construction-implementation-methodology-4}

\begin{itemize}
\item
  Site Preparation: Multi-pass forestry mowing followed by
  aquatic-approved glyphosate herbicide treatments to eradicate invasive
  root mats.
\item
  Shelter Tube Installation: Hardwood saplings protected with 48- to
  60-in rigid, ventilated photodegradable polyethylene tree shelters
  staked with rot-resistant white oak stakes to eliminate deer browse
  and rodent girdling.
\end{itemize}

\subsection{Quality Control, Inspection \& Long-Term
Maintenance}\label{quality-control-inspection-long-term-maintenance-4}

\begin{itemize}
\item
  Survival Censuses: Fixed-radius inventory plots (0.1 acre) evaluated
  at months 12, 24, and 36 post-planting. Target standard is ≥80\%
  survival density.
\item
  Herbicide Touch-Ups: Backpack spot-spraying or mechanical
  brush-cutting around tree collars for 3 consecutive growing seasons
  until sapling crowns overtop competing weeds.
\end{itemize}

\section{Spatial Feature Attribute}\label{spatial-feature-attribute-3}

\subsection{Definition}\label{definition-4}

Flora Planting or Seeding is a floodplain restoration measure in which
vegetation is intentionally planted or seeded to improve natural
floodplain function and habitat.

\subsection{Purpose}\label{purpose-6}

Flora planting or seeding is used to establish or enhance native plant
communities, stabilize soils, increase structural and species diversity,
and provide food and cover for wildlife. These actions support long-term
habitat goals, improve resilience to flooding and disturbance, and help
restore natural floodplain processes such as nutrient cycling and
sediment retention.

\subsection{Applicability}\label{applicability-4}

Flora planting or seeding is applicable wherever vegetation needs to be
established, enhanced, or converted within Upper Mississippi River
System (UMRS) restoration projects. Typical applications include
revegetating disturbed areas, restoring floodplain forests, wetlands,
and prairies, improving shoreline and island stability, and enhancing
terrestrial and aquatic edge habitats. (UMRR Design Handbook, 2012)

\subsection{Notes}\label{notes-6}

Restoration features now have a Program\_UID field to distinguish
between NESP and UMRR features. Aliases without program identification
will be created for any field that currently has UMRR in the name, so
they can be used for NESP or UMRR data.

Attribute criteria for each field includes the following:

\begin{itemize}
\tightlist
\item
  Attribute Status -- defines whether the field requires a value or if
  it is optional based on available information.
\item
  Phases -- indicates the earliest time in the project lifecycle values
  for this attribute are available.
\item
  Digitizing requirements -- indicates when the field value should be
  entered in the attribute table.
\end{itemize}

Geometry Note: No Z values are needed for this feature.

\subsection{FloraPlantingOrSeeding}\label{floraplantingorseeding}

(polygon) -- A location where vegetation is planted or seeded.

\begin{center}\rule{0.5\linewidth}{0.5pt}\end{center}

\section{Feature Attribute Criteria}\label{feature-attribute-criteria-4}

\#

Attributes

Attribute Status

Phase

Digitizing Requirement

1

Program\_UID

Required

Feasibility

At feature creation

2

PROJECT\_ID

Required

Feasibility

At feature creation

3

Type

Required

Feasibility

At feature creation

4

Subtype

Required

Feasibility

At feature creation

5

RestorFeatStatus

Required

Design

At feature creation

6

sdsFeatureName

Optional

Design

As required and available

7

sdsFeatureDescription

Optional

Design

As required and available

8

RestorFeatVdatum

Optional

Design

As required and available

9

PlantingDate

Optional

Design

As required and available

10

DominantSpecies1

Optional

Construction

As required and available

11

DominantSpecies2

Optional

Construction

As required and available

12

DominantSpecies3

Optional

Construction

As required and available

13

MastPlanting

Optional

Construction

As required and available

14

sdsMetaDataID

Optional

Construction

As required and available

15

RestorFeatIDPK

Optional

Construction

As required and available

16

sdsID

Optional

Construction

As required and available

17

mediaIDFK

Optional

Construction

As required and available

\begin{center}\rule{0.5\linewidth}{0.5pt}\end{center}

\section{Parameters}\label{parameters-4}

\subsection{Program\_UID}\label{program_uid-4}

Designate program NESP or UMRR.

\subsection{Project\_ID (alias:
PROJECT\_ID)}\label{project_id-alias-project_id-4}

Foreign key to the unique NESP/UMRR project ID.

\subsection{Type (alias)}\label{type-alias-4}

A broad classification used to group NESP/UMRR restoration features into
categories of similar function.

Value

Description

Shoreline Protection

Measures taken to reduce erosion from banks and islands due to wind,
wave, and ice action.

Floodplain Restoration

Measures taken to improve natural floodplain function and habitat.

\subsection{SubType (alias: Subtype)}\label{subtype-alias-subtype-4}

A classification used to group NESP/UMRR restoration features into
specific planting or seeding design feature categories.

Value

Description

Tree planting

Planting of tree species in floodplain areas that have appropriate soil
types and water regimes.

Wetland planting

Planting of wetland species in areas with suitable hydrology and soils
to establish or enhance wetland habitat.

Prairie Planting

Planting of prairie species in floodplain areas with appropriate soil
types and water regimes.

Floodplain Forest Planting

Planting or augmentation of native floodplain tree and shrub species to
establish or restore forested floodplain habitat.

Timber Stand Improvement

Silvicultural treatments within existing forest stands (e.g., thinning,
release, underplanting) to improve species composition, structure, and
habitat quality.

Aquatic Vegetation

Planting or seeding of emergent, submergent, or floating-leaved aquatic
vegetation to improve aquatic habitat and stabilize substrates.

\subsection{RestorFeatStatus}\label{restorfeatstatus-4}

A description of the status of the individual restoration. It also
includes the party responsible for the individual feature.

Value

Description

USACE Planned

Feature is in a USACE planning stage.

USACE Built

Feature has been implemented (planted/seeded) by USACE. Once implemented
by USACE, it remains classified as USACE Built even if later maintained
or repaired.

USACE Enhanced

The function of an existing planting or habitat area not originally
established by USACE has been enhanced by a USACE action.

USACE Repaired

An existing planting area not established by USACE has been directly
improved or repaired by USACE (e.g., replanting after failure).

Other Planned

Planting/seeding feature planned by another organization or government
agency.

Other Built

Planting/seeding feature implemented by another organization or
government agency and integral to the UMRR project.

Other Enhanced

Function of an existing planting, built by another organization,
enhanced by a non-USACE organization or government agency.

Other Repaired

Existing planting directly improved or repaired by another organization
or government agency.

\subsection{sdsFeatureName}\label{sdsfeaturename-4}

The common or management name of the planting area.

\subsection{sdsFeatureDescription}\label{sdsfeaturedescription-4}

A narrative describing the planting or seeding feature, including
objectives, methods, and any notable site conditions.

\subsection{RestorFeatVdatum}\label{restorfeatvdatum-4}

Vertical datum of any elevations recorded in the attribute table (if
used for related features).

Value

Description

MSL 1912

Mean Sea Level of 1912.

NGVD 1929

National Geodetic Vertical Datum of 1929.

NAVD 1988

North American Vertical Datum of 1988.

Other

Other; note datum in sdsFeatureDescription.

Unknown

Unknown vertical datum.

\subsection{PlantingDate}\label{plantingdate}

The date the planting or seeding was carried out (or first completed).

\subsection{DominantSpecies1}\label{dominantspecies1}

Dominant or co-dominant species planted or seeded within the feature.
Percentage of the planting mix can be included as part of the value or
in sdsFeatureDescription.

\subsection{DominantSpecies2}\label{dominantspecies2}

Second dominant or co-dominant species planted or seeded. Percentage of
the planting mix can be included.

\subsection{DominantSpecies3}\label{dominantspecies3}

Third dominant or co-dominant species planted or seeded. Percentage of
the planting mix can be included.

\subsection{MastPlanting}\label{mastplanting}

Indicates whether the planting provides mast (nuts, acorns, berries,
etc.) important to wildlife species.

Value

Description

Yes

Planting intended to provide mast for wildlife.

No

Planting not specifically intended as a mast-producing planting.

Unknown

Mast function not specified.

\subsection{sdsMetaDataID}\label{sdsmetadataid-4}

References attached metadata (e.g., seed mix specifications, planting
plans, monitoring protocols).

\subsection{RestorFeatIDPK}\label{restorfeatidpk-4}

Primary key used to uniquely identify the restoration feature record.

\subsection{sdsID}\label{sdsid-4}

Identifier used to link the feature to supporting SDS (spatial data
system) records or related datasets.

\subsection{mediaIDFK}\label{mediaidfk-4}

Foreign key used to reference associated media (e.g., photographs,
planting plans, seed tags, monitoring reports) documenting the planting
or seeding.

\chapter{Flow Channel}\label{flow-channel}

\section{Spatial \& Geometric
Definition}\label{spatial-geometric-definition-5}

\begin{itemize}
\item
  Geometry Type: Polygon (Flow\_Channel\_Footprint).
\item
  Spatial Delineation: Channel boundaries of excavated aquatic chutes,
  secondary conveyance corridors, or backwater flow-supply corridors
  reconnecting the main river to isolated interior wetlands.
\item
  Hydraulic Classification: Continuous low-velocity flow channels,
  seasonal skimming channels, or flood pulse distributors.
\end{itemize}

\subsection{Primary Restoration \& Engineering
Objectives}\label{primary-restoration-engineering-objectives-4}

\begin{itemize}
\item
  Oxygenated Water Delivery: Introduce low volumes of oxygen-rich river
  water into stagnant interior backwaters to prevent winter and summer
  anoxic fish kills.
\item
  Biological Corridor Connectivity: Provide passage routes for migratory
  riverine fish traveling between river feeding reaches and sheltered
  nursery lakes.
\item
  Controlled Energy Conveyance: Route discharge without inducing bed
  scour, toe destabilization, or excessive sediment imports into fragile
  interior basins.
\end{itemize}

\subsection{Typical Cross-Sections, Dimensions \& Sizing
Criteria}\label{typical-cross-sections-dimensions-sizing-criteria-4}

\begin{itemize}
\item
  Bottom Widths: 20 to 60 ft; side slopes stabilized at 1V:3H or
  flatter.
\item
  Longitudinal Bed Slope: Very flat gradient (S≤0.0002 to 0.0005) to
  constrain cross-sectional average velocities beneath critical erosive
  limits (v≤1.0 to 1.5 ft/s).
\item
  Inlet Control Sills: Hard-armored rock sills placed at the channel
  inlet set 0.5 to 1.5 ft above normal pool to block bedload sediment
  entry during low-water regimes.
\end{itemize}

\subsection{Construction \& Implementation
Methodology}\label{construction-implementation-methodology-5}

\begin{itemize}
\item
  Excavation Fleet: Low-ground-pressure (LGP) amphibious excavators or
  spud-barge mechanical dredge rigs working from downstream upstream to
  maintain positive gravity dewatering.
\item
  Bank Armor Integration: Riprap placement along channel entrance mouths
  and sharp bendway curves subject to centrifugal shear forces.
\end{itemize}

\subsection{Quality Control, Inspection \& Long-Term
Maintenance}\label{quality-control-inspection-long-term-maintenance-5}

\begin{itemize}
\item
  Acoustic Flow Gauging: Seasonal ADCP velocity transects to confirm
  flow splits match calibrated hydrodynamic model predictions (typically
  1\% to 5\% of main river discharge).
\item
  Debris Clearance: Annual post-flood inspection to clear lodged drift
  logs and alluvial sand wedges from channel entrances.
\end{itemize}

\section{Spatial Feature Attribute}\label{spatial-feature-attribute-4}

\subsection{Definition}\label{definition-5}

A FlowChannel is an excavated channel, created by mechanical or
hydraulic means, intended to allow the flow of water across, into, or
out of the project area or any subcomponent of a project area.

\subsection{Purpose}\label{purpose-7}

Flow channels are constructed to direct, convey, or redistribute water
in support of restoration objectives. They can be used to reconnect
floodplains or backwaters, improve circulation, manage water levels,
enhance habitat connectivity, and reduce stagnant or low-quality water
conditions by providing designed flow paths.

\subsection{Applicability}\label{applicability-5}

Flow channels are applicable wherever controlled or enhanced conveyance
of water is needed within Upper Mississippi River System (UMRS)
restoration projects. They are used in both terrestrial and aquatic
settings to route water between project features (e.g., MSMUs, backwater
lakes, side channels), restore or mimic natural flow patterns, and
support habitat, sediment, and water quality objectives. (UMRR Design
Handbook, 2012)

\subsection{Notes}\label{notes-7}

Restoration features now have a Program\_UID field to distinguish
between NESP and UMRR features. Aliases without program identification
will be created for any field that currently has UMRR in the name, so
they can be used for NESP or UMRR data.

Attribute criteria for each field includes the following:

\begin{itemize}
\tightlist
\item
  Attribute Status -- defines whether the field requires a value or if
  it is optional based on available information.
\item
  Phases -- indicates the earliest time in the project lifecycle values
  for this attribute are available.
\item
  Digitizing requirements -- indicates when the field value should be
  entered in the attribute table.
\end{itemize}

Geometry Note: Polygon nodes should have z values; elevations should be
taken from as-built survey or design documents.

\subsection{FlowChannel}\label{flowchannel}

(polygon) -- An excavated channel (via mechanical or hydraulic means)
intended to allow for the flow of water across, into, or out of the
project area or any subcomponent of a project area.

\begin{center}\rule{0.5\linewidth}{0.5pt}\end{center}

\section{Feature Attribute Criteria}\label{feature-attribute-criteria-5}

\#

Attribute

Attribute Status

Phase

Digitizing Requirements

1

Program\_UID

Required

Feasibility

At feature creation

2

PROJECT\_ID (Alias)

Required

Feasibility

At feature creation

3

Type (Alias)

Required

Feasibility

At feature creation

4

Subtype (Alias)

Required

Feasibility

At feature creation

5

RestorFeatStatus

Required

Feasibility

At feature creation

6

sdsFeatureName

Optional

Design

As required and available

7

sdsFeatureDescription

Optional

Design

As required and available

8

RestorFeatVdatum

Optional

Design

As required and available

9

ConstructionMethod

Optional

Design

As required and available

10

ConstructionMethod

Optional

Design

As required and available

11

VolumeCY

Optional

Design

As required and available

12

CutWidthFT

Optional

Design

As required and available

13

CutLengthFT

Optional

Design

As required and available

14

CutDepthFT

Optional

Construction

As required and available

15

sdsMetaDataID

Optional

Construction

As required and available

16

RestorFeatIDPK

Optional

Construction

As required and available

17

sdsID

Optional

Construction

As required and available

18

mediaIDFK

Optional

Construction

As required and available

\begin{center}\rule{0.5\linewidth}{0.5pt}\end{center}

\section{Parameters}\label{parameters-5}

\subsection{Program\_UID}\label{program_uid-5}

Designate program NESP or UMRR.

\subsection{Project\_ID (alias:
PROJECT\_ID)}\label{project_id-alias-project_id-5}

Foreign key to the unique NESP/UMRR project ID.

\subsection{Type (alias: Type)}\label{type-alias-type}

A broad classification used to group NESP/UMRR restoration features into
categories of similar function.

Value

Description

Flow Channel

An excavated channel (via mechanical or hydraulic means) intended to
allow for the flow of water across, into, or out of the project area or
any subcomponent of a project area.

\subsection{SubType (alias: Subtype)}\label{subtype-alias-subtype-5}

A classification used to group NESP/UMRR restoration features into
specific design feature categories.

Value

Description

Terrestrial

Removal of sediment to create channels to allow for flow of water across
or through predominantly terrestrial areas.

Aquatic

Removal of sediment to create channels to allow for flow of water within
predominantly aquatic areas.

\subsection{RestorFeatStatus}\label{restorfeatstatus-5}

A description of the status of the individual restoration. It also
includes the party responsible for the individual feature.

Value

Description

USACE Planned

Feature is in a USACE planning stage.

USACE Built

Feature has been built by USACE. If a feature has been planned and built
by USACE, once constructed, it will always be considered USACE built
even if it has been repaired under later project stages.

USACE Enhanced

The function of an existing feature that was not built by USACE has been
enhanced by a feature that was USACE built or repaired.

USACE Repaired

An existing feature that was not built by USACE has been directly
improved or repaired.

Other Planned

A feature that is planned by another organization or government agency.

Other Built

Feature has been built by another organization or government agency and
is an integral part of the UMRR project.

Other Enhanced

The function of an existing feature, that was built by another
organization or government agency, has been enhanced by a feature that
was built or repaired by a non-USACE organization or government agency.

Other Repaired

An existing feature that has been directly improved or repaired by
another organization or government agency.

\subsection{sdsFeatureName}\label{sdsfeaturename-5}

The common name of the feature.

\subsection{sdsFeatureDescription}\label{sdsfeaturedescription-5}

A narrative describing the feature.

\subsection{RestorFeatVdatum}\label{restorfeatvdatum-5}

Vertical datum of elevations recorded in the attribute table or used as
feature z values.

Value

Description

MSL 1912

Mean Sea Level of 1912.

NGVD 1929

National Geodetic Vertical Datum of 1929.

NAVD 1988

North American Vertical Datum of 1988.

Other

Other; note datum in sdsFeatureDescription.

Unknown

Unknown vertical datum.

\subsection{ConstructionMethod}\label{constructionmethod}

Method used to remove sediment and construct the flow channel.

Value

Description

Mechanical

Removal method using mechanical equipment (e.g., excavators, draglines).

Hydraulic

Removal method using hydraulic equipment (e.g., cutterhead or plain
suction dredges).

\subsection{VolumeCY}\label{volumecy-1}

Volume of material removed in cubic yards.

\subsection{CutWidthFT}\label{cutwidthft-1}

Width of the excavated channel in feet.

\subsection{CutLengthFT}\label{cutlengthft-1}

Length of the excavated channel in feet.

\subsection{CutDepthFT}\label{cutdepthft-1}

Average depth of the excavated channel in feet.

\subsection{sdsMetaDataID}\label{sdsmetadataid-5}

References attached metadata.

\subsection{RestorFeatIDPK}\label{restorfeatidpk-5}

Primary key used to uniquely identify the restoration feature record.

\subsection{sdsID}\label{sdsid-5}

Identifier used to link the feature to supporting SDS (spatial data
system) records or related datasets.

\subsection{mediaIDFK}\label{mediaidfk-5}

Foreign key used to reference associated media (e.g., design drawings,
dredging records, as-built plans, photographs) documenting the flow
channel.

\chapter{Island Centerline}\label{island-centerline}

\section{Spatial \& Geometric
Definition}\label{spatial-geometric-definition-6}

\begin{itemize}
\item
  Geometry Type: Polyline (Island\_Crest\_Line or Island\_Baseline).
\item
  Spatial Delineation: Centerline axis, design baseline, or high-crest
  ridge running continuously along the longitudinal alignment of a
  constructed island.
\item
  Data Fields: Island segment ID, design crest profile grade, baseline
  stationing, revetment armoring class, and design overtopping
  threshold.
\end{itemize}

\subsection{Primary Restoration \& Engineering
Objectives}\label{primary-restoration-engineering-objectives-5}

\begin{itemize}
\item
  Survey Baseline Control: Serves as the primary reference vector from
  which cross-sectional design offset templates (perpendicular
  cross-sections) are generated and staked.
\item
  Hydraulic Ridge Integrity: Tracks the continuous elevation of the
  island's primary windbreak and flood-deflection backbone.
\end{itemize}

\subsection{Typical Specifications \& Geometric
Properties}\label{typical-specifications-geometric-properties}

\begin{itemize}
\item
  Stationing System: Formatted identically to civil alignments (0+00 to
  L), with stationing typically progressing from the downstream trailing
  tip to the upstream prow.
\item
  Elevation Control: Tied strictly to the project vertical control
  network (NAVD88) with construction staking tolerances of ±0.1 ft.
\end{itemize}

\subsection{Construction \& Implementation
Methodology}\label{construction-implementation-methodology-6}

\begin{itemize}
\tightlist
\item
  GPS Navigation Strings: Exported as civil design alignment files into
  real-time barge GPS tracking suites to guide placement derrick
  operators and bulldozer operators.
\end{itemize}

\subsection{Quality Control, Inspection \& Long-Term
Maintenance}\label{quality-control-inspection-long-term-maintenance-6}

\begin{itemize}
\tightlist
\item
  Crest Profile Resurveys: Periodic RTK-GPS profile walks along the
  island line following major flood inundation cycles to calculate
  post-consolidation subgrade settlement.
\end{itemize}

\section{Spatial Feature Attribute}\label{spatial-feature-attribute-5}

\subsection{Definition}\label{definition-6}

IslandLine represents the linear footprint or alignment of an island or
island segment within a river or backwater area when the feature is too
narrow or detailed to be mapped effectively as a polygon at project
scale. It captures the engineered or restored island crest or centerline
used for analysis, design reference, or cartographic depiction.

\subsection{Purpose}\label{purpose-8}

IslandLine features are used to represent long, narrow, or small island
structures where polygon mapping is impractical, while still documenting
their location, extent, and functional role in influencing flow, wave
action, and habitat conditions. They provide a simplified linear
representation that can be used in hydraulic analysis, design layouts,
and monitoring of island performance over time.

\subsection{Applicability}\label{applicability-6}

IslandLine is applicable where island features are present but are
either too small to digitize as polygons at the project mapping scale or
where a linear representation (e.g., crest line, centerline, or control
alignment) is needed for modeling, design, or management. Typical uses
include narrow barrier islands, training-island alignments, or thin
island ridges within Upper Mississippi River System (UMRS) restoration
projects. (UMRR Design Handbook, 2012)

\subsection{Notes}\label{notes-8}

Restoration features now have a Program\_UID field to distinguish
between NESP and UMRR features. Aliases without program identification
will be created for any field that currently has UMRR in the name, so
they can be used for NESP or UMRR data.

Attribute criteria for each field includes the following:

\begin{itemize}
\tightlist
\item
  Attribute Status -- defines whether the field requires a value or if
  it is optional based on available information.
\item
  Phases -- indicates the earliest time in the project lifecycle values
  for this attribute are available.
\item
  Digitizing requirements -- indicates when the field value should be
  entered in the attribute table.
\end{itemize}

Geometry Note: No Z values are needed for this feature.

\subsection{IslandLine}\label{islandline}

(line) -- Linear representation of island structures engineered with a
variety of techniques and materials that serve as barriers to flow,
wind, and waves.

\begin{center}\rule{0.5\linewidth}{0.5pt}\end{center}

\section{Feature Attribute Criteria}\label{feature-attribute-criteria-6}

\#

Attributes

Attribute Status

Phase

Digitizing Requirement

1

Program\_UID

Required

Feasibility

At feature creation

2

PROJECT\_ID

Required

Feasibility

At feature creation

3

Type

Required

Feasibility

At feature creation

4

Subtype

Required

Feasibility

At feature creation

5

RestorFeatStatus

Required

Design

At feature creation

6

sdsFeatureName

Optional

Design

As required and available

7

sdsFeatureDescription

Optional

Design

As required and available

8

RestorFeatVdatum

Optional

Design

As required and available

9

NumberPumps

Optional

Design

As required and available

10

PumpConfiguration

Optional

Construction

As required and available

11

PumpDrive

Optional

Construction

As required and available

12

PumpType

Optional

Construction

As required and available

13

MaxDesignHeadFT

Optional

Construction

As required and available

14

PumpCapacityCFS

Optional

Construction

As required and available

15

PumpElevation

Optional

Construction

As required and available

16

InvertElevation

Optional

Construction

As required and available

17

PumpAgeYRS

Optional

Construction

As required and available

18

DateBuilt

Optional

Construction

As required and available

19

Station

Optional

Construction

As required and available

20

sdsMetaDataID

Optional

Construction

As required and available

21

RestorFeatIDPK

Optional

Construction

As required and available

22

sdsID

Optional

Construction

As required and available

23

mediaIDFK

Optional

Construction

As required and available

\begin{center}\rule{0.5\linewidth}{0.5pt}\end{center}

\section{Parameters}\label{parameters-6}

\subsection{Program\_UID}\label{program_uid-6}

Designate program NESP or UMRR.

\subsection{Project\_ID (alias:
PROJECT\_ID)}\label{project_id-alias-project_id-6}

Foreign key to the unique NESP/UMRR project ID.

\subsection{Type (alias)}\label{type-alias-5}

A broad classification used to group NESP/UMRR restoration features into
categories of similar function.

\begin{quote}
\textbf{Note:} The Type/Subtype section below appears to have been
copied from PumpStation.\\
For IslandLine, you will likely want to change these to
\texttt{Island\ Creation} and appropriate island subtypes to match the
polygon Island feature class. They are shown here as provided.
\end{quote}

Value

Description

Water Level Management

Measures taken to imitate natural low flow river conditions with the
goal of improving wetland habitats.

\subsection{SubType (alias: Subtype)}\label{subtype-alias-subtype-6}

A classification used to group NESP/UMRR restoration features into
specific design feature categories.

Value

Description

Pump Station

Water pump infrastructure, comprised of a pad, pump, and pump enclosure,
used to introduce or remove water from a MSMU or backwater lake.

\subsection{RestorFeatStatus}\label{restorfeatstatus-6}

A description of the status of the individual restoration. It also
includes the party responsible for the individual feature.

Value

Description

USACE Planned

Feature is in a USACE planning stage.

USACE Built

Feature has been built by USACE. If a feature has been planned and built
by USACE, once constructed, it will always be considered USACE built
even if it has been repaired under later project stages.

USACE Enhanced

The function of an existing feature that was not built by USACE has been
enhanced by a feature that was USACE built or repaired.

USACE Repaired

An existing feature that was not built by USACE has been directly
improved or repaired.

Other Planned

A feature that is planned by another organization or government agency.

Other Built

Feature has been built by another organization or government agency and
is a integral part of the UMRR project.

Other Enhanced

The function of an existing feature, that was built by another
organization or government agency, has been enhanced by a feature that
was built or repaired by a non-USACE organization or government agency.

Other Repaired

An existing feature that has been directly improved or repaired by
another organization or government agency.

\subsection{sdsFeatureName}\label{sdsfeaturename-6}

The common name of the feature.

\subsection{sdsFeatureDescription}\label{sdsfeaturedescription-6}

A narrative describing the feature.

\subsection{RestorFeatVdatum}\label{restorfeatvdatum-6}

Vertical datum of elevations recorded in the attribute table or used as
feature z values.

Value

Description

MSL 1912

Mean Sea Level of 1912.

NGVD 1929

National Geodetic Vertical Datum of 1929.

NAVD 1988

North American Vertical Datum of 1988.

Other

Other, note datum in sdsFeatureDescription.

Unknown

Unknown vertical datum.

\begin{center}\rule{0.5\linewidth}{0.5pt}\end{center}

\section{Pump-related Attributes (as currently
defined)}\label{pump-related-attributes-as-currently-defined}

\begin{quote}
These attributes are included because they appear in your provided
IslandLine schema, but they are typically associated with PumpStation.
If IslandLine is not intended to store pump information, these may be
removed or repurposed in your data model.
\end{quote}

\subsection{NumberPumps}\label{numberpumps}

Number of a specific type of pumps in the pump station.

\subsection{PumpConfiguration}\label{pumpconfiguration}

Pump configuration generally referring to a horizontally or vertically
oriented pump. Area available for pump placement usually dictates this.

Value

Description

Horizontal

Horizontal pump configuration.

Vertical Line Shaft

Vertical line shaft pump configuration.

Vertical Submersible

Vertical submersible pump configuration.

Unknown

Unknown pump configuration.

NA

Not applicable.

\subsection{PumpDrive}\label{pumpdrive}

The method used to drive the pump action.

Value

Description

Diesel engine

Diesel engine pump drive.

Electric motor

Electric motor pump drive.

Natural gas engine

Natural gas engine pump drive.

Unknown

Unknown.

NA

Not applicable.

\subsection{PumpType}\label{pumptype}

The common pump type.

Value

Description

Axial flow

Axial flow pump type.

Centrifugal

Centrifugal pump type.

Mixed flow

Mixed flow pump type.

Unknown

Unknown.

NA

Not applicable.

\subsection{MaxDesignHeadFT}\label{maxdesignheadft}

The water elevation difference corresponding to the maximum design head
of the pump (feet).

\subsection{PumpCapacityCFS}\label{pumpcapacitycfs}

Pumping capacity in cubic feet per second (if multiple pumps, this
should be the total capacity).

\subsection{PumpElevation}\label{pumpelevation}

The midpoint elevation of the pump impeller or propeller.

\subsection{InvertElevation}\label{invertelevation}

The lowest sill elevation (e.g., intake or discharge invert).

\subsection{PumpAgeYRS}\label{pumpageyrs}

Age of the pump in years.

\subsection{DateBuilt}\label{datebuilt}

Date the pump station (or installation) was built.

\subsection{Station}\label{station}

Stationing of the feature, typically referenced to a project baseline or
alignment (e.g., 10+00).

\subsection{Horsepower}\label{horsepower}

Horsepower of the pump.

\subsection{FlowRate}\label{flowrate}

Flow rate of the pump (GPM).

\subsection{PumpStationType}\label{pumpstationtype}

Type of pump station installation.

Value

Description

Permanent Pump Station

Permanent pump station.

Portable Pump

Portable pump station.

\subsection{sdsMetaDataID}\label{sdsmetadataid-6}

References attached metadata.

\subsection{RestorFeatIDPK}\label{restorfeatidpk-6}

Primary key used to uniquely identify the restoration feature record.

\subsection{sdsID}\label{sdsid-6}

Identifier used to link the feature to supporting SDS (spatial data
system) records or related datasets.

\subsection{mediaIDFK}\label{mediaidfk-6}

Foreign key used to reference associated media (e.g., photographs,
design drawings, as-built plans).

\chapter{Island}\label{island}

\section{Spatial \& Geometric
Definition}\label{spatial-geometric-definition-7}

\begin{itemize}
\item
  Geometry Type: Polygon (Restoration\_Island\_Polygon).
\item
  Spatial Delineation: Subaerial landmass polygon constructed above
  normal flat pool within broad, shallow river expanses.
\item
  Architectural Typology: Barrier islands, multi-island archipelagos,
  horseshoe bays, and interior spits.
\end{itemize}

\subsection{Primary Restoration \& Engineering
Objectives}\label{primary-restoration-engineering-objectives-6}

\begin{itemize}
\item
  Wind-Wave Fetch Reduction: Break wind-wave generation across shallow
  open bays, reducing bottom orbital wave resuspension and restoring
  water clarity for submersed aquatic vegetation (SAV).
\item
  Terrestrial Habitat Creation: Provide secure, predator-isolated
  nesting habitat for migratory waterfowl, shorebirds, turtles, and
  bottomland hardwood trees.
\item
  Sediment Trapping: Create sheltered leeward lagoons that promote fine
  sediment drop-out and organic soil building.
\end{itemize}

\subsection{Typical Cross-Sections, Dimensions \& Sizing
Criteria}\label{typical-cross-sections-dimensions-sizing-criteria-5}

\begin{itemize}
\item
  Crest Elevations: Tiered multi-level profiles:
\item
  High Ridges (5 to 10 ft above pool): Designed above the 10- to 25-year
  flood stage to preserve terrestrial flood refugia and bottomland
  hardwood mast forests.
\item
  Low Sills / Wetland Benches (1 to 3 ft above pool): Frequently
  overtopped to sustain emergent marsh fringes.
\item
  Base Widths \& Footprints: Base widths of 100 to 400 ft; total
  individual footprints ranging from 2 to 50+ acres.
\item
  Side Slopes: Windward wave-impact faces sloped at 1V:3H to 1V:4H;
  leeward slopes shaped gently at 1V:8H to 1V:20H down to shallow
  littoral shelves.
\end{itemize}

\subsection{Construction \& Implementation
Methodology}\label{construction-implementation-methodology-7}

\begin{itemize}
\item
  Material Sourcing \& Sequencing: Fine to medium granular dredged sand
  mechanically or hydraulically placed to form the structural core.
\item
  Cohesive Topsoil Caps: Capping the sand core with a 1.0- to 2.0-ft
  layer of fine-grained cohesive dredge silt/clay to support plant
  rooting and resist wind deflation.
\item
  Windward Revetment Placement: Placing non-woven geotextile and Class
  II--III riprap along exposed upstream and windward shorelines
  concurrently with fill operations.
\end{itemize}

\subsection{Quality Control, Inspection \& Long-Term
Maintenance}\label{quality-control-inspection-long-term-maintenance-7}

\begin{itemize}
\item
  Topographic Lidar \& RTK Surveys: Post-construction as-built terrain
  modeling to verify settled elevations against contract cut/fill
  balances.
\item
  Shoreline Stability Audits: Biannual visual inspections to detect
  riprap displacement, wave notch cutting, or flank erosion following
  major spring river floods.
\end{itemize}

\section{Spatial Feature Attributes}\label{spatial-feature-attributes-1}

\subsection{Definition}\label{definition-7}

Island features are engineered or restored landforms within a river or
backwater area, constructed or modified using a variety of techniques
and materials to provide physical barriers to flow, wind, and waves, and
to support habitat and geomorphic functions.

\subsection{Purpose}\label{purpose-9}

The loss of islands in the Upper Mississippi River is primarily due to
human-induced alterations to the river's natural hydrology and sediment
dynamics. Dams, channelization, sediment loss, and altered floodplain
connectivity have disrupted the processes that create and maintain
islands. As a result, islands are increasingly vulnerable to erosion,
lack of sediment deposition, and other environmental stresses, leading
to their gradual disappearance. This loss impacts biodiversity and the
overall health of the river ecosystem, emphasizing the need for
restoration and management efforts to protect and rebuild these
important habitats. Island creation serves to restore structural
diversity, reduce wave and current energy, improve sediment processes,
and provide critical habitat for fish, wildlife, and plants.

\subsection{Applicability}\label{applicability-7}

Island creation serves to restore and enhance river ecosystems,
providing diverse habitats for wildlife, improving sediment and water
quality processes, supporting floodplain and aquatic species health, and
promoting biodiversity. Islands can be incorporated into restoration
designs to protect shorelines, create sheltered backwaters, influence
flow paths, and support vegetation and nesting habitat. (UMRR Design
Handbook, 2012)

\subsection{Notes}\label{notes-9}

Restoration features now have a Program\_UID field to distinguish
between NESP and UMRR features. Aliases without program identification
will be created for any field that currently has UMRR in the name, so
they can be used for NESP or UMRR data.

Attribute criteria for each field includes the following:

\begin{itemize}
\tightlist
\item
  Attribute Status -- defines whether the field requires a value or if
  it is optional based on available information.
\item
  Phases -- indicates the earliest time in the project lifecycle values
  for this attribute are available.
\item
  Digitizing requirements -- indicates when the field value should be
  entered in the attribute table.
\end{itemize}

Geometry Note: Polygon nodes should have z values; elevations should be
taken from as-built survey or design documents.

\subsection{Island}\label{island-1}

(polygon) -- Island structures engineered with a variety of techniques
and materials that serve as barriers to flow, wind, and waves.

\includegraphics[width=0.75\linewidth,height=\textheight,keepaspectratio]{images/island.jpg}

\begin{center}\rule{0.5\linewidth}{0.5pt}\end{center}

\section{Feature Attribute Criteria}\label{feature-attribute-criteria-7}

\#

Attributes

Attribute Status

Phase

Digitizing Requirement

1

Program\_UID

Required

Feasibility

At feature creation

2

PROJECT\_ID

Required

Feasibility

At feature creation

3

Type

Required

Feasibility

At feature creation

4

Subtype

Required

Feasibility

At feature creation

5

RestorFeatStatus

Required

Design

At feature creation

6

sdsFeatureName

Optional

Design

As required and available

7

sdsFeatureDescription

Optional

Design

As required and available

8

RestorFeatVdatum

Optional

Design

As required and available

9

WaterOvertopping

Optional

Design

As required and available

10

PrimaryMaterial

Optional

Construction

As required and available

11

SecondaryMaterial

Optional

Construction

As required and available

12

DredgeMaterialPlacement

Optional

Construction

As required and available

13

Acres

Optional

Construction

As required and available

14

sdsMetaDataID

Optional

Construction

As required and available

15

RestorFeatIDPK

Optional

Construction

As required and available

16

sdsID

Optional

Construction

As required and available

17

mediaIDFK

Optional

Construction

As required and available

\begin{center}\rule{0.5\linewidth}{0.5pt}\end{center}

\section{Parameters}\label{parameters-7}

\subsection{Program\_UID}\label{program_uid-7}

Designate program NESP or UMRR.

\subsection{Project\_ID (alias:
PROJECT\_ID)}\label{project_id-alias-project_id-7}

Foreign key to the unique NESP/UMRR project ID.

\subsection{Type (alias)}\label{type-alias-6}

A broad classification used to group NESP/UMRR restoration features into
categories of similar function.

Value

Description

Island Creation

Measures taken to construct islands in rivers.

\subsection{SubType (alias: Subtype)}\label{subtype-alias-subtype-7}

A classification used to group NESP/UMRR restoration features into
specific island design feature categories.

Value

Description

Barrier island

Island structures engineered with a variety of techniques and materials
that serve as barriers to flow, wind, and waves.

Seed island

Rock or log structures placed perpendicular to currents in areas where
coarse sediment transport is occurring, used to create islands via
sediment accumulation due to reduced velocities.

Nourished seed island

Rock or log structures perpendicular to currents with sand placed on the
upstream or downstream side; currents reshape sand naturally, used where
coarse sediment transport is low.

Rock sill

Short, steep ramp comprised of boulders created to circumvent in-channel
barriers to natural fish movement.

Sand tip

Low structures overtopped more frequently than islands, used to increase
flow into areas during floods or to maintain flood flow into secondary
channels; may include geotechnical membrane in the upstream sill to
reduce seepage.

Sand flat

Areas of sand along islands precisely graded to be seasonally inundated
to very shallow depths (0.1--0.2 ft) to create microtopography and
habitat diversity.

Loafing structure

Large woody debris anchored into rock structures or placed along
shorelines to provide loafing habitat for waterfowl, turtles, and other
fauna, as well as fish habitat and enhanced shoreline aesthetics.

Rock/Log island

Structures constructed using rock to anchor large woody debris in the
river channel, providing habitat diversity and loafing structures while
improving aesthetics.

\subsection{RestorFeatStatus}\label{restorfeatstatus-7}

A description of the status of the individual restoration. It also
includes the party responsible for the individual feature.

Value

Description

USACE Planned

Feature is in a USACE planning stage.

USACE Built

Feature has been built by USACE. If a feature has been planned and built
by USACE, once constructed, it will always be considered USACE built
even if it has been repaired under later project stages.

USACE Enhanced

The function of an existing feature that was not built by USACE has been
enhanced by a feature that was USACE built or repaired.

USACE Repaired

An existing feature that was not built by USACE has been directly
improved or repaired.

Other Planned

A feature that is planned by another organization or government agency.

Other Built

Feature has been built by another organization or government agency and
is an integral part of the UMRR project.

Other Enhanced

The function of an existing feature, built by another organization or
government agency, has been enhanced by a feature that was built or
repaired by a non-USACE organization or government agency.

Other Repaired

An existing feature that has been directly improved or repaired by
another organization or government agency.

\subsection{sdsFeatureName}\label{sdsfeaturename-7}

The common name of the feature.

\subsection{sdsFeatureDescription}\label{sdsfeaturedescription-7}

A narrative describing the feature.

\subsection{RestorFeatVdatum}\label{restorfeatvdatum-7}

Vertical datum of elevations recorded in the attribute table or used as
feature z values.

Value

Description

MSL 1912

Mean Sea Level of 1912.

NGVD 1929

National Geodetic Vertical Datum of 1929.

NAVD 1988

North American Vertical Datum of 1988.

Other

Other; note datum in sdsFeatureDescription.

Unknown

Unknown vertical datum.

\subsection{WaterOvertopping}\label{waterovertopping}

Description of the relationship between the island feature and water
level.

Value

Description

AlwaysDry

Not covered at high water under average meteorological conditions.

AlwaysSubmerged

Remains covered by water at all times under average meteorological
conditions.

Awash

Flush with or washed by waves at low water under average meteorological
conditions.

CoversAndUncovers

Projection from the bottom that periodically extends above and is
submerged below the water surface.

Floating

Resting or moving on the surface of a water body.

PartlySubmergedAtHighWater

Partially covered and partially dry at high water.

SubjectToInundationOrFlooding

Area or structure periodically covered by flood waters.

Unknown

Relationship to water level is unknown.

NA

Not applicable.

\subsection{PrimaryMaterial}\label{primarymaterial}

Primary construction material used for the island feature.

Value

Description

Sand

Sand.

Rock

Rock.

Fines

Fine sediment including clay and clay mixtures.

Timber/Wood

Timber or wood, including logs or root wads.

Earthen

Mixture of sand and fines.

Concrete

Concrete.

\subsection{SecondaryMaterial}\label{secondarymaterial}

Secondary construction material, if applicable.

Value

Description

Sand

Sand.

Rock

Rock.

Fines

Fine sediment including clay and clay mixtures.

Timber/Wood

Timber or wood, including logs or root wads.

Earthen

Mixture of sand and fines.

Concrete

Concrete.

\subsection{DredgeMaterialPlacement}\label{dredgematerialplacement}

Describes whether dredged material was placed to create or augment the
island feature.

Value

Description

Yes

Dredged material was placed to create or augment the feature.

No

No dredged material was used.

NA

Not applicable.

Unknown

Use of dredged material is unknown.

\subsection{Acres}\label{acres}

Acreage of the island feature.

\subsection{sdsMetaDataID}\label{sdsmetadataid-7}

References attached metadata.

\subsection{RestorFeatIDPK}\label{restorfeatidpk-7}

Primary key used to uniquely identify the restoration feature record.

\subsection{sdsID}\label{sdsid-7}

Identifier used to link the feature to supporting SDS (spatial data
system) records or related datasets.

\subsection{mediaIDFK}\label{mediaidfk-7}

Foreign key used to reference associated media (e.g., photographs,
design drawings, construction records, as-built plans) documenting the
island feature.

\chapter{Moist Soil Management Unit}\label{moist-soil-management-unit}

\section{Spatial \& Geometric
Definition}\label{spatial-geometric-definition-8}

\begin{itemize}
\item
  Geometry Type: Polygon (MSU\_Cell\_Boundary).
\item
  Spatial Delineation: Enclosed, impounded floodplain basin bounded by
  perimeter containment embankments, internal ditches, and water control
  structures.
\item
  Ecological Classification: Shallow seasonally flooded palustrine
  emergent wetland complex.
\end{itemize}

\subsection{Primary Restoration \& Engineering
Objectives}\label{primary-restoration-engineering-objectives-7}

\begin{itemize}
\item
  Targeted Seed \& Tuber Production: Provide precise, independent
  hydrologic drawdowns and reflooding to stimulate high-energy forage
  plants (smartweeds, wild millets, sedges, sprangletop).
\item
  Migratory Waterfowl Refueling: Provide shallow foraging water depths
  (2.0 to 8.0 in) timed precisely to match spring and autumn waterfowl
  migration pulses.
\item
  Shorebird Staging Flats: Expose extensive mudflats during late summer
  (July--August) for migrating northbound and southbound shorebirds.
\end{itemize}

\section{Typical Cross-Sections, Dimensions \& Sizing
Criteria}\label{typical-cross-sections-dimensions-sizing-criteria-6}

\begin{itemize}
\item
  Cell Surface Areas: Individual management cells typically scale from
  10 to 80 acres.
\item
  Topographic Grading: Graded with gentle interior slopes
  (\textless0.1\% grade) toward a central drainage swale network
  (bottomwidths10 to 20 ft) to ensure complete gravity drainage without
  water pocketing.
\item
  Operating Depth Envelope: Maximum impounded water depth of 1.5 to 3.0
  ft; drawdowns expose 100\% of cell floor to atmospheric drying.
\end{itemize}

\subsection{Construction \& Implementation
Methodology}\label{construction-implementation-methodology-8}

\begin{itemize}
\item
  Precision Earthmoving: Laser-guided bulldozers and low-ground-pressure
  earthmovers shaping interior swales and cell bottom micro-topography.
\item
  Clay Liner Compaction: Where underlying alluvium contains excessive
  sand lenses, interior floors are sealed with a compacted 1.0- to
  2.0-ft cohesive clay blanket to prevent downward water loss.
\end{itemize}

\subsection{Quality Control, Inspection \& Long-Term
Maintenance}\label{quality-control-inspection-long-term-maintenance-8}

\begin{itemize}
\item
  Annual Water Management Plans: Rigorous hydro-period cycling: slow
  spring drawdowns (March--May), summer soil drying and germination
  (June--August), and gradual fall re-flooding (September--November).
\item
  Discing \& Succession Setback: Rotational tractor discing or
  controlled burning every 3 to 5 years to reset plant succession and
  eliminate monotypic willow or cattail stands.
\end{itemize}

Spatial Feature Attributes \#\#\# Definition Areas or units designed to
enhance wetland habitat through the use of water containment, water
supply, and water control structures.

\subsection{Purpose}\label{purpose-10}

The basic operating plan for a Moist Soil Management Unit (MSMU) is to
keep water out in the late spring and summer and to gradually flood the
area in the fall. In a multiple-cell system, independent control of
water levels in each cell is preferred to allow targeted habitat
management.

\subsection{Applicability}\label{applicability-8}

Moist Soil Management Units are typically designed to include water
containment, water supply, and water control structures. Water
containment is provided by construction of exterior berms, interior
berms, and overflow spillways, which are used to impound water during
seasonal waterfowl migrations or to keep water out of the impounded
area. Water supply may be provided by either river water or groundwater
through the use of a pump station or well, respectively. Water control
structures are utilized to maintain desired water elevations throughout
the year. (UMRR Design Handbook, 2012)

\subsection{Notes}\label{notes-10}

Restoration features now have a Program\_UID field to distinguish
between NESP and UMRR features. Aliases without program identification
will be created for any field that currently has UMRR in the name, so
they can be used for NESP or UMRR data.

Attribute criteria for each field includes the following:

\begin{itemize}
\tightlist
\item
  Attribute Status -- defines whether the field requires a value or if
  it is optional based on available information.
\item
  Phases -- indicates the earliest time in the project lifecycle values
  for this attribute are available.
\item
  Digitizing requirements -- indicates when the field value should be
  entered in the attribute table.
\end{itemize}

Geometry Note: No Z values are needed for this feature.

\subsection{MoistSoilMgmtUnit}\label{moistsoilmgmtunit}

(polygon) -- Area or area subdivision that is managed to create moist
soil units.

\begin{center}\rule{0.5\linewidth}{0.5pt}\end{center}

\section{Feature Attribute Criteria}\label{feature-attribute-criteria-8}

\#

Attribute

Attribute Status

Phase

Digitizing Requirements

1

Program\_UID

Required

Feasibility

At feature creation

2

PROJECT\_ID (Alias)

Required

Feasibility

At feature creation

3

Type (Alias)

Required

Feasibility

At feature creation

4

Subtype (Alias)

Required

Feasibility

At feature creation

5

RestorFeatStatus

Required

Feasibility

At feature creation

6

sdsFeatureName

Optional

Design

As required and available

7

sdsFeatureDescription

Optional

Design

As required and available

8

RestorFeatVdatum

Optional

Design

As required and available

9

Acres

Optional

Design

As required and available

10

Depth

Optional

Design

As required and available

11

WaterControlType

Optional

Design

As required and available

12

WaterControlMethod

Optional

Design

As required and available

13

sdsMetaDataID

Optional

Construction

As required and available

14

RestorFeatIDPK

Optional

Construction

As required and available

15

sdsID

Optional

Construction

As required and available

16

mediaIDFK

Optional

Construction

As required and available

\begin{center}\rule{0.5\linewidth}{0.5pt}\end{center}

\section{Parameters}\label{parameters-8}

\subsection{Program\_UID}\label{program_uid-8}

Designate program NESP or UMRR.

\subsection{Project\_ID (alias:
PROJECT\_ID)}\label{project_id-alias-project_id-8}

Foreign key to the unique NESP/UMRR project ID.

\subsection{Type (alias)}\label{type-alias-7}

A broad classification used to group NESP/UMRR restoration features into
categories of similar function.

Value

Description

Water Level Management

Measures taken to manage and manipulate water levels with the goal of
improving wetland habitats.

\subsection{SubType (alias: Subtype)}\label{subtype-alias-subtype-8}

A classification used to group NESP/UMRR restoration features into
specific design feature categories.

Value

Description

Moist soil management unit

Units designed to enhance wetland habitat through use of water
containment, water supply, and water control structures.

Backwater lake

Units designed to enhance aquatic habitat through use of water
containment, water supply, and water control structures.

Cell

Sub-unit of a moist soil unit, bounded by exterior and/or interior
berms, allowing targeted water level management within a MSMU.

\subsection{RestorFeatStatus}\label{restorfeatstatus-8}

A description of the status of the individual restoration. It also
includes the party responsible for the individual feature.

Value

Description

USACE Planned

Feature is in a USACE planning stage.

USACE Built

Feature has been built by USACE. If a feature has been planned and built
by USACE, once constructed, it will always be considered USACE built
even if repaired under later project stages.

USACE Enhanced

The function of an existing feature not built by USACE has been enhanced
by a feature that was USACE built or repaired.

USACE Repaired

An existing feature not built by USACE has been directly improved or
repaired.

Other Planned

A feature that is planned by another organization or government agency.

Other Built

Feature has been built by another organization or government agency and
is an integral part of the UMRR project.

Other Enhanced

The function of an existing feature, built by another organization or
government agency, has been enhanced by a feature that was built or
repaired by a non-USACE organization or government agency.

Other Repaired

An existing feature that has been directly improved or repaired by
another organization or government agency.

\subsection{sdsFeatureName}\label{sdsfeaturename-8}

The common name of the feature.

\subsection{sdsFeatureDescription}\label{sdsfeaturedescription-8}

A narrative describing the feature.

\subsection{RestorFeatVdatum}\label{restorfeatvdatum-8}

Vertical datum of elevations recorded in the attribute table or used as
feature z values.

Value

Description

MSL 1912

Mean Sea Level of 1912.

NGVD 1929

National Geodetic Vertical Datum of 1929.

NAVD 1988

North American Vertical Datum of 1988.

Other

Other; note datum in sdsFeatureDescription.

Unknown

Unknown vertical datum.

\subsection{Acres}\label{acres-1}

Acreage of the feature.

\subsection{Depth}\label{depth}

Typical or design depth of the feature (e.g., target water depth),
recorded in appropriate units (commonly feet).

\subsection{WaterControlType}\label{watercontroltype}

General classification of how water control is achieved.

Value

Description

Passive

Water control depends on natural processes (e.g., seasonal inundation)
or passive engineering structures (e.g., spillway or notch).

Active

Water control depends on human actions (e.g., pump station operation,
stoplog structures).

\subsection{WaterControlMethod}\label{watercontrolmethod}

Specific water control technique used.

Value

Description

Overflow spillway

Section of embankment engineered to provide controlled overtopping of an
embankment system, typically at the downstream end of a site at an
elevation lower than the exterior berm.

Pump

Water pump infrastructure (pad, pump, enclosure) used to introduce or
remove water from a MSMU or backwater lake.

Well

Vertical hole drilled into the ground to access or influence groundwater
movement, used to introduce or remove water from a MSMU or backwater
lake.

Stoplog

Water management structure that allows control of water level or flow
rate into or out of MSMUs and backwater lakes by adding or removing
stoplogs.

Other

Other method; specify in sdsFeatureDescription.

\subsection{sdsMetaDataID}\label{sdsmetadataid-8}

References attached metadata.

\subsection{RestorFeatIDPK}\label{restorfeatidpk-8}

Primary key used to uniquely identify the restoration feature record.

\subsection{sdsID}\label{sdsid-8}

Identifier used to link the feature to supporting SDS (spatial data
system) records or related datasets.

\subsection{mediaIDFK}\label{mediaidfk-8}

Foreign key used to reference associated media (e.g., photographs,
design drawings, operating plans, as-built plans) documenting the moist
soil management unit.

\chapter{Notch}\label{notch}

\section{Spatial \& Geometric
Definition}\label{spatial-geometric-definition-9}

\begin{itemize}
\item
  Geometry Type: Polygon or Point (Notch\_Geometry).
\item
  Spatial Delineation: Geometric cut, depression, or gap engineered into
  rock training structures, wing dams, closing dams, or earthen levee
  crests.
\item
  Cross-Sectional Typology: Trapezoidal, parabolic, or rectangular weir
  cuts.
\end{itemize}

\subsection{Primary Restoration \& Engineering
Objectives}\label{primary-restoration-engineering-objectives-8}

\begin{itemize}
\item
  Secondary Channel Revitalization: Pass controlled baseflow through
  closing dams to refresh side channels without destabilizing
  main-channel navigation drafts.
\item
  Bathymetric Scour Generation: Concentrate localized flow vectors
  through the notch throat, carving deep downstream plunge pools (10 to
  25 ft deep) that function as prime overwintering fish holes.
\item
  Sediment Purging: Re-energize sediment transport capacity through
  backwaters to purge deposited silt plugs.
\end{itemize}

\subsection{Typical Cross-Sections, Dimensions \& Sizing
Criteria}\label{typical-cross-sections-dimensions-sizing-criteria-7}

\begin{itemize}
\item
  Throat Widths: Bottom widths typically range from 10 to 50 ft; side
  slopes cut at 1V:2H to 1V:3H.
\item
  Depth of Invert: Invert set 1.0 to 4.0 ft below the parent weir crest
  or normal flat pool elevation.
\item
  Hydraulic Conveyance: Sized using broad-crested weir formulas (Q=C d
  LH 3/2 ) to restrict diverted river discharge to calibrated project
  thresholds (typically \textless5\% of total channel discharge).
\end{itemize}

\subsection{Construction \& Implementation
Methodology}\label{construction-implementation-methodology-9}

\begin{itemize}
\item
  Marine Excavation: Existing stone removed using barge-mounted
  hydraulic excavators equipped with rock buckets or heavy clamshells.
\item
  Bed \& Abutment Armoring: Installing heavy riprap transitions along
  notch edges and lining the downstream plunge impact zone with Class
  IV--V stone to halt headcut unraveling.
\end{itemize}

\subsection{Quality Control, Inspection \& Long-Term
Maintenance}\label{quality-control-inspection-long-term-maintenance-9}

\begin{itemize}
\tightlist
\item
  Single-Beam Sonar Soundings: Annual hydrographic cross-sections
  through the notch and downstream scour basin to verify that scour hole
  depth remains stable and does not undermine the main weir abutments.
\end{itemize}

\section{Spatial Feature Attribute}\label{spatial-feature-attribute-6}

\subsection{Definition}\label{definition-8}

A notch is a deliberate opening cut into a riverine flow structure or
levee (typically a dike) used to control the flow of water into and out
of a low-flow area or adjacent floodplain.

\subsection{Purpose}\label{purpose-11}

A notch creates a controlled gap or opening in the structure. Its
purpose is to regulate water flow, manage river levels, and prevent
damage from excessive water buildup. Notching can support flow
regulation, sediment management, habitat creation, and flood control by
selectively passing water and sediment and reconnecting channels and
floodplains.

\subsection{Applicability}\label{applicability-9}

Notches are applicable where targeted hydraulic connections are needed
between the main channel and backwaters, side channels, or floodplain
areas. They are commonly used to improve habitat, reduce sedimentation,
restore natural inundation patterns, and modify river training
structures or levees to better meet environmental objectives within
Upper Mississippi River System (UMRS) restoration projects. (UMRR Design
Handbook, 2012)

\subsection{Notes}\label{notes-11}

Restoration features now have a Program\_UID field to distinguish
between NESP and UMRR features. Aliases without program identification
will be created for any field that currently has UMRR in the name, so
they can be used for NESP or UMRR data.

Attribute criteria for each field includes the following:

\begin{itemize}
\tightlist
\item
  Attribute Status -- defines whether the field requires a value or if
  it is optional based on available information.
\item
  Phases -- indicates the earliest time in the project lifecycle values
  for this attribute are available.
\item
  Digitizing requirements -- indicates when the field value should be
  entered in the attribute table.
\end{itemize}

Geometry Note: No Z values are needed for this feature.

\subsection{Notch}\label{notch-1}

(point) -- A notch in a riverine flow structure or levee used to control
the flow of water into and out of the low-flow area.

\begin{center}\rule{0.5\linewidth}{0.5pt}\end{center}

\section{Feature Attribute Criteria}\label{feature-attribute-criteria-9}

\#

Attribute

Attribute Status

Phase

Digitizing Requirements

1

Program\_UID

Required

Feasibility

At feature creation

2

PROJECT\_ID (Alias)

Required

Feasibility

At feature creation

3

Type (Alias)

Required

Feasibility

At feature creation

4

Subtype (Alias)

Required

Feasibility

At feature creation

5

RestorFeatStatus

Required

Feasibility

At feature creation

6

sdsFeatureName

Optional

Design

As required and available

7

sdsFeatureDescription

Optional

Design

As required and available

8

RestorFeatVdatum

Optional

Design

As required and available

9

ConstructionDate

Optional

Design

As required and available

10

NotchElevation

Optional

Design

As required and available

11

NotchWidthFT

Optional

Design

As required and available

12

Station

Optional

Construction

As required and available

13

sdsMetaDataID

Optional

Construction

As required and available

14

RestorFeatIDPK

Optional

Construction

As required and available

15

sdsID

Optional

Construction

As required and available

16

mediaIDFK

Optional

Construction

As required and available

\begin{center}\rule{0.5\linewidth}{0.5pt}\end{center}

\section{Parameters}\label{parameters-9}

\subsection{Program\_UID}\label{program_uid-9}

Designate program NESP or UMRR.

\subsection{Project\_ID (alias:
PROJECT\_ID)}\label{project_id-alias-project_id-9}

Foreign key to the unique NESP/UMRR project ID.

\subsection{Type (alias)}\label{type-alias-8}

A broad classification used to group NESP/UMRR restoration features into
categories of similar function.

Value

Description

Channel Modifications

Measures taken to train the main channel of a river and to modify
secondary channels.

Floodplain Restoration

Measures taken to improve natural floodplain function and habitat.

\subsection{SubType (alias: Subtype)}\label{subtype-alias-subtype-9}

A classification used to group NESP/UMRR restoration features into
specific notch design feature categories.

Value

Description

Wing dam notch

Wing dam modification that significantly lowers a section of the dam to
support diverse habitats while maintaining navigation channel function.

Closure structure notch

Closure structure modification that significantly lowers a section of
the structure, allowing decreased sedimentation and enhanced habitat.

Dike notch

Dike modification that significantly lowers a section of the structure,
allowing decreased sedimentation and enhanced habitat.

Levee notch

A low spot excavated in a frontline levee to allow inundation of the
area behind the levee above a specific flow.

\subsection{RestorFeatStatus}\label{restorfeatstatus-9}

A description of the status of the individual restoration. It also
includes the party responsible for the individual feature.

Value

Description

USACE Planned

Feature is in a USACE planning stage.

USACE Built

Feature has been built by USACE. If a feature has been planned and built
by USACE, once constructed, it will always be considered USACE built
even if repaired under later project stages.

USACE Enhanced

The function of an existing feature not built by USACE has been enhanced
by a USACE-built or repaired feature.

USACE Repaired

An existing feature not built by USACE has been directly improved or
repaired.

Other Planned

Feature planned by another organization or government agency.

Other Built

Feature built by another organization or government agency and integral
to the UMRR project.

Other Enhanced

Function of a feature built by another organization has been enhanced by
a feature built or repaired by a non-USACE organization or government
agency.

Other Repaired

Existing feature directly improved or repaired by another organization
or government agency.

\subsection{sdsFeatureName}\label{sdsfeaturename-9}

The common name of the feature.

\subsection{sdsFeatureDescription}\label{sdsfeaturedescription-9}

A narrative describing the feature.

\subsection{RestorFeatVdatum}\label{restorfeatvdatum-9}

Vertical datum of elevations recorded in the attribute table or used as
feature z values.

Value

Description

MSL 1912

Mean Sea Level of 1912.

NGVD 1929

National Geodetic Vertical Datum of 1929.

NAVD 1988

North American Vertical Datum of 1988.

Other

Other; note datum in sdsFeatureDescription.

Unknown

Unknown vertical datum.

\subsection{ConstructionDate}\label{constructiondate}

The date of construction (or reconstruction) of the notch.

\subsection{NotchElevation}\label{notchelevation}

Elevation of the bottom of the notch, referenced to the vertical datum
specified in RestorFeatVdatum.

\subsection{NotchWidthFT}\label{notchwidthft}

Width of the notch in feet.

\subsection{Station}\label{station-1}

Stationing of the feature, typically referenced to a project baseline or
alignment (e.g., 10+00).

\subsection{sdsMetaDataID}\label{sdsmetadataid-9}

References attached metadata.

\subsection{RestorFeatIDPK}\label{restorfeatidpk-9}

Primary key used to uniquely identify the restoration feature record.

\subsection{sdsID}\label{sdsid-9}

Identifier used to link the feature to supporting SDS (spatial data
system) records or related datasets.

\subsection{mediaIDFK}\label{mediaidfk-9}

Foreign key used to reference associated media (e.g., design drawings,
construction photos, as-built plans) documenting the notch.

\chapter{Pump Station}\label{pump-station}

\section{Spatial \& Geometric
Definition}\label{spatial-geometric-definition-10}

\begin{itemize}
\item
  Geometry Type: Point or Polygon (Pump\_Station\_Structure).
\item
  Spatial Delineation: Footprint encompassing cast-in-place concrete
  wet-wells, intake bays, discharge manifolds, electrical control
  buildings, and fuel/power infrastructure.
\item
  Operational Classification: High-capacity, low-head axial-flow pumping
  systems.
\end{itemize}

\subsection{Primary Restoration \& Engineering
Objectives}\label{primary-restoration-engineering-objectives-9}

\begin{itemize}
\item
  Hydrologic Independence: Allow rapid filling or complete dewatering of
  moist soil units and managed wetland cells regardless of ambient river
  flood or drought stages.
\item
  Dewatering Reliability: Drain impounded areas during spring flood
  stages when external river levels sit higher than the internal wetland
  ground elevation.
\item
  Rapid Flood Protection: Evacuate internal runoff following heavy storm
  events to prevent premature drown-out of maturing wetland seed crops.
\end{itemize}

\subsection{Typical Mechanical \& Structural Engineering
Specifications}\label{typical-mechanical-structural-engineering-specifications}

\begin{itemize}
\item
  Pumping Capacity: Total station flow capacities range between 2,000
  and 25,000 gpm (5 to 55 cfs), operating against total dynamic heads
  (TDH) of 6 to 18 ft.
\item
  Pump Mechanics: Vertical axial-flow or mixed-flow electric submersible
  or lineshaft pumps. Sized with variable frequency drives (VFD) for
  adjustable flow rates.
\item
  Reversible Manifold Piping: Multi-valve directional piping manifolds
  allowing a single pump train to draw water from the river to fill the
  cell, or draw from the cell to discharge out into the river.
\end{itemize}

\subsection{Construction \& Implementation
Methodology}\label{construction-implementation-methodology-10}

\begin{itemize}
\item
  Deep Foundation Systems: Reinforced concrete substructure founded on
  steel H-piling or driven ductile iron piles keyed into underlying
  load-bearing river sands to resist uplift hydrostatic pressures.
\item
  Cofferdam Dewatering: Wet-well construction executed within sheetpile
  cofferdams dewatered by deep-well point systems.
\end{itemize}

\subsection{Quality Control, Inspection \& Long-Term
Maintenance}\label{quality-control-inspection-long-term-maintenance-10}

\begin{itemize}
\item
  Intake Screening: Trash racks equipped with automated mechanical rakes
  or heavy bar screens (2- to 3-in clear spacing) to prevent log and
  debris ingestion.
\item
  Winterization: Annual late-autumn dewatering of pump bowls,
  lubrication checks, and cycling of emergency backup diesel generators
  under full load.
\end{itemize}

\section{Spatial Feature Attribute}\label{spatial-feature-attribute-7}

\subsection{Definition}\label{definition-9}

A pump station is a facility that pumps water between two storage areas,
a storage area and a river reach, or between different river reaches. It
can also be used to transfer water within the same 2D Flow Area or
between different cells within the same 2D Flow Area.

\subsection{Purpose}\label{purpose-12}

Pump stations are used to introduce or remove water from backwater
lakes, MSMUs, or other project areas using surface water (typically from
a river) or groundwater. They provide controlled water level management,
allowing project operators to manipulate hydrologic conditions to
support habitat objectives, flood risk management, and operational
needs.

\subsection{Applicability}\label{applicability-10}

Pumping stations are used for removal of interior drainage from areas
protected by local flood protection works and for water supply. Pumping
stations can and should have negligible adverse environmental impacts.
They can also be used to connect pumps, lift flows from lower to higher
elevations, increase pressure in gravity flow distribution systems, and
improve the selection and operation of water level management facilities
within restoration projects. (UMRR Design Handbook, 2012)

\subsection{Notes}\label{notes-12}

Restoration features now have a Program\_UID field to distinguish
between NESP and UMRR features. Aliases without program identification
will be created for any field that currently has UMRR in the name, so
they can be used for NESP or UMRR data.

Attribute criteria for each field includes the following:

\begin{itemize}
\tightlist
\item
  Attribute Status -- defines whether the field requires a value or if
  it is optional based on available information.
\item
  Phases -- indicates the earliest time in the project lifecycle values
  for this attribute are available.
\item
  Digitizing requirements -- indicates when the field value should be
  entered in the attribute table.
\end{itemize}

Geometry Note: No Z values are needed for this feature.

\subsection{PumpStation}\label{pumpstation}

(point) -- The location of a pump for the purpose of pumping water from
the interior of a levee or embankment.

\begin{center}\rule{0.5\linewidth}{0.5pt}\end{center}

\section{Feature Attribute
Criteria}\label{feature-attribute-criteria-10}

\#

Attributes

Attribute Status

Phase

Digitizing Requirement

1

Program\_UID

Required

Feasibility

At feature creation

2

PROJECT\_ID

Required

Feasibility

At feature creation

3

Type

Required

Feasibility

At feature creation

4

Subtype

Required

Feasibility

At feature creation

5

RestorFeatStatus

Required

Design

At feature creation

6

sdsFeatureName

Optional

Design

As required and available

7

sdsFeatureDescription

Optional

Design

As required and available

8

RestorFeatVdatum

Optional

Design

As required and available

9

NumberPumps

Optional

Design

As required and available

10

PumpConfiguration

Optional

Construction

As required and available

11

PumpDrive

Optional

Construction

As required and available

12

PumpType

Optional

Construction

As required and available

13

MaxDesignHeadFT

Optional

Construction

As required and available

14

PumpCapacityCFS

Optional

Construction

As required and available

15

PumpElevation

Optional

Construction

As required and available

16

InvertElevation

Optional

Construction

As required and available

17

PumpAgeYRS

Optional

Construction

As required and available

18

DateBuilt

Optional

Construction

As required and available

19

Station

Optional

Construction

As required and available

20

sdsMetaDataID

Optional

Construction

As required and available

21

RestorFeatIDPK

Optional

Construction

As required and available

22

sdsID

Optional

Construction

As required and available

23

mediaIDFK

Optional

Construction

As required and available

\begin{center}\rule{0.5\linewidth}{0.5pt}\end{center}

\section{Parameters}\label{parameters-10}

\subsection{Program\_UID}\label{program_uid-10}

Designate program NESP or UMRR.

\subsection{Project\_ID (alias:
PROJECT\_ID)}\label{project_id-alias-project_id-10}

Foreign key to the unique NESP/UMRR project ID.

\subsection{Type (alias)}\label{type-alias-9}

A broad classification used to group NESP/UMRR restoration features into
categories of similar function.

Value

Description

Water Level Management

Measures taken to imitate natural low flow river conditions with the
goal of improving wetland habitats.

\subsection{SubType (alias: Subtype)}\label{subtype-alias-subtype-10}

A classification used to group NESP/UMRR restoration features into
specific design feature categories.

Value

Description

Pump Station

Water pump infrastructure, which may include a pad, pump(s), and pump
enclosure, used to introduce or remove water from a MSMU or backwater
lake.

\subsection{RestorFeatStatus}\label{restorfeatstatus-10}

A description of the status of the individual restoration. It also
includes the party responsible for the individual feature.

Value

Description

USACE Planned

Feature is in a USACE planning stage.

USACE Built

Feature has been built by USACE. If a feature has been planned and built
by USACE, once constructed, it will always be considered USACE built
even if it has been repaired under later project stages.

USACE Enhanced

The function of an existing feature that was not built by USACE has been
enhanced by a feature that was USACE built or repaired.

USACE Repaired

An existing feature that was not built by USACE has been directly
improved or repaired.

Other Planned

A feature that is planned by another organization or government agency.

Other Built

Feature has been built by another organization or government agency and
is an integral part of the UMRR project.

Other Enhanced

The function of an existing feature, that was built by another
organization or government agency, has been enhanced by a feature that
was built or repaired by a non-USACE organization or government agency.

Other Repaired

An existing feature that has been directly improved or repaired by
another organization or government agency.

\subsection{sdsFeatureName}\label{sdsfeaturename-10}

The common name of the feature.

\subsection{sdsFeatureDescription}\label{sdsfeaturedescription-10}

A narrative describing the feature.

\subsection{RestorFeatVdatum}\label{restorfeatvdatum-10}

Vertical datum of elevations recorded in the attribute table or used as
feature z values.

Value

Description

MSL 1912

Mean Sea Level of 1912.

NGVD 1929

National Geodetic Vertical Datum of 1929.

NAVD 1988

North American Vertical Datum of 1988.

Other

Other; note datum in sdsFeatureDescription.

Unknown

Unknown vertical datum.

\subsection{NumberPumps}\label{numberpumps-1}

Number of pumps of a specific type in the pump station.

\subsection{PumpConfiguration}\label{pumpconfiguration-1}

Pump configuration, generally referring to a horizontally or vertically
oriented pump. The area available for pump placement usually dictates
configuration.

Value

Description

Horizontal

Horizontal pump configuration.

Vertical Line Shaft

Vertical line shaft pump configuration.

Vertical Submersible

Vertical submersible pump configuration.

Unknown

Unknown pump configuration.

NA

Not applicable.

\subsection{PumpDrive}\label{pumpdrive-1}

The method used to drive the pump action.

Value

Description

Diesel engine

Diesel engine pump drive.

Electric motor

Electric motor pump drive.

Natural gas engine

Natural gas engine pump drive.

Unknown

Unknown pump drive.

NA

Not applicable.

\subsection{PumpType}\label{pumptype-1}

The common pump type.

Value

Description

Axial flow

Axial flow pump type.

Centrifugal

Centrifugal pump type.

Mixed flow

Mixed flow pump type.

Unknown

Unknown pump type.

NA

Not applicable.

\subsection{MaxDesignHeadFT}\label{maxdesignheadft-1}

The maximum design head of the pump, expressed as the water surface
elevation difference (in feet) between suction and discharge at design
conditions.

\subsection{PumpCapacityCFS}\label{pumpcapacitycfs-1}

Pumping capacity of the station in cubic feet per second (CFS). If
multiple pumps are present, this should represent the total station
capacity.

\subsection{PumpElevation}\label{pumpelevation-1}

Midpoint elevation of the pump impeller or propeller, referenced to
RestorFeatVdatum.

\subsection{InvertElevation}\label{invertelevation-1}

Lowest sill or invert elevation of the intake or discharge conduit,
referenced to RestorFeatVdatum.

\subsection{PumpAgeYRS}\label{pumpageyrs-1}

Age of the pump in years (may be derived from DateBuilt or last
replacement date).

\subsection{DateBuilt}\label{datebuilt-1}

Date the pump station (or pump installation) was built or substantially
completed.

\subsection{Station}\label{station-2}

Stationing of the feature, typically referenced to a project baseline or
alignment (e.g., 10+00).

\subsection{Horsepower}\label{horsepower-1}

Horsepower rating of the pump or pump motor (if tracked as a separate
attribute in implementation).

\subsection{FlowRate}\label{flowrate-1}

Flow rate of the pump in gallons per minute (GPM) for individual pump
units, if applicable.

\subsection{PumpStationType}\label{pumpstationtype-1}

Type of pump station installation.

Value

Description

Permanent Pump Station

Fixed, permanent pump station installation.

Portable Pump

Portable or temporary pump installation.

\subsection{sdsMetaDataID}\label{sdsmetadataid-10}

References attached metadata.

\subsection{RestorFeatIDPK}\label{restorfeatidpk-10}

Primary key used to uniquely identify the restoration feature record.

\subsection{sdsID}\label{sdsid-10}

Identifier used to link the feature to supporting SDS (spatial data
system) records or related datasets.

\subsection{mediaIDFK}\label{mediaidfk-10}

Foreign key used to reference associated media (e.g., photographs,
design drawings, pump curves, as-built plans) documenting the pump
station.

\chapter{Revetment}\label{revetment}

\section{Spatial \& Geometric
Definition}\label{spatial-geometric-definition-11}

\begin{itemize}
\item
  Geometry Type: Polygon or Polyline(Revetment\_Armor\_Band).
\item
  Spatial Delineation: Armored protective blanket placed along littoral
  shoreline banks, island provar edges, wing dam abutments, and channel
  margins.
\item
  Materials: Heavy quarried angular stone riprap, crushed aggregate
  bedding, and non-woven geotextile filter fabric.
\end{itemize}

\subsection{Primary Restoration \& Engineering
Objectives}\label{primary-restoration-engineering-objectives-10}

\begin{itemize}
\item
  Shoreline Erosion Arrest: Neutralize high hydraulic shear stresses
  (τ=γRS) induced by commercial towboat propeller wash, vessel wakes,
  and flood currents.
\item
  Wind-Wave Dissipation: Protect constructed earthen island fills and
  natural delta edges from wind-fetch wave impact and slumping.
\item
  Littoral Micro-Cover: Stone voids provide interstitial refuge and
  feeding zones for macroinvertebrates and small fish species.
\end{itemize}

\subsection{Typical Cross-Sections, Dimensions \& Sizing
Criteria}\label{typical-cross-sections-dimensions-sizing-criteria-8}

\begin{itemize}
\item
  Armor Sizing: USACE Class II to Class IV stone riprap (d 50 =12 to 18
  in; top weight up to 1,500 lbs).
\item
  Blanket Thickness: 1.5× to 2.0× the d 100 stone size (typically 18 to
  30 in thick).
\item
  Slope Geometry: Maximum stable slope angle is 1V:2H (1V:3H preferred).
\item
  Toe Key-In: Armored toe trench excavated 3.0 to 5.0 ft into the
  channel bed, or a wide stone launching apron (5 to 10 ft wide)
  designed to self-launch and protect the bank toe against scour.
\end{itemize}

\subsection{Construction \& Implementation
Methodology}\label{construction-implementation-methodology-11}

\begin{itemize}
\item
  Marine Placement: Work conducted via spud barges using long-reach
  excavators equipped with hydraulic stone grapples or clamshell buckets
  to place stone smoothly without tearing geotextile.
\item
  Filter Layer Installation: High-survivability non-woven geotextile
  placed under the stone to stop fine-grained subgrade soils from
  leaching out through riprap voids.
\end{itemize}

\subsection{Quality Control, Inspection \& Long-Term
Maintenance}\label{quality-control-inspection-long-term-maintenance-11}

\begin{itemize}
\item
  Gradation Testing: Quarry riprap testing per ASTM standards to ensure
  durability, bulk specific gravity (≥2.60), and resistance to
  freeze-thaw cycles.
\item
  Post-Flood Audits: Boat-based visual and sonar scans to detect riprap
  displacement, toe undermining, or ice plucking along the waterline.
\end{itemize}

\section{Spatial Feature Attribute}\label{spatial-feature-attribute-8}

\subsection{Definition}\label{definition-10}

A revetment is a cover or facing of erosion-resistant material placed
directly on an existing slope, embankment, or dike. Unlike seawalls,
revetments are not self-supporting and depend on the underlying soil for
support.

\subsection{Purpose}\label{purpose-13}

Revetments are used to stabilize river banks and protect shorelines from
strong waves and currents. Bank revetments prevent erosion, maintain a
stable channel, and support new habitat creation. (UMRR Design Handbook,
2012)

\subsection{Applicability}\label{applicability-11}

The primary forces that affect shorelines are river currents and
wind-driven wave action; ice action and waves created by towboats or
recreational boats can also cause erosion. Revetment stabilization
techniques can be employed singly or in combination to protect shoreline
and add habitat diversity. Gradual side slopes and sand or mud soils can
benefit turtles and waterbirds that nest, feed, and loaf on shorelines.
Native plantings can be more aesthetically pleasing than traditional
riprap, and traditional stabilization techniques can be refined (e.g.,
larger and mixed-grade rock) to improve habitat diversity for fish and
invertebrates. (Report to Congress, 2004)

\subsection{Notes}\label{notes-13}

Restoration features now have a Program\_UID field to distinguish
between NESP and UMRR features. Aliases without program identification
will be created for any field that currently has UMRR in the name, so
they can be used for NESP or UMRR data.

Attribute criteria for each field includes the following:

\begin{itemize}
\tightlist
\item
  Attribute Status -- defines whether the field requires a value or if
  it is optional based on available information.
\item
  Phases -- indicates the earliest time in the project lifecycle values
  for this attribute are available.
\item
  Digitizing requirements -- indicates when the field value should be
  entered in the attribute table.
\end{itemize}

Geometry Note: Points should have well top elevation as Z value if
available.

\subsection{Revetment}\label{revetment-1}

(polygon) -- An area faced with stone or other material along the
shoreline designed to stabilize the bank or protect it from erosion.

\section{\texorpdfstring{\protect\includegraphics[width=0.75\linewidth,height=\textheight,keepaspectratio]{images/revetment.jpg}}{}}\label{section-2}

\section{Feature Attribute
Criteria}\label{feature-attribute-criteria-11}

\#

Atribute

Attribute Status

Phase

Digitizing Requirement

1

Program\_UID

Required

Feasibility

At feature creation

2

PROJECT\_ID

Required

Feasibility

At feature creation

3

Type

Required

Feasibility

At feature creation

4

Subtype

Required

Feasibility

At feature creation

5

RestorFeatStatus

Required

Design

At feature creation

6

sdsFeatureName

Optional

Design

As required and available

7

sdsFeatureDescription

Optional

Design

As required and available

8

RestorFeatVdatum

Optional

Design

As required and available

9

QuantityTONS

Optional

Design

As required and available

10

MaterialThicknessIN

Optional

Construction

As required and available

11

BeddingThicknessIN

Optional

Construction

As required and available

12

GeoTextile

Optional

Construction

As required and available

13

StartStation

Optional

Construction

As required and available

14

EndStation

Optional

Construction

As required and available

15

sdsMetaDataID

Optional

Construction

As required and available

16

RestorFeatIDPK

Optional

Construction

As required and available

17

sdsID

Required

Construction

As required and available

18

mediaIDFK

Required

Construction

As required and available

\begin{center}\rule{0.5\linewidth}{0.5pt}\end{center}

\section{Parameters}\label{parameters-11}

\subsection{Program\_UID}\label{program_uid-11}

Designate program NESP or UMRR.

\subsection{Project\_ID (alias:
PROJECT\_ID)}\label{project_id-alias-project_id-11}

Foreign key to the unique NESP/UMRR project ID.

\subsection{Type (alias)}\label{type-alias-10}

A broad classification used to group NESP/UMRR restoration features into
categories of similar function.

Value

Description

Shoreline Protection

Measures taken to reduce erosion from banks and islands due to wind,
wave, and ice action.

Channel Modifications

Measures taken to train the main channel of a river and to modify
secondary channels.

Water Level Management

Measures taken to imitate natural low-flow river conditions with the
goal of improving wetland habitats.

\subsection{SubType (alias: Subtype)}\label{subtype-alias-subtype-11}

A classification used to group NESP/UMRR restoration features into
specific revetment design feature categories.

Value

Description

Rock revetment

Placement of riprap along a shoreline to a specified depth; may include
a geotextile filter beneath rock. Generally used to reduce shoreline
erosion from waves and ice.

Off-bankline revetment

Rock structure placed parallel to a caving river bank, which can be
notched; used to reduce erosion and provide diverse habitat.

Root-wad revetment

Revetment constructed by keying logs into the bankline with root wads
extending into the channel; used to reduce stream bank erosion, provide
fish cover, shade, detritus, insect habitat, and a diversity of
habitats.

Concrete Revetment

Revetment constructed from concrete.

Sheet pile revetment

Revetment constructed using interlocking sheet piles to form abutments
or retaining walls.

Gabion Wall

Wall constructed from large stones or concrete blocks placed in a wire
mesh cage for protection against waves or water turbulence.

\subsection{RestorFeatStatus}\label{restorfeatstatus-11}

A description of the status of the individual restoration, including the
party responsible for the feature.

Value

Description

USACE Planned

Feature is in a USACE planning stage.

USACE Built

Feature has been built by USACE. Once constructed by USACE, it remains
classified as USACE Built even if later repaired.

USACE Enhanced

Function of an existing feature not built by USACE has been enhanced by
a USACE-built or repaired feature.

USACE Repaired

Existing feature not built by USACE has been directly improved or
repaired by USACE.

Other Planned

Feature planned by another organization or government agency.

Other Built

Feature built by another organization or government agency and integral
to the UMRR project.

Other Enhanced

Function of a feature built by another organization has been enhanced by
a non-USACE organization or government agency.

Other Repaired

Existing feature directly improved or repaired by another organization
or government agency.

\subsection{sdsFeatureName}\label{sdsfeaturename-11}

The common name of the feature.

\subsection{sdsFeatureDescription}\label{sdsfeaturedescription-11}

A narrative describing the feature.

\subsection{RestorFeatVdatum}\label{restorfeatvdatum-11}

Vertical datum of elevations recorded in the attribute table or used as
feature z values.

Value

Description

MSL 1912

Mean Sea Level of 1912.

NGVD 1929

National Geodetic Vertical Datum of 1929.

NAVD 1988

North American Vertical Datum of 1988.

Other

Other; note datum in sdsFeatureDescription.

Unknown

Unknown vertical datum.

\subsection{QuantityTONS}\label{quantitytons}

Quantity of revetment material placed, measured in tons.

\subsection{MaterialThicknessIN}\label{materialthicknessin}

Thickness of revetment material placed, measured in inches.

\subsection{BeddingThicknessIN}\label{beddingthicknessin}

Thickness of bedding or underlying layer (e.g., bedding stone or sand)
in inches.

\subsection{GeoTextile}\label{geotextile}

Indicates if geotextile fabric was used beneath or within the revetment.

Value

Description

Yes

Geotextile fabric was used.

No

Geotextile fabric was not used.

Unknown

Use of geotextile fabric is unknown.

\subsection{StartStation}\label{startstation-1}

Beginning station of the revetment feature, typically referenced to a
project baseline or alignment (e.g., 10+00).

\subsection{EndStation}\label{endstation-1}

Ending station of the revetment feature.

\subsection{sdsMetaDataID}\label{sdsmetadataid-11}

References attached metadata.

\subsection{RestorFeatIDPK}\label{restorfeatidpk-11}

Primary key used to uniquely identify the restoration feature record.

\subsection{sdsID}\label{sdsid-11}

Identifier used to link the feature to supporting SDS (spatial data
system) records or related datasets.

\subsection{mediaIDFK}\label{mediaidfk-11}

Foreign key used to reference associated media (e.g., photographs,
design drawings, as-built plans) documenting the revetment.

\chapter{River Flow Structure}\label{river-flow-structure}

\section{Spatial \& Geometric
Definition}\label{spatial-geometric-definition-12}

\begin{itemize}
\item
  Geometry Type: Polygon (Flow\_Structure\_Polygon).
\item
  Spatial Delineation: 3D structural stone footprint resting on the
  riverbed, representing river training dikes, chevrons, bendway weirs,
  closing dams, and hard-points.
\item
  Structural Material: 100\% quarried limestone/dolomite riprap and
  stone core.
\end{itemize}

\subsection{Primary Restoration \& Engineering
Objectives}\label{primary-restoration-engineering-objectives-11}

\begin{itemize}
\item
  Channel Steering \& Velocity Partitioning: Direct river discharge into
  primary navigation lines while partitioning secondary flows into
  restoration backwaters.
\item
  Scour Hole Generation: Carve deep structural aquatic habitat niches
  through localized flow acceleration around structure tips and heads.
\item
  Island \& Shoreline Sheltering: Act as rigid upstream hydrodynamic
  deflection shields protecting soft downstream wetland deposits.
\end{itemize}

\subsection{Typical Cross-Sections, Dimensions \& Sizing
Criteria}\label{typical-cross-sections-dimensions-sizing-criteria-9}

\begin{itemize}
\item
  Crest Widths: 6.0 to 12.0 ft along the top weir crest.
\item
  Side Slopes: Heavy stable stone slopes of 1V:1.5H to 1V:2H.
\item
  Crest Elevation: Range from submerged low sills (1.0 to 3.0 ft below
  pool) up to emergent structures (2.0 to 5.0 ft above normal pool).
\item
  Structural Prow / Nose: Heavily reinforced rounded apex with stone
  sizing upgraded to Class V (d 50 ≥24 in) to withstand direct river ice
  sheets.
\end{itemize}

\subsection{Construction \& Implementation
Methodology}\label{construction-implementation-methodology-12}

\begin{itemize}
\item
  Marine Derrick Operations: Heavy rock barges positioned with tugboats;
  stone placed using clamshell buckets or hydraulic excavators working
  against river currents.
\item
  Upstream Root Tie-In: Structure landward roots deeply keyed into
  adjacent banks (15 to 30 ft) to prevent high-water flanking.
\end{itemize}

\subsection{Quality Control, Inspection \& Long-Term
Maintenance}\label{quality-control-inspection-long-term-maintenance-12}

\begin{itemize}
\tightlist
\item
  Acoustic Bathymetric Scans: Side-scan and multi-beam hydrographic
  surveys monitoring bed scour development and verifying structural mass
  stability.
\end{itemize}

\section{Spatial Feature Attributes}\label{spatial-feature-attributes-2}

\subsection{Definition}\label{definition-11}

River Flow Structure polygon is a natural or man-made flow (or sediment)
control structure in a water course or water body such as a dike or
weir. This feature should not be used to model a levee. The polygon
feature class is for features large enough that a polygon can be drawn
around them and reasonably seen at a project-level extent.

\subsection{Purpose}\label{purpose-14}

Riverine flow structures, also called ``river training structures,'' are
man-made constructions designed to modify the flow of a river, typically
to control erosion, manage sediment transport, improve navigation,
protect riverbanks, and sometimes enhance aquatic habitats by
influencing water flow dynamics within the river channel.

\subsection{Applicability}\label{applicability-12}

Training structures can be used to alter hydrodynamic conditions, the
sediment transport regime, and ultimately habitat conditions on the
UMRS. The 2010 EMP Report to Congress described how river training
structures meet UMRS Environmental Objectives by improving fish habitat
and water quality, stabilizing eroding channels, reducing sediment load
to backwaters by reducing flow velocities, and maintaining water
temperature and rock substrate. (UMRR Design Handbook, 2012)

\subsection{Notes}\label{notes-14}

Restoration features now have a Program\_UID field to distinguish
between NESP and UMRR features. Aliases without program identification
will be created for any field that currently has UMRR in the name, so
they can be used for NESP or UMRR data.

Attribute criteria for each field includes the following:

\begin{itemize}
\tightlist
\item
  Attribute Status -- defines whether the field requires a value or if
  it is optional based on available information.
\item
  Phases -- indicates the earliest time in the project lifecycle values
  for this attribute are available.
\item
  Digitizing requirements -- indicates when the field value should be
  entered in the attribute table.
\end{itemize}

Geometry Note: Points should have well top elevation as Z value if
available.

\subsection{RiverineFlowStructure}\label{riverineflowstructure}

(polygon) -- A natural or man-made flow (or sediment) control structure
in a water course or water body such as a dike or weir. This feature
should be used to represent the footprint of the feature.

\begin{center}\rule{0.5\linewidth}{0.5pt}\end{center}

\section{Feature Attribute
Criteria}\label{feature-attribute-criteria-12}

\#

Attribute

Attribute Status

Phase

Digitizing Requirements

1

Program\_UID

Required

Feasibility

At feature creation

2

PROJECT\_ID (Alias)

Required

Feasibility

At feature creation

3

Type (Alias)

Required

Feasibility

At feature creation

4

Subtype (Alias)

Required

Feasibility

At feature creation

5

RestorFeatStatus

Required

Feasibility

At feature creation

6

sdsFeatureName

Optional

Design

As required and available

7

sdsFeatureDescription

Optional

Design

As required and available

8

RestorFeatVdatum

Optional

Design

As required and available

9

WaterOvertopping

Optional

Design

As required and available

10

PrimaryMaterial

Optional

Design

As required and available

11

SecondaryMaterial

Optional

Design

As required and available

12

DredgeMaterialPlacement

Optional

Design

As required and available

13

Acres

Optional

Design

As required and available

14

GeoTextileMat

Optional

Construction

As required and available

15

sdsMetaDataID

Optional

Construction

As required and available

16

RestorFeatIDPK

Optional

Construction

As required and available

17

sdsID

Optional

Construction

As required and available

18

mediaIDFK

Optional

Construction

As required and available

\begin{center}\rule{0.5\linewidth}{0.5pt}\end{center}

\section{Parameters}\label{parameters-12}

\subsection{Program\_UID}\label{program_uid-12}

Designate program NESP or UMRR.

\subsection{Project\_ID (alias:
PROJECT\_ID)}\label{project_id-alias-project_id-12}

Foreign key to the unique NESP/UMRR project ID.

\subsection{Type (alias)}\label{type-alias-11}

A broad classification used to group NESP/UMRR restoration features into
categories of similar function.

Value

Description

Shoreline Protection

Measures taken to reduce erosion from banks and islands due to wind,
wave, and ice action.

Channel Modifications

Measures taken to train the main channel of a river and to modify
secondary channels.

\subsection{SubType (alias: Subtype)}\label{subtype-alias-subtype-12}

A classification used to group NESP/UMRR restoration features into
specific design feature categories.

Value

Description

Rock groin

Placement of rock extending from shore into shallow water, perpendicular
to shore. Encourages deposition of sediment along the groin and erosion
between groins to form stable scallop shapes; appropriate where wave
action and littoral drift are dominant processes.

Rock vane

Placement of rock extending from shore into water at 45--60 degrees
upstream, away from shore; features a sloping top and directs current
away from shore; can also reduce littoral drift.

J-hook vane

Rock vane with a ``J''-shaped hook whose outside curve faces upstream.

Cross vane

U-shaped vane keyed into both banklines and sloping down to the top of
the U pointing upstream; used to offset adverse effects of straight
weirs and check dams.

Root-wad vane

Placement of root wads extending from shore into water at 45--60 degrees
upstream, away from shore.

Spur dike

Series of river training structures perpendicular to shore used to
establish normal channel width, direct flow, promote scour and
deposition, and trap bedload to build banks.

L-head dike

Training dike with a perpendicular dike structure attached at the
channel end (shoreward), forming an ``L'' shape to control scour
patterns at the riverward end.

Alternating dike

Structures constructed in straight side channels to encourage
development of sinuous flow.

Stepped-up dike

Field of dikes that increase in elevation (generally \textasciitilde2 ft
dike-to-dike) to change morphology of natural alluvial waterways.

Chevron dike

V- or U-shaped rock dike pointing upstream and unattached to bankline,
typically built above flat-pool elevation to \textasciitilde2-year flood
elevation; used in wide reaches to split flow and enhance habitat and to
provide dredged-material placement area.

Longitudinal peak stone toe

Continuous stone dike keyed into bank at upstream and downstream ends,
placed at the toe of an eroding bend bankline; can be offset from bank
to smooth alignment through a bend.

Bendway weir

Low-level, completely submerged rock structure placed on the outside
bankline of a bend and angled upstream toward flow, extending into the
navigation channel; prevents excess deepening and reduces outside-bank
erosion.

W-weir

W-shaped rock structure in the channel, center of W facing downstream;
ends tied into bank at bankfull elevation; slopes down to two upstream
points and up slightly to center; provides bed control and enhanced fish
habitat.

Closure structure

Rock, earth, or wood structures constructed across secondary channels to
reduce floodplain conveyance and increase main channel depths; may be
submerged or emerged.

Rock-Log structure

Economical alternative to offshore rock mounds where ice impacts are
minimal; logs anchored at both ends by rock.

Wood pile structure

Logs driven into the river bed to create roughness and form river
training structures, benefiting macro- and microinvertebrates.

Offshore rock mounds

Rock placed offshore from shorelines where shoreline protection is
needed, particularly where nearshore bathymetry is shallow, shorelines
are poorly defined, or direct rock placement is difficult.

Boulder clusters

Large stones placed slightly below base-flow water surface elevation to
diversify flow regime and provide in-channel refuges.

Multiple round points

Field of round stone structures, generally slightly emerged, used to
create bathymetric and flow diversity.

Side channel hard point

Stone structures placed at regular intervals along an eroding bank,
trenched, keyed, or dumped on shore to resist forces associated with
bank failure.

Woody debris

Trees at least 25 ft long placed perpendicular or at an angle to
shoreline; trunks may be buried or rock-keyed, typically in
wave-sheltered areas.

Spillway

Section of embankment engineered to provide controlled overtopping of an
embankment system, typically at the downstream end at a lower elevation
than the exterior berm.

Diversion

Structure intended to divert flow from the main channel, typically into
a bypass or cut-off channel.

Other

Other type of riverine flow structure; must be specified in
sdsFeatureDescription.

\subsection{RestorFeatStatus}\label{restorfeatstatus-12}

A description of the status of the individual restoration, including the
responsible party.

Value

Description

USACE Planned

Feature is in a USACE planning stage.

USACE Built

Feature has been built by USACE. Once constructed by USACE, it remains
classified as USACE Built even if later repaired.

USACE Enhanced

Function of a non-USACE feature enhanced by a USACE-built or repaired
feature.

USACE Repaired

Non-USACE feature directly improved or repaired by USACE.

Other Planned

Feature planned by another organization or government agency.

Other Built

Feature built by another organization or government agency and integral
to the UMRR project.

Other Enhanced

Function of a feature built by another organization enhanced by a
non-USACE organization or government agency.

Other Repaired

Existing feature directly improved or repaired by another organization
or government agency.

\subsection{sdsFeatureName}\label{sdsfeaturename-12}

The common name of the feature.

\subsection{sdsFeatureDescription}\label{sdsfeaturedescription-12}

A narrative describing the feature.

\subsection{RestorFeatVdatum}\label{restorfeatvdatum-12}

Vertical datum of elevations recorded in the attribute table or used as
feature z values.

Value

Description

MSL 1912

Mean Sea Level of 1912.

NGVD 1929

National Geodetic Vertical Datum of 1929.

NAVD 1988

North American Vertical Datum of 1988.

Other

Other; note datum in sdsFeatureDescription.

Unknown

Unknown vertical datum.

\subsection{WaterOvertopping}\label{waterovertopping-1}

Description of the relationship between the feature and water level.

Value

Description

AlwaysDry

Not covered at high water under average meteorological conditions.

AlwaysSubmerged

Remains covered by water at all times under average meteorological
conditions.

Awash

Flush with or washed by waves at low water under average meteorological
conditions.

CoversAndUncovers

Projection that periodically extends above and is submerged below the
water surface.

Floating

Resting or moving on the surface of a water body.

PartlySubmergedAtHighWater

Partially covered and partially dry at high water.

SubjectToInundationOrFlooding

Area or structure periodically covered by flood waters.

Unknown

Relationship to water level is unknown.

NA

Not applicable.

\subsection{PrimaryMaterial}\label{primarymaterial-1}

Primary construction material of the structure.

Value

Description

Sand

Sand.

Rock

Rock.

Fines

Fine sediment including clay and clay mixtures.

Timber/Wood

Timber or wood, including logs or root wads.

Earthen

Mixture of sand and fines.

Concrete

Concrete.

\subsection{SecondaryMaterial}\label{secondarymaterial-1}

Secondary construction material, if applicable.

Value

Description

Sand

Sand.

Rock

Rock.

Fines

Fine sediment including clay and clay mixtures.

Timber/Wood

Timber or wood, including logs or root wads.

Earthen

Mixture of sand and fines.

Concrete

Concrete.

\subsection{DredgeMaterialPlacement}\label{dredgematerialplacement-1}

Describes whether dredged material was placed to create the feature.

Value

Description

Yes

Dredged material was placed to create or augment the feature.

No

No dredged material was used.

NA

Not applicable.

Unknown

Use of dredged material is unknown.

\subsection{Acres}\label{acres-2}

Area of the feature in acres.

\subsection{GeoTextileMat}\label{geotextilemat}

Indicates whether geotextile mat or similar fabric was used beneath or
within the structure.

Value

Description

Yes

Geotextile mat was used.

No

Geotextile mat was not used.

Unknown

Use of geotextile mat is unknown.

\subsection{sdsMetaDataID}\label{sdsmetadataid-12}

References attached metadata.

\subsection{RestorFeatIDPK}\label{restorfeatidpk-12}

Primary key used to uniquely identify the restoration feature record.

\subsection{sdsID}\label{sdsid-12}

Identifier used to link the feature to supporting SDS (spatial data
system) records or related datasets.

\subsection{mediaIDFK}\label{mediaidfk-12}

Foreign key used to reference associated media (e.g., photographs,
design drawings, as-built plans) documenting the riverine flow
structure.

\chapter{River Flow Structure
Centerline}\label{river-flow-structure-centerline}

\section{Spatial \& Geometric
Definition}\label{spatial-geometric-definition-13}

\begin{itemize}
\item
  Geometry Type: Polyline (Flow\_Structure\_Centerline).
\item
  Spatial Delineation: Longitudinal linear crest axis of a riverine
  training dike, chevron, or rock closure structure.
\item
  Data Fields: Structure ID, structure class (wing dike, closing dam,
  chevron), crest elevation, top width, upstream azimuth, and baseline
  stationing.
\end{itemize}

\subsection{Primary Restoration \& Engineering
Objectives}\label{primary-restoration-engineering-objectives-12}

\begin{itemize}
\item
  Alignment Guidance: Guides derrick barge positioning during underwater
  rock dumping operations.
\item
  Hydraulic Modeling Reference: Defines the exact mathematical weir
  boundary line in 2D and 3D hydrodynamic models (e.g., Adaptive
  Hydraulics - AdH, HEC-RAS 2D).
\end{itemize}

\subsection{Typical Specifications \& Geometric
Properties}\label{typical-specifications-geometric-properties-1}

\begin{itemize}
\item
  Survey Baseline: Linear axis delineating the centerline of stone
  placement; elevation profile reflects the exact design crest grade
  without sagging.
\item
  Coordinate Accuracy: Sub-decimeter horizontal precision tied into
  local State Plane coordinates.
\end{itemize}

\subsection{Construction \& Implementation
Methodology}\label{construction-implementation-methodology-13}

\begin{itemize}
\tightlist
\item
  DGPS Marine Guidance Systems: Centerline profile strings loaded into
  onboard barge positioning systems (e.g., Trimble Marine Construction
  software) to ensure stone dumping remains centered within design
  tolerances (±1.0 ft).
\end{itemize}

\subsection{Quality Control, Inspection \& Long-Term
Maintenance}\label{quality-control-inspection-long-term-maintenance-13}

\begin{itemize}
\tightlist
\item
  Crest Alignment Verification: High-precision RTK or lidar scans along
  the exposed crest during low-flow river stages to confirm structure
  crest stability.
\end{itemize}

\section{Spatial Feature Attribute}\label{spatial-feature-attribute-9}

\subsection{Definition}\label{definition-12}

Riverine Flow Structure Line is a natural or man-made flow (or sediment)
control structure in a water course or water body such as a dike or
weir. This feature should not be used to model a levee. The line feature
class is for features that are too small to be drawn as polygons and are
better represented as lines.

\subsection{Purpose}\label{purpose-15}

Riverine flow structures, also called ``river training structures,'' are
man-made constructions designed to modify the flow of a river, typically
to control erosion, manage sediment transport, improve navigation,
protect riverbanks, and sometimes enhance aquatic habitats by
influencing water flow dynamics within the river channel.

\subsection{Applicability}\label{applicability-13}

Refined construction techniques improve habitat through site-specific
modifications. Successfully combining a broad range of construction
techniques with approaches that strive to use or mimic natural riverine
processes provides benefits to the river at system, reach, pool, and
local scales. Programs continually build upon lessons learned in
constructing and managing prior projects. (UMRR Design Handbook, 2012)

\subsection{Notes}\label{notes-15}

Restoration features now have a Program\_UID field to distinguish
between NESP and UMRR features. Aliases without program identification
will be created for any field that currently has UMRR in the name, so
they can be used for NESP or UMRR data.

Attribute criteria for each field includes the following:

\begin{itemize}
\tightlist
\item
  Attribute Status -- defines whether the field requires a value or if
  it is optional based on available information.
\item
  Phases -- indicates the earliest time in the project lifecycle values
  for this attribute are available.
\item
  Digitizing requirements -- indicates when the field value should be
  entered in the attribute table.
\end{itemize}

Geometry Note: Points should have well top elevation as Z value if
available.

\subsection{RiverineFlowStructureLine}\label{riverineflowstructureline}

(line) -- A natural or man-made flow (or sediment) control structure in
a water course or water body such as a dike or weir. This feature should
be used to represent the footprint of the feature.

\begin{center}\rule{0.5\linewidth}{0.5pt}\end{center}

\section{Feature Attribute
Criteria}\label{feature-attribute-criteria-13}

\#

Attribute

Attribute Status

Phase

Digitizing Requirements

1

Program\_UID

Required

Feasibility

At feature creation

2

PROJECT\_ID (Alias)

Required

Feasibility

At feature creation

3

Type (Alias)

Required

Feasibility

At feature creation

4

Subtype (Alias)

Required

Feasibility

At feature creation

5

RestorFeatStatus

Required

Feasibility

At feature creation

6

sdsFeatureName

Optional

Design

As required and available

7

sdsFeatureDescription

Optional

Design

As required and available

8

RestorFeatVdatum

Optional

Design

As required and available

9

WaterOvertopping

Optional

Design

As required and available

10

PrimaryMaterial

Optional

Design

As required and available

11

SecondaryMaterial

Optional

Design

As required and available

12

DredgeMaterialPlacement

Optional

Design

As required and available

13

Acres

Optional

Design

As required and available

14

GeoTextileMat

Optional

Construction

As required and available

15

sdsMetaDataID

Optional

Construction

As required and available

16

RestorFeatIDPK

Optional

Construction

As required and available

17

sdsID

Optional

Construction

As required and available

18

mediaIDFK

Optional

Construction

As required and available

\begin{center}\rule{0.5\linewidth}{0.5pt}\end{center}

\section{Parameters}\label{parameters-13}

\subsection{Program\_UID}\label{program_uid-13}

Designate program NESP or UMRR.

\subsection{Project\_ID (alias:
PROJECT\_ID)}\label{project_id-alias-project_id-13}

Foreign key to the unique NESP/UMRR project ID.

\subsection{Type (alias)}\label{type-alias-12}

A broad classification used to group NESP/UMRR restoration features into
categories of similar function.

Value

Description

Shoreline Protection

Measures taken to reduce erosion from banks and islands due to wind,
wave, and ice action.

Channel Modifications

Measures taken to train the main channel of a river and to modify
secondary channels.

\subsection{SubType (alias: Subtype)}\label{subtype-alias-subtype-13}

A classification used to group NESP/UMRR restoration features into
specific design feature categories.

Value

Description

Rock groin

Placement of rock extending from shore into shallow water, perpendicular
to shore. Encourages deposition of sediment along the groin and erosion
in the reach between two groins to form stable scallop shapes.
Appropriate where wave action and littoral drift are dominant processes.

Rock vane

Placement of rock extending from shore into water at angles from 45--60
degrees upstream, away from shore. Features a sloping top and functions
to direct current away from shore; can also reduce littoral drift.

J-hook vane

Placement of rock extending from shore into water at angles from 45--60
degrees upstream, away from shore, with a ``J''-shaped hook whose
outside curve faces upstream.

Cross vane

U-shaped vane with either end keyed into banklines and sloping down to
the top of the U pointing upstream, used to offset adverse effects of
straight weirs and check dams.

Root-wad vane

Placement of root wads extending from shore into water at angles from
45--60 degrees upstream, away from shore.

Spur dike

River training structures generally constructed in a series
perpendicular to shore to establish normal channel width, direct the
axis of flow, promote scour and sediment deposition, and trap bedload to
build banks.

L-head dike

Training dike with a perpendicular dike structure attached at the
channel end, creating an ``L'' shape; controls scour patterns at the
riverward end of the training dike.

Alternating dike

River training structures usually constructed in straight side channels
to encourage development of sinuous flow.

Stepped-up dike

Field of dikes that increase in elevation (generally \textasciitilde2 ft
from dike to dike) used to change morphology of natural alluvial
waterways.

Chevron dike

V- or U-shaped rock dike pointing upstream and unattached to the
bankline, typically built above normal flat pool elevation to
approximately a 2-year flood elevation; used in wider reaches to provide
flow split and habitat and to provide dredged-material placement area.

Longitudinal peak stone toe

Continuous stone dike keyed into the bank at both upstream and
downstream ends, placed at the toe of an eroding bend bankline; may be
placed slightly away from the bank to smooth alignment through the bend.

Bendway weir

Low-level, completely submerged rock structure placed on the outside
bankline of a river bend and angled upstream toward flow, extending into
the navigation channel, used to prevent excess channel deepening and
reduce erosion on the outside of the bendway.

W-weir

Placement of rock in a W shape in the channel, center point facing
downstream; each end is tied into the bank at bankfull elevation and
slopes down to the two upstream points and up slightly to the center
point, providing bed control and enhanced fish habitat.

Closure structure

Rock, earth, or wood structures constructed across secondary channels to
reduce floodplain conveyance and increase main channel depths; can be
submerged or emerged.

Rock-Log structure

Economical alternative to offshore rock mounds where ice impacts are
minimal; piles of rock anchor both ends of logs.

Wood pile structure

Logs driven into the river bed to create roughness and form river
training structures, providing benefit to macro- and microinvertebrate
species.

Offshore rock mounds

Placement of rock offshore from natural shorelines where shoreline
protection is needed, typically where nearshore bathymetry is shallow or
shorelines are poorly defined, or where direct placement of rock is
difficult.

Boulder clusters

Large stone placed slightly below base-flow water surface elevation to
diversify flow regime and provide in-channel refuges for aquatic
species.

Multiple round points

Field of round stone structures, generally slightly emerged, used to
create bathymetric and flow diversity.

Side channel hard point

Stone structures placed at regular intervals along an eroding bank,
which may be trenched in, keyed in, or dumped on shore to resist forces
associated with bank failure.

Woody debris

Placement of trees at least 25 ft long perpendicular or at an angle to
the shoreline; trunks can be buried in shore sand or anchored using a
rock key-in, usually in areas sheltered from wind-generated waves.

Spillway

Section of embankment engineered to provide controlled overtopping of an
embankment system, typically at the downstream end of a site at an
elevation lower than the exterior berm.

Diversion

Structure intended to divert flow from the main channel, typically into
a bypass or cut-off channel.

Other

Other type of riverine flow structure; must be specified in the
sdsFeatureDescription field.

\subsection{RestorFeatStatus}\label{restorfeatstatus-13}

A description of the status of the individual restoration. It also
includes the party responsible for the individual feature.

Value

Description

USACE Planned

Feature is in a USACE planning stage.

USACE Built

Feature has been built by USACE. If a feature has been planned and built
by USACE, once constructed, it will always be considered USACE built
even if repaired later.

USACE Enhanced

The function of an existing feature not built by USACE has been enhanced
by a feature that was USACE built or repaired.

USACE Repaired

An existing feature not built by USACE has been directly improved or
repaired.

Other Planned

Feature planned by another organization or government agency.

Other Built

Feature built by another organization or government agency and is an
integral part of the UMRR project.

Other Enhanced

Function of an existing feature built by another organization or
government agency has been enhanced by a feature built or repaired by a
non-USACE organization or government agency.

Other Repaired

Existing feature directly improved or repaired by another organization
or government agency.

\subsection{sdsFeatureName}\label{sdsfeaturename-13}

The common name of the feature.

\subsection{sdsFeatureDescription}\label{sdsfeaturedescription-13}

A narrative describing the feature.

\subsection{RestorFeatVdatum}\label{restorfeatvdatum-13}

Vertical datum of elevations recorded in the attribute table or used as
feature z values.

Value

Description

MSL 1912

Mean Sea Level of 1912.

NGVD 1929

National Geodetic Vertical Datum of 1929.

NAVD 1988

North American Vertical Datum of 1988.

Other

Other; note datum in sdsFeatureDescription.

Unknown

Unknown vertical datum.

\subsection{WaterOvertopping}\label{waterovertopping-2}

Description of the relationship between the feature and water level.

Value

Description

AlwaysDry

Not covered at high water under average meteorological conditions.

AlwaysSubmerged

Remains covered by water at all times under average meteorological
conditions.

Awash

Flush with or washed by waves at low water under average meteorological
conditions.

CoversAndUncovers

Projection from the bottom that periodically extends above and is
submerged below the water surface.

Floating

Resting or moving on the surface of a water body.

PartlySubmergedAtHighWater

Partially covered and partially dry at high water.

SubjectToInundationOrFlooding

Area or structure periodically covered by flood waters.

Unknown

Relationship to water level is unknown.

NA

Not applicable.

\subsection{PrimaryMaterial}\label{primarymaterial-2}

Primary construction material of the structure.

Value

Description

Sand

Sand.

Rock

Rock.

Fines

Fine sediment including clay and clay mixtures.

Timber/Wood

Timber or wood, including logs or root wads.

Earthen

Mixture of sand and fines.

Concrete

Concrete.

\subsection{SecondaryMaterial}\label{secondarymaterial-2}

Secondary construction material, if applicable.

Value

Description

Sand

Sand.

Rock

Rock.

Fines

Fine sediment including clay and clay mixtures.

Timber/Wood

Timber or wood, including logs or root wads.

Earthen

Mixture of sand and fines.

Concrete

Concrete.

\subsection{DredgeMaterialPlacement}\label{dredgematerialplacement-2}

Describes whether dredged material was placed to create the feature.

Value

Description

Yes

Dredged material was placed to create or augment the feature.

No

No dredged material was used.

NA

Not applicable.

Unknown

Use of dredged material is unknown.

\subsection{Acres}\label{acres-3}

Area associated with the feature in acres (for reference where area is
calculated or tracked).

\subsection{GeoTextileMat}\label{geotextilemat-1}

Indicates whether geotextile mat or similar fabric was used beneath or
within the structure.

Value

Description

Yes

Geotextile mat was used.

No

Geotextile mat was not used.

Unknown

Use of geotextile mat is unknown.

\subsection{sdsMetaDataID}\label{sdsmetadataid-13}

References attached metadata.

\subsection{RestorFeatIDPK}\label{restorfeatidpk-13}

Primary key used to uniquely identify the restoration feature record.

\subsection{sdsID}\label{sdsid-13}

Identifier used to link the feature to supporting SDS (spatial data
system) records or related datasets.

\subsection{mediaIDFK}\label{mediaidfk-13}

Foreign key used to reference associated media (e.g., photographs,
design drawings, as-built plans) documenting the riverine flow
structure.

\chapter{Water Control Structure}\label{water-control-structure}

\section{Spatial \& Geometric
Definition}\label{spatial-geometric-definition-14}

\begin{itemize}
\item
  Geometry Type: Point or Polygon (Water\_Control\_Structure\_Location).
\item
  Spatial Delineation: Footprint encompassing reinforced concrete
  headwalls, drop inlets, culvert barrels, slide gates, stoplog bays,
  and energy dissipation aprons.
\item
  Structural Typology: Gravity flow stoplog structures, flap-gated
  culverts, screw-jack sluice gates, and radial gates.
\end{itemize}

\subsection{Primary Restoration \& Engineering
Objectives}\label{primary-restoration-engineering-objectives-13}

\begin{itemize}
\item
  Water Surface Regulation: Provide bi-directional water-level control
  across moist soil cells, backwater lakes, and floodplain basins.
\item
  Bi-directional Flow Management: Retain managed pools during low river
  stages, permit rapid gravity drainage, and allow controlled gravity
  inflow during high-water pulses.
\item
  Fish Passage Control: Incorporate low-velocity bays or operational
  settings that permit native fish ingress while blocking rough fish
  (common carp).
\end{itemize}

\subsection{Typical Structural \& Hydraulic Engineering
Criteria}\label{typical-structural-hydraulic-engineering-criteria}

\begin{itemize}
\item
  Stoplog Bays: Multi-bay reinforced concrete structures (4 to 6 ft bay
  widths) with embedded stainless steel or aluminum slide channels
  housing treated timber or structural aluminum stoplogs.
\item
  Culvert Barrels: Heavy-duty reinforced concrete pipe (RCP) or
  structural aluminum/HDPE conduits (24 to 48 in diameter) passing
  through embankment bodies.
\item
  Seepage Cutoff Collars: Minimum of two concrete anti-seep collars or
  bentonite cutoff rings placed around culvert barrels within the
  embankment core to prevent internal seepage piping.
\item
  Scour Dissipation Aprons: Concrete stilling basins or Class III riprap
  aprons underlain by geotextile at structure outlets to absorb exit
  velocities (v\textgreater5.0 ft/s).
\end{itemize}

\subsection{Construction \& Implementation
Methodology}\label{construction-implementation-methodology-14}

\begin{itemize}
\item
  Cast-in-Place Concrete Construction: Structural concrete engineered
  with minimum 28-day compressive strengths of f c′ ≥4,000 psi and
  air-entrainment (5\% to 8\%) for extreme freeze-thaw durability.
\item
  Sheetpile Cutoff Walls: Steel sheetpile cutoff walls driven 10 to 20
  ft beneath concrete aprons to prevent under-seepage and structural
  piping.
\end{itemize}

\subsection{Quality Control, Inspection \& Long-Term
Maintenance}\label{quality-control-inspection-long-term-maintenance-14}

\begin{itemize}
\tightlist
\item
  Stoplog Seal Integrity: Annual inspection of stoplog rubber seals,
  lifting tracks, and slide gate screw-jacks. Lubrication of gear
  assemblies and clearing of silt accumulation within stoplog guide
  slots.
\end{itemize}

\section{Spatial Features Attributes}\label{spatial-features-attributes}

\subsection{Definition}\label{definition-13}

Water control structures are structures within a water management system
that convey water, control the direction or rate of flow, maintain a
desired water surface elevation, or measure water levels or flows.

\subsection{Purpose}\label{purpose-16}

Water control structures are used to regulate water levels, flow rates,
and distribution within managed areas such as MSMUs, backwater lakes,
and wetlands. They can be used to maintain or increase water depths to
create open water areas, periodically drain or flood wetlands to enhance
plant germination, and provide operational flexibility to meet multiple
habitat objectives (e.g., for waterfowl, shorebirds, wading birds,
reptiles, amphibians, and fisheries).

\subsection{Applicability}\label{applicability-14}

Water control structures are applicable wherever active water level or
flow management is required in Upper Mississippi River System (UMRS)
restoration projects. They are commonly used to manage wetlands,
backwaters, and managed floodplain units, and are integral components of
water level management strategies implemented under HREP/UMRR projects.
(UMRR Design Handbook, 2012)

\subsection{Notes}\label{notes-16}

Restoration features now have a Program\_UID field to distinguish
between NESP and UMRR features. Aliases without program identification
will be created for any field that currently has UMRR in the name, so
they can be used for NESP or UMRR data.

Attribute criteria for each field includes the following:

\begin{itemize}
\tightlist
\item
  Attribute Status -- defines whether the field requires a value or if
  it is optional based on available information.
\item
  Phases -- indicates the earliest time in the project lifecycle values
  for this attribute are available.
\item
  Digitizing requirements -- indicates when the field value should be
  entered in the attribute table.
\end{itemize}

Geometry Note: No Z values are needed for this feature.

\subsection{WaterControlStructure}\label{watercontrolstructure}

(point) -- A structure designed to control the flow of water from the
interior of an embankment to the waterside or from the waterside to the
interior.

\section{\texorpdfstring{\protect\includegraphics[width=0.75\linewidth,height=\textheight,keepaspectratio]{images/watercontrolstructure.jpg}}{}}\label{section-3}

\section{Feature Attribute
Criteria}\label{feature-attribute-criteria-14}

\#

Attribute

Attribute Status

Phase

Digitizing Requirements

1

Program\_UID

Required

Feasibility

At feature creation

2

PROJECT\_ID (Alias)

Required

Feasibility

At feature creation

3

Type (Alias)

Required

Feasibility

At feature creation

4

Subtype (Alias)

Required

Feasibility

At feature creation

5

RestorFeatStatus

Required

Feasibility

At feature creation

6

sdsFeatureName

Optional

Design

As required and available

7

sdsFeatureDescription

Optional

Design

As required and available

8

RestorFeatVdatum

Optional

Design

As required and available

9

ControlMethod

Optional

Design

As required and available

10

StructureCapacityCFS

Optional

Design

As required and available

11

LengthFT

Optional

Design

As required and available

12

DiameterIN

Optional

Design

As required and available

13

StructureMaterial

Optional

Design

As required and available

14

InletInvertElevation

Optional

Construction

As required and available

15

OutletInvertElevation

Optional

Construction

As required and available

16

BarrelNumber

Optional

Construction

As required and available

17

OperationType

Optional

Construction

As required and available

18

StructureType

Optional

Construction

As required and available

19

BottomElevation

Optional

Construction

As required and available

20

Width

Optional

Construction

As required and available

21

Height

Optional

Construction

As required and available

22

Station

Optional

Construction

As required and available

23

sdsMetaDataID

Optional

Construction

As required and available

24

RestorFeatIDPK

Optional

Construction

As required and available

25

sdsID

Optional

Construction

As required and available

26

mediaIDFK

Optional

Construction

As required and available

\begin{center}\rule{0.5\linewidth}{0.5pt}\end{center}

\section{Parameters}\label{parameters-14}

\subsection{Program\_UID}\label{program_uid-14}

Designate program NESP or UMRR.

\subsection{Project\_ID (alias:
PROJECT\_ID)}\label{project_id-alias-project_id-14}

Foreign key to the unique NESP/UMRR project ID.

\subsection{Type (alias)}\label{type-alias-13}

A broad classification used to group NESP/UMRR restoration features into
categories of similar function.

Value

Description

Water Level Management

Measures taken to imitate natural low flow river conditions with the
goal of improving wetland habitats.

\subsection{SubType (alias: Subtype)}\label{subtype-alias-subtype-14}

A classification used to group NESP/UMRR restoration features into
specific design feature categories.

Value

Description

Gravity Drain

Gravity drain feature consisting of a discharge pipe and an
opening/closing mechanism, used to introduce or remove water from a MSMU
or backwater lake.

Inlet Structure

Large structure, generally concrete, used to allow large amounts of
water into a MSMU or backwater area.

Outlet Structure

Large structure, generally concrete, used to allow large amounts of
water out of a MSMU or backwater area.

\subsection{RestorFeatStatus}\label{restorfeatstatus-14}

A description of the status of the individual restoration. It also
includes the party responsible for the individual feature.

Value

Description

USACE Planned

Feature is in a USACE planning stage.

USACE Built

Feature has been built by USACE. If a feature has been planned and built
by USACE, once constructed, it will always be considered USACE built
even if it has been repaired under later project stages.

USACE Enhanced

The function of an existing feature that was not built by USACE has been
enhanced by a feature that was USACE built or repaired.

USACE Repaired

An existing feature that was not built by USACE has been directly
improved or repaired.

Other Planned

A feature that is planned by another organization or government agency.

Other Built

Feature has been built by another organization or government agency and
is an integral part of the UMRR project.

Other Enhanced

The function of an existing feature, that was built by another
organization or government agency, has been enhanced by a feature that
was built or repaired by a non-USACE organization or government agency.

Other Repaired

An existing feature that has been directly improved or repaired by
another organization or government agency.

\subsection{sdsFeatureName}\label{sdsfeaturename-14}

The common name of the feature.

\subsection{sdsFeatureDescription}\label{sdsfeaturedescription-14}

A narrative describing the feature.

\subsection{RestorFeatVdatum}\label{restorfeatvdatum-14}

Vertical datum of elevations recorded in the attribute table or used as
feature z values.

Value

Description

MSL 1912

Mean Sea Level of 1912

NGVD 1929

National Geodetic Vertical Datum of 1929

NAVD 1988

North American Vertical Datum of 1988

Other

Other, note datum in sdsFeatureDescription

Unknown

Unknown vertical datum

\subsection{ControlMethod}\label{controlmethod}

The type of flow control used.

Value

Description

Flap Gate

Top-hinged gate, can be manually or automatically operated.

Gate Wells

Concrete enclosure used to house a manually operated drain closure valve
or another device.

Sluice Gate

Gate lowered into place, generally using a threaded stem system,
operated manually or automatically.

Flap gate and sluice gate

Combination flap and sluice gate.

Stoplog

Stoplog structure.

Weir Gate

Gate that slides down to allow water to flow over the top of the gate.

Other

Other gate type; specify in sdsFeatureDescription.

NA

Not applicable.

Unknown

Control method unknown.

\subsection{StructureCapacityCFS}\label{structurecapacitycfs}

Design capacity of the structure in cubic feet per second (CFS).

\subsection{LengthFT}\label{lengthft}

Length of the gravity drain or pipe in feet.

\subsection{DiameterIN}\label{diameterin}

Diameter of the drainpipe in inches.

\subsection{StructureMaterial}\label{structurematerial}

Material the structure (typically barrel/pipe) is made of.

Value

Description

Aluminum

Aluminum pipe.

Concrete

Concrete pipe.

Fiberglass

Fiberglass pipe.

Steel

Steel pipe.

PVC

Polyvinyl chloride pipe.

Corrugated metal

Corrugated metal pipe.

Copper

Copper pipe.

HDPE

High density polyethylene pipe.

Other

Other material; specify in sdsFeatureDescription.

NA

Not applicable.

Unknown

Material unknown.

\subsection{InletInvertElevation}\label{inletinvertelevation}

Elevation of the inlet invert (inside bottom) of the structure,
referenced to RestorFeatVdatum.

\subsection{OutletInvertElevation}\label{outletinvertelevation}

Elevation of the outlet invert (inside bottom) of the structure,
referenced to RestorFeatVdatum.

\subsection{BarrelNumber}\label{barrelnumber}

Number of barrels, if the gravity drain is a multi-barrel drain.

\subsection{OperationType}\label{operationtype}

Type of operation required for gate closing or opening.

Value

Description

Manual

Manual operation required.

Automatic

Automatic operation required.

Both

Manual and automatic operation can be used.

Unknown

Operation unknown.

NA

Not applicable.

\subsection{StructureType}\label{structuretype}

Type of water control structure.

Value

Description

Pipe

Pipe-type structure.

Gate

Gate structure (non-pipe).

Box Culvert

Box culvert (non-pipe).

Other

Other structure type; specify in sdsFeatureDescription.

\subsection{BottomElevation}\label{bottomelevation}

Bottom elevation of non-pipe type structures, referenced to
RestorFeatVdatum.

\subsection{Width}\label{width}

Width of non-pipe type structures in feet.

\subsection{Height}\label{height}

Height of non-pipe type structures in feet.

\subsection{Station}\label{station-3}

Stationing of the feature, typically referenced to a project baseline or
alignment (e.g., 10+00).

\subsection{sdsMetaDataID}\label{sdsmetadataid-14}

References attached metadata.

\subsection{RestorFeatIDPK}\label{restorfeatidpk-14}

Primary key used to uniquely identify the restoration feature record.

\subsection{sdsID}\label{sdsid-14}

Identifier used to link the feature to supporting SDS (spatial data
system) records or related datasets.

\subsection{mediaIDFK}\label{mediaidfk-14}

Foreign key used to reference associated media (e.g., structure details,
as-built drawings, photographs) documenting the water control structure.

\chapter{Well}\label{well}

\section{Spatial \& Geometric
Definition}\label{spatial-geometric-definition-15}

\begin{itemize}
\item
  Geometry Type: Point (Groundwater\_Well\_Location).
\item
  Spatial Delineation: Point location of the drilled wellhead, surface
  sanitary seal, power connection, and discharge connection pipe.
\item
  Hydrogeologic Target: Deep unconsolidated alluvial sand and gravel
  aquifers within the modern river floodplain.
\end{itemize}

\subsection{Primary Restoration \& Engineering
Objectives}\label{primary-restoration-engineering-objectives-14}

\begin{itemize}
\item
  Sediment-Free Water Delivery: Supply clean, silt-free groundwater to
  moist soil units, eliminating the high turbidity and sediment loads
  inherent to river surface water.
\item
  Drought-Proof Water Supply: Provide reliable filling capacity for
  managed waterfowl impoundments during severe late-summer droughts when
  river stages are depressed.
\item
  Winter Open-Water Maintenance: Discharge warm groundwater (50 ∘ F to
  55 ∘ F) into impoundment pockets during midwinter freezes to maintain
  open-water foraging zones for wintering waterfowl.
\end{itemize}

\subsection{Typical Hydrogeologic \& Mechanical
Specifications}\label{typical-hydrogeologic-mechanical-specifications}

\begin{itemize}
\item
  Drilling Depths: Typically drilled 80 to 250 ft into basal alluvial
  gravel formations.
\item
  Casing \& Screen Geometry: Heavy-duty steel or PVC casing (12 to 20 in
  diameter) with continuous-slot stainless steel well screens (0.040- to
  0.080-in slot sizes) sized to match aquifer grain-size distributions.
\item
  Gravel Pack: Graded silica sand/gravel filter pack placed around the
  screen to stop fine sand intrusion.
\item
  Yield \& Mechanics: High-capacity line-shaft or submersible turbine
  pumps delivering yields of 1,000 to 4,000 gpm (2.2 to 8.9 cfs).
\end{itemize}

\subsection{Construction \& Implementation
Methodology}\label{construction-implementation-methodology-15}

\begin{itemize}
\item
  Mud-Rotary or Reverse-Circulation Drilling: Boreholes drilled into
  unconsolidated alluvium using heavy truck-mounted drilling rigs.
\item
  Well Development: High-energy surge-block surging, air-jetting, and
  continuous over-pumping until discharge water tests free of suspended
  sand grains (\textless5 mg/L).
\end{itemize}

\subsection{Quality Control, Inspection \& Long-Term
Maintenance}\label{quality-control-inspection-long-term-maintenance-15}

\begin{itemize}
\item
  Step-Drawdown Aquifer Testing: 24- to 72-hour continuous pump tests to
  establish baseline specific capacity (gpm/ft of drawdown) and ensure
  long-term sustainability without well-field interference.
\item
  Water Quality Testing: Routine chemical analysis testing for dissolved
  iron, manganese, and total dissolved solids (TDS) to track potential
  screen encrustation.
\end{itemize}

\section{Spatial Feature Attribute}\label{spatial-feature-attribute-10}

\subsection{Definition}\label{definition-14}

A well is a shaft dug or drilled into the Earth for the purpose of
extracting fluids from the subsurface, collecting environmental samples,
injecting fluids into the subsurface, or extracting contamination or
other impurities from the subsurface. Wells typically operate in
conjunction with pumps and associated infrastructure.

\subsection{Purpose}\label{purpose-17}

Wells are used to manage groundwater conditions, provide water supply,
support environmental monitoring, and implement remediation or treatment
measures in support of restoration projects. They enable controlled
introduction or removal of water or other fluids, as well as sampling
and observation, to inform design, operations, and adaptive management
of project features.

\subsection{Applicability}\label{applicability-15}

Water can be introduced or removed from a MSMU or backwater lake through
the use of a well. The volume of water required will generally dictate
whether a groundwater well can be feasibly constructed. Groundwater
wells are limited in capacity due to available well yield from the
aquifer, construction limitations, commercially available well pump
size, and availability of utility power. Wells are applicable wherever
groundwater interaction is a key component of water level management,
habitat restoration, contaminant control, or environmental monitoring
within Upper Mississippi River System (UMRS) restoration projects. (UMRR
Design Handbook, 2012)

\subsection{Notes}\label{notes-17}

Restoration features now have a Program\_UID field to distinguish
between NESP and UMRR features. Aliases without program identification
will be created for any field that currently has UMRR in the name, so
they can be used for NESP or UMRR data.

Attribute criteria for each field includes the following:

\begin{itemize}
\tightlist
\item
  Attribute Status -- defines whether the field requires a value or if
  it is optional based on available information.
\item
  Phases -- indicates the earliest time in the project lifecycle values
  for this attribute are available.
\item
  Digitizing requirements -- indicates when the field value should be
  entered in the attribute table.
\end{itemize}

Geometry Note: No Z values are needed for this feature.

\subsection{Well}\label{well-1}

(point) -- A shaft dug or drilled into the Earth for the purpose of
extracting fluids from the subsurface, collecting environmental samples,
injecting fluids into the subsurface or extracting contamination or
other impurities from the subsurface.

\begin{center}\rule{0.5\linewidth}{0.5pt}\end{center}

\section{Feature Attribute
Criteria}\label{feature-attribute-criteria-15}

\#

Attributes

Attribute Status

Phase

Digitizing Requirement

1

Program\_UID

Required

Feasibility

At feature creation

2

PROJECT\_ID

Required

Feasibility

At feature creation

3

Type

Required

Feasibility

At feature creation

4

Subtype

Required

Feasibility

At feature creation

5

RestorFeatStatus

Required

Design

At feature creation

6

sdsFeatureName

Optional

Design

As required and available

7

sdsFeatureDescription

Optional

Design

As required and available

8

RestorFeatVdatum

Optional

Design

As required and available

9

WellPurpose

Optional

Design

As required and available

10

WellResource

Optional

Construction

As required and available

11

WellTopElev

Optional

Construction

As required and available

12

Station

Optional

Construction

As required and available

13

sdsMetaDataID

Optional

Construction

As required and available

14

RestorFeatIDPK

Optional

Construction

As required and available

15

sdsID

Optional

Construction

As required and available

16

mediaIDFK

Optional

Construction

As required and available

\begin{center}\rule{0.5\linewidth}{0.5pt}\end{center}

\section{Parameters}\label{parameters-15}

\subsection{Program\_UID}\label{program_uid-15}

Designate program NESP or UMRR.

\subsection{Project\_ID (alias:
PROJECT\_ID)}\label{project_id-alias-project_id-15}

Foreign key to the unique NESP/UMRR project ID.

\subsection{Type (alias)}\label{type-alias-14}

A broad classification used to group NESP/UMRR restoration features into
categories of similar function.

Value

Description

Water Level Management

Measures taken to manage or influence water levels and hydrologic
conditions for habitat and floodplain function.

\subsection{SubType (alias: Subtype)}\label{subtype-alias-subtype-15}

A classification used to group NESP/UMRR restoration features into
specific design feature categories.

Value

Description

Well

Vertical hole drilled into the ground to access or influence groundwater
movement, used to introduce or remove water from a MSMU or backwater
lake.

\subsection{RestorFeatStatus}\label{restorfeatstatus-15}

A description of the status of the individual restoration. It also
includes the party responsible for the individual feature.

Value

Description

USACE Planned

Feature is in a USACE planning stage.

USACE Built

Feature has been built by USACE. If a feature has been planned and built
by USACE, once constructed, it will always be considered USACE built
even if it has been repaired under later project stages.

USACE Enhanced

The function of an existing feature that was not built by USACE has been
enhanced by a feature that was USACE built or repaired.

USACE Repaired

An existing feature that was not built by USACE has been directly
improved or repaired.

Other Planned

A feature that is planned by another organization or government agency.

Other Built

Feature has been built by another organization or government agency and
is an integral part of the UMRR project.

Other Enhanced

The function of an existing feature, that was built by another
organization or government agency, has been enhanced by a feature that
was built or repaired by a non-USACE organization or government agency.

Other Repaired

An existing feature that has been directly improved or repaired by
another organization or government agency.

\subsection{sdsFeatureName}\label{sdsfeaturename-15}

The common name of the feature.

\subsection{sdsFeatureDescription}\label{sdsfeaturedescription-15}

A narrative describing the feature.

\subsection{RestorFeatVdatum}\label{restorfeatvdatum-15}

Vertical datum of elevations recorded in the attribute table or used as
feature z values.

Value

Description

MSL 1912

Mean Sea Level of 1912

NGVD 1929

National Geodetic Vertical Datum of 1929

NAVD 1988

North American Vertical Datum of 1988

Other

Other, note datum in sdsFeatureDescription

Unknown

Unknown vertical datum

\subsection{WellPurpose}\label{wellpurpose}

Purpose of use for the well.

Value

Description

Extraction

Well used to withdraw groundwater or other fluids from the subsurface.

Injection

Well used to introduce fluids (e.g., water, treatment solutions) into
the subsurface.

Monitoring

Well used for observation and sampling (e.g., groundwater level,
quality, or contaminant monitoring).

NA

Not applicable.

TBD

To be determined.

\subsection{WellResource}\label{wellresource}

The resource type which is being extracted from or injected into the
earth.

Value

Description

Air

Air.

NA

Not applicable.

Natural\_gas

Natural gas.

None

No resource currently associated (e.g., inactive well).

oil

Oil.

Steam

Steam.

Sulphur

Sulphur.

TBD

To be determined.

Vapor

Vapor.

Water

Water.

\subsection{WellTopElev}\label{welltopelev}

Elevation at the top of the well (e.g., top of casing or protective
cover), referenced to the vertical datum specified in RestorFeatVdatum.

\subsection{Station}\label{station-4}

Stationing of the feature, typically referenced to a project baseline or
alignment (e.g., 10+00).

\subsection{sdsMetaDataID}\label{sdsmetadataid-15}

References attached metadata.

\subsection{RestorFeatIDPK}\label{restorfeatidpk-15}

Primary key used to uniquely identify the restoration feature record.

\subsection{sdsID}\label{sdsid-15}

Identifier used to link the feature to supporting SDS (spatial data
system) records or related datasets.

\subsection{mediaIDFK}\label{mediaidfk-15}

Foreign key used to reference associated media (e.g., well construction
logs, pumping test records, photographs, as-built drawings) documenting
the well.

\chapter{Wetland}\label{wetland}

\section{Spatial \& Geometric
Definition}\label{spatial-geometric-definition-16}

\begin{itemize}
\item
  Geometry Type: Polygon (Wetland\_Restoration\_Polygon).
\item
  Spatial Delineation: Shallow-water aquatic marsh and bottomland
  wetland zones restored through grading, embankment construction, or
  river stage management.
\item
  Wetland Classes: Emergent hemi-marsh, deep-water marsh, scrub-shrub
  wetland, and wet meadow.
\end{itemize}

\subsection{Primary Restoration \& Engineering
Objectives}\label{primary-restoration-engineering-objectives-15}

\begin{itemize}
\item
  Hemi-Marsh Restoration: Create balanced 50:50 ratios of dense emergent
  vegetation to open-water pools, maximizing wildlife diversity and
  nesting cover.
\item
  Nutrient \& Silt Cycling: Serve as natural bio-retention filters that
  sequester suspended sediment, nitrogen, and phosphorus from passing
  floodwaters.
\item
  Amphibian \& Reptile Nursery: Provide predator-sheltered, warm, sunlit
  nursery waters for amphibians, shorebirds, and native aquatic
  reptiles.
\end{itemize}

\subsection{Typical Cross-Sections, Dimensions \& Sizing
Criteria}\label{typical-cross-sections-dimensions-sizing-criteria-10}

\begin{itemize}
\item
  Water Depth Zonation:Littoral Shelf (0.0 to 1.0 ft depth): Supports
  emergent marsh species (Typha latifolia, Sagittaria latifolia, Scirpus
  spp.).
\item
  Intermediate Zone (1.0 to 3.0 ft depth): Promotes diverse beds of
  submersed aquatic vegetation (SAV) (Vallisneria americana, Potamogeton
  spp.).
\item
  Deep Pools (4.0 to 6.0 ft depth): Overwintering survival pockets for
  turtles and amphibians.
\item
  Micro-Topographic Diversity: Irregular curving shorelines, isolated
  mud islands, and variable bottom elevations to eliminate uniform
  bathymetry.
\end{itemize}

\subsection{Construction \& Implementation
Methodology}\label{construction-implementation-methodology-16}

\begin{itemize}
\item
  Low-Ground-Pressure Earthmoving: Wide-track bulldozers (\textless3 psi
  ground pressure) shaping undulating bottoms and shallow perimeter
  benches.
\item
  Vegetation Transplants \& Hydroseeding: Installation of pre-rooted
  emergent wetland plugs and broadcast wetland seed mixes immediately
  following final earthwork.
\end{itemize}

\subsection{Quality Control, Inspection \& Long-Term
Maintenance}\label{quality-control-inspection-long-term-maintenance-16}

\begin{itemize}
\item
  Hydro-Period Monitoring: Continuous pressure transducer stage
  recorders logging water depths throughout the year to confirm
  hydrologic permanence.
\item
  Invasive Species Eradication: Rigorous annual mechanical mowing and
  aquatic herbicide treatments targeting narrowleaf cattail and purple
  loosestrife.
\end{itemize}

\section{Spatial Feature Attribute}\label{spatial-feature-attribute-11}

\subsection{Definition}\label{definition-15}

Geospatial feature composed of one unique wetland class (no complexes),
used to determine USACE jurisdiction and regulatory impacts. Each
feature represents a single, homogeneous wetland type suitable for
regulatory and design analysis.

\subsection{Purpose}\label{purpose-18}

Wetland features document the location, extent, and type of wetlands
affected by or created through restoration projects. They support
regulatory determinations, impact assessments, mitigation planning, and
design of floodplain restoration measures intended to improve natural
floodplain function and habitat.

\subsection{Applicability}\label{applicability-16}

Wetland features are applicable wherever wetlands occur within Upper
Mississippi River System (UMRS) restoration project areas or where
wetlands are created, restored, or enhanced as part of project actions.
They are used to assess regulatory jurisdiction, quantify impacts and
benefits, and guide design and management of floodplain restoration
measures. (UMRR Design Handbook, 2012)

\subsection{Notes}\label{notes-18}

Restoration features now have a Program\_UID field to distinguish
between NESP and UMRR features. Aliases without program identification
will be created for any field that currently has UMRR in the name, so
they can be used for NESP or UMRR data.

Attribute criteria for each field includes the following:

\begin{itemize}
\tightlist
\item
  Attribute Status -- defines whether the field requires a value or if
  it is optional based on available information.
\item
  Phases -- indicates the earliest time in the project lifecycle values
  for this attribute are available.
\item
  Digitizing requirements -- indicates when the field value should be
  entered in the attribute table.
\end{itemize}

Geometry Note: No Z values are needed for this feature.

\subsection{Wetland}\label{wetland-1}

(polygon) -- Geospatial feature composed of one unique wetland class (no
complexes) used to determine USACE jurisdiction and regulatory impacts.

\begin{center}\rule{0.5\linewidth}{0.5pt}\end{center}

\section{Feature Attribute
Criteria}\label{feature-attribute-criteria-16}

\#

Attributes

Attribute Status

Phase

Digitizing Requirement

1

Program\_UID

Required

Feasibility

At feature creation

2

PROJECT\_ID

Required

Feasibility

At feature creation

3

Type

Required

Feasibility

At feature creation

4

Subtype

Required

Feasibility

At feature creation

5

RestorFeatStatus

Required

Design

At feature creation

6

sdsFeatureName

Optional

Design

As required and available

7

sdsFeatureDescription

Optional

Design

As required and available

8

RestorFeatVdatum

Optional

Design

As required and available

9

HydroGeoMorphType

Optional

Design

As required and available

10

DepthMaxFT

Optional

Construction

As required and available

11

DepthAverageFT

Optional

Construction

As required and available

12

Acres

Optional

Construction

As required and available

13

sdsMetaDataID

Optional

Construction

As required and available

14

RestorFeatIDPK

Optional

Construction

As required and available

15

sdsID

Optional

Construction

As required and available

16

mediaIDFK

Optional

Construction

As required and available

\begin{center}\rule{0.5\linewidth}{0.5pt}\end{center}

\section{Parameters}\label{parameters-16}

\subsection{Program\_UID}\label{program_uid-16}

Designate program NESP or UMRR.

\subsection{Project\_ID (alias:
PROJECT\_ID)}\label{project_id-alias-project_id-16}

Foreign key to the unique NESP/UMRR project ID.

\subsection{Type (alias)}\label{type-alias-15}

A broad classification used to group NESP/UMRR restoration features into
categories of similar function.

Value

Description

Floodplain Restoration

Measures taken to improve natural floodplain function and habitat.

Island Creation

Measures taken to create or restore islands in rivers.

\subsection{SubType (alias: Subtype)}\label{subtype-alias-subtype-16}

A classification used to group NESP/UMRR restoration features into
specific wetland design feature categories.

Value

Description

Depressional wetland

Wetland occurring in a topographic depression, typically ponded or
seasonally flooded, with limited surface outlet and water supplied by
precipitation, runoff, and/or groundwater.

Mudflat

Shallow, frequently exposed wetland surface characterized by fine
sediments and periodic inundation, often used by shorebirds and other
wildlife.

Swale

Shallow, elongated wetland feature that conveys or temporarily stores
water, often oriented along low-lying topographic features within the
floodplain.

\subsection{RestorFeatStatus}\label{restorfeatstatus-16}

A description of the status of the individual restoration. It also
includes the party responsible for the individual feature.

Value

Description

USACE Planned

Feature is in a USACE planning stage.

USACE Built

Feature has been built by USACE. If a feature has been planned and built
by USACE, once constructed, it will always be considered USACE built
even if it has been repaired under later project stages.

USACE Enhanced

The function of an existing feature that was not built by USACE has been
enhanced by a feature that was USACE built or repaired.

USACE Repaired

An existing feature that was not built by USACE has been directly
improved or repaired.

Other Planned

A feature that is planned by another organization or government agency.

Other Built

Feature has been built by another organization or government agency and
is an integral part of the UMRR project.

Other Enhanced

The function of an existing feature, that was built by another
organization or government agency, has been enhanced by a feature that
was built or repaired by a non-USACE organization or government agency.

Other Repaired

An existing feature that has been directly improved or repaired by
another organization or government agency.

\subsection{sdsFeatureName}\label{sdsfeaturename-16}

The common name of the feature.

\subsection{sdsFeatureDescription}\label{sdsfeaturedescription-16}

A narrative describing the feature.

\subsection{RestorFeatVdatum}\label{restorfeatvdatum-16}

Vertical datum of elevations recorded in the attribute table or used as
feature z values.

Value

Description

MSL 1912

Mean Sea Level of 1912

NGVD 1929

National Geodetic Vertical Datum of 1929

NAVD 1988

North American Vertical Datum of 1988

Other

Other, note datum in sdsFeatureDescription

Unknown

Unknown vertical datum

\subsection{HydroGeoMorphType}\label{hydrogeomorphtype}

The hydrogeomorphology of the wetland.

Value

Description

Depressional

Characterized by a water source consisting of return flow from
groundwater and interflow with primarily vertical hydrodynamics.

Estuarine Fringe

Water source consists of overbank flow from estuaries, with
bidirectional and horizontal hydrodynamics being dominant.

Lacustrine Fringe

Dominant water source is lake overbank flow, and the dominant
hydrodynamics are bidirectional and horizontal.

Mineral Soil Flats

Water source is precipitation, and vertical hydrodynamics are dominant.

Organic Soil Flats

Water source is precipitation, and the dominant hydrodynamics are
vertical.

Riverine

Characterized by a water source of overbank flow from a channel, with
predominantly unidirectional and horizontal hydrodynamics.

Slope

Characterized by a water source of return flow from groundwater, with
principally unidirectional and horizontal hydrodynamics.

Other

Other hydrogeomorphic setting; specify in sdsFeatureDescription.

NA

Not applicable.

\subsection{DepthMaxFT}\label{depthmaxft}

Maximum depth in feet of the constructed wetland.

\subsection{DepthAverageFT}\label{depthaverageft}

Average depth in feet of the constructed wetland.

\subsection{Acres}\label{acres-4}

Area of the wetland in acres.

\subsection{sdsMetaDataID}\label{sdsmetadataid-16}

References attached metadata.

\subsection{RestorFeatIDPK}\label{restorfeatidpk-16}

Primary key used to uniquely identify the restoration feature record.

\subsection{sdsID}\label{sdsid-16}

Identifier used to link the feature to supporting SDS (spatial data
system) records or related datasets.

\subsection{mediaIDFK}\label{mediaidfk-16}

Foreign key used to reference associated media (e.g., photographs,
design drawings, as-built plans) documenting the wetland.

\part{Restoration Support Features}

\chapter{Restoration Support
Features}\label{restoration-support-features-1}

\section{Purpose}\label{purpose-19}

Restoration Support Features are a topical group of feature classes used
to represent, at various phases, engineering features constructed or
implemented to directly meet the goals and objectives of programs and
the projects they support.

\section{Notes}\label{notes-19}

Data Sources - If a project has been constructed, CAD drawings or
As-Built drawings should be used to guide feature digitizing. If a
project is in the planning stages, the schema can be used to record
potential features and track which planning alternatives they
participate in. Elevation data should be attached to geometries as
specified in this document.

These feature classes are stored in the MVR UMRR Oracle database to
ensure that authoritative data is available to all. See these reference
guides to access the oracle database for editing features.

\section{Restoration Support
Features}\label{restoration-support-features-2}

\begin{itemize}
\tightlist
\item
  DredgingPlacement (polygon) - The extent of the area covered by
  dredged material.
\item
  BoreholePoint (point) -- A shaft drilled, driven or dug into the earth
  for the purpose of subsurface exploration.
\item
  ElectricalUtilityNode (point) -- A subdivision of an electrical
  distribution network, particularly an asset that participates in the
  transmission of electricity but is not a line. For example, a
  transformer, fuse, generator, meter, or switch.
\item
  ElectricalUtilitySegment (line) -- A subdivision of an electrical
  distribution network, particularly a line for the transmission of
  electricity.
\item
  FeatureStationing (point) -- Reference points, generally in feet, that
  allow identification of specific locations on engineered structures.
\item
  PavementSection (polygon) -- A pavement section is a portion of a
  pavement branch that differs in some aspect from other sections such
  that further segmentation is required to uniquely identify that
  section. For example, if the only difference in a road is the pavement
  material type -- concrete for part of the road and asphalt for the
  remainder of the road -- that road branch has two sections.
\item
  RoadSegment (line) -- A portion of a road with similar characteristics
  such as the number of lanes or surface type.
\item
  Sign (point) -- A structure that conveys directional, warning or other
  information.
\item
  WaterUtilityNode (point) -- A subdivision of a water distribution
  network, particularly an asset that participates in the transmission
  of water but that is not a pipeline. For example, a pump, valve or
  fitting.
\item
  WaterUtilitySegment (line) -- A subdivision of a wastewater network,
  particularly a pipeline for the transport of stormwater, sewage or
  industrial waste between the source, holding facilities, and/or
  treatment facilities.
\end{itemize}

\chapter{Dredging Placement}\label{dredging-placement}

\section{Spatial \& Geometric
Definition}\label{spatial-geometric-definition-17}

\begin{itemize}
\item
  Geometry Type: Polygon (Dredge\_Placement\_Site).
\item
  Spatial Delineation: Bounded footprint of upland disposal sites,
  confined disposal facilities (CDFs), in-water containment cells, or
  marsh nourishment zones receiving dredged sediments.
\item
  Placement Typology: Terrestrial containment basins, in-water island
  fill cells, marsh thin-layer nourishment zones, and beneficial-use
  stockpiles.
\end{itemize}

\subsection{Primary Engineering \& Operational
Objectives}\label{primary-engineering-operational-objectives}

\begin{itemize}
\item
  Dredged Material Retention: Securely hold thousands of cubic yards of
  slurried or mechanically excavated silts and sands without
  catastrophic containment failure.
\item
  Effluent Water Clarification: Allow sufficient detention settling time
  for hydraulic slurries to drop suspended solids, ensuring return water
  meets strict Clean Water Act Section 401 state water quality
  standards.
\item
  Beneficial Reuse \& Topographic Elevation: Elevate low-lying
  floodplain terrains to create flood-refugia ridges, island cores, or
  agricultural soil enhancements.
\end{itemize}

\subsection{Typical Cross-Sections, Dimensions \& Sizing
Criteria}\label{typical-cross-sections-dimensions-sizing-criteria-11}

\begin{itemize}
\item
  Cell Capacities: Sized between 20,000 and 500,000+ cubic yards of
  material.
\item
  Perimeter Containment Dikes: Compacted earthen containment dikes
  (heights6 to 15 ft, side slopes 1V:3H) with minimum 3.0-ft interior
  freeboard.
\item
  Weir Dewatering Structures: Adjustable rectangular overflow weir boxes
  (morning-glory or stoplog weirs) sized to govern effluent decant rates
  based on Stokes' Law hydraulic settling velocities.
\end{itemize}

\subsection{Construction \& Implementation
Methodology}\label{construction-implementation-methodology-17}

\begin{itemize}
\item
  Hydraulic Discharge Management: Pipeline discharge placed at the
  furthest point from the overflow weir structure to maximize hydraulic
  residence time and particle settling.
\item
  Polymer / Flocculant Injection: When dredging fine colloidal clays,
  organic flocculants may be injected into the slurry line to accelerate
  settling.
\end{itemize}

\subsection{Quality Control, Inspection \& Long-Term
Maintenance}\label{quality-control-inspection-long-term-maintenance-17}

\begin{itemize}
\item
  Effluent Turbidity Monitoring: Continuous water testing at the weir
  discharge for Total Suspended Solids (TSS) and turbidity (NTU),
  maintaining compliance with state mixing-zone permits.
\item
  Dewatering \& Crusting Management: Post-dredging trenching with
  specialized marsh excavators to promote rapid drying, desiccation, and
  soil consolidation.
\end{itemize}

\section{Spatial Feature Attribute}\label{spatial-feature-attribute-12}

\subsection{Definition}\label{definition-16}

The extent of the area covered by dredged material.

\subsection{Purpose}\label{purpose-20}

Dredging placement features are used to document and manage the
location, extent, and characteristics of areas receiving dredged
material. Mapping these features supports design coordination,
environmental compliance, maintenance planning, and evaluation of
habitat restoration outcomes.

\subsection{Applicability}\label{applicability-17}

This applies to all dredging placement features (UMRR Design Handbook,
2012)

\subsection{Notes}\label{notes-20}

Restoration features now have a Program\_UID field to distinguish
between NESP and UMRR features. Aliases without program identification
will be created for any field that currently has UMRR in the name, so
they can be used for NESP or UMRR data.

Attribute criteria for each field includes the following:

\begin{itemize}
\tightlist
\item
  Attribute Status -- defines whether the field requires a value or if
  it is optional based on available information.
\item
  Phases -- indicates the earliest time in the project lifecycle values
  for this attribute as available.
\item
  Digitizing requirements -- indicates when the field value should be
  entered in the attribute table.
\item
  Geometry Note: No Z values are needed for this feature.
\end{itemize}

\subsection{DredgingPlacement}\label{dredgingplacement}

(polygon) - The extent of the area covered by dredged material. When the
use is Island, should we be putting polygons for both in the same area?

\section{Feature Attribute
Criteria}\label{feature-attribute-criteria-17}

\#

Attribute

Attribute Status

Phase

Digitizing Requirements

1

Program\_UID

Required

Feasibility

At feature creation

2

PROJECT\_ID

Required

Feasibility

At feature creation

3

Type

Required

Feasibility

At feature creation

4

Subtype

Required

Feasibility

At feature creation

5

RestorFeatStatus

Required

Design

At feature creation

6

sdsFeatureName

Optional

Design

As required and available

7

sdsFeatureDescription

Optional

Design

As required and available

8

RestorFeatVdatum

Optional

Design

As required and available

9

VolumeCY

Optional

Design

As required and available

10

Acres

Optional

Construction

As required and available

11

sdsMetaDataID

Optional

Construction

As required and available

12

RestorFeatIDPK

Optional

Construction

As required and available

13

sdsID

Optional

Construction

As required and available

14

mediaIDFK

Optional

Construction

As required and available

\section{Parameters}\label{parameters-17}

\subsection{Program\_UID}\label{program_uid-17}

Designate program NESP or UMRR.

\subsection{Project\_ID (alias:
PROJECT\_ID)}\label{project_id-alias-project_id-17}

Foreign key to the unique NESP/UMRR project ID.

\subsection{Type (alias)}\label{type-alias-16}

A broad classification used to group NESP/UMRR restoration features into
categories of similar function.

Value

Description

Dredge Placement

Add text for description.

\subsection{SubType (alias: Subtype)}\label{subtype-alias-subtype-17}

A classification used to group NESP/UMRR restoration features into
specific design feature categories.

Value

Description

Island

Island building.

Levee

Build a new levee or strengthen and existing levee as part of a moist
soil unit.

Floodplain Depth Diversity

Placement to increase floodplain topographic diversity and habitat.

Unknown

Undocumented use.

\subsection{RestorFeatStatus}\label{restorfeatstatus-17}

A description of the status of the individual restoration. It also
includes the party responsible for the individual feature.

Value

Description

USACE Planned

Feature is in a USACE planning stage.

USACE Built

Feature has been built by USACE. If a feature has been planned and built
by USACE, once constructed, will always be considered USACE built even
if it has been repaired under later project stages.

USACE Enhanced

The function of an existing feature that was not built by USACE has been
enhanced by a feature that was USACE built or repaired.

USACE Repaired

An existing feature that was not built by USACE has been directly
improved or repaired.

Other Planned

A feature that is planned by another organization or government agency.

Other Built

Feature has been built by another organization or government agency and
is a integral part of the UMRR project.

Other Enhanced

The function of an existing feature, that was built by another
organization or government agency, has been enhanced by a feature that
was built or repaired by a non-USACE organization or government agency.

Other Repaired

An existing feature that has been directly improved or repaired by
another organization or government agency.

\subsection{sdsFeatureName}\label{sdsfeaturename-17}

The common name of the feature.

\subsection{sdsFeatureDescription}\label{sdsfeaturedescription-17}

A narrative describing the feature.

\subsection{RestorFeatVdatum}\label{restorfeatvdatum-17}

Vertical datum of elevations recorded in attribute table or used as
feature z values.

Value

Description

MSL 1912

Mean Sea Level of 1912

NGVD 1929

National Geodetic Vertical Datum of 1929

NAVD 1988

North American Vertical Datum of 1988

Other

Other, note datum in sdsFeatureDescription

Unknown

Unknown vertical datum

\subsection{VolumeCY}\label{volumecy-2}

Volume of material removed in cubic yards.

\subsection{Acres}\label{acres-5}

Acreage of feature.

\subsection{sdsMetaDataID}\label{sdsmetadataid-17}

References attached metadata.

\subsection{RestorFeatIDPK}\label{restorfeatidpk-17}

References attached media.

\subsection{sdsID}\label{sdsid-17}

References attached media.

\subsection{mediaIDFK}\label{mediaidfk-17}

References attached media.

\chapter{Borehole Points}\label{borehole-points}

\section{Spatial \& Geometric
Definition}\label{spatial-geometric-definition-18}

\begin{itemize}
\item
  Geometry Type: Point (Geotechnical\_Borehole\_Point).
\item
  Spatial Delineation: Exact geographic surface coordinates of
  subsurface geotechnical exploration borings drilled during project
  investigation and design phases.
\item
  Associated Data Fields: Hole ID, surface ground elevation, depth to
  bedrock, depth to groundwater, sampling method, and laboratory test
  links.
\end{itemize}

\subsection{Primary Engineering \& Geotechnical
Objectives}\label{primary-engineering-geotechnical-objectives}

\begin{itemize}
\tightlist
\item
  Subsurface Stratigraphy Profiling: Delineate alluvium strata, identify
  compressible soft organic clay/silt layers, and map bedrock depth.
\end{itemize}

Foundation Bearing Capacity Assessment: Determine in-situ shear
strength, cohesion, and internal friction angles required to design
stable embankments, concrete pump stations, and stone dikes.

Seepage \& Settlement Modeling: Obtain soil samples for laboratory
testing to calculate consolidation settlement (C c \hspace{0pt} ),
compaction density, and hydraulic permeability (k).

Typical Drilling \& Testing Specifications Drilling Methodology:
Hollow-stem auger (HSA), mud-rotary drilling, or rock coring (NQ/HQ
wireline barrels) deployed from truck rigs or marine spud barges.

Standard Penetration Testing (SPT): Conducted per ASTM D1586 at
continuous 2.5- to 5.0-ft intervals, recording blow counts (N-values)
with standard 140-lb split-spoon hammers.

Thin-Walled Shelby Tube Sampling: Undisturbed cohesive soil samples
collected per ASTM D1587 for triaxial shear and consolidation testing.

Field Logging \& Data Delivery Standard Boring Logs: Field logs classify
all soil horizons according to the Unified Soil Classification System
(USCS) (e.g., CH, CL, ML, SP, SM).

Grouting \& Abandonment: All holes tremie-grouted from the bottom to the
ground surface with high-solids bentonite slurry per state well
abandonment regulations.

Quality Control \& Long-Term Data Archiving Enterprise GIS Integration:
All borehole logs digitized and georeferenced within the USACE MVR
Geoportal database, preserving geotechnical baseline archives for future
construction and maintenance.

\section{Definition}\label{definition-17}

A borehole point represents the location of a shaft drilled, driven, or
dug into the earth for the purpose of subsurface exploration, sampling,
or testing (e.g., geotechnical, environmental, or hydrogeologic
investigations).

\section{Purpose}\label{purpose-21}

Borehole points are used to document the location and basic
characteristics of subsurface investigations that support planning,
design, construction, and performance evaluations of restoration
features. They provide reference for logs, laboratory results, and
groundwater or soil data that inform foundation design, material
suitability, and risk assessment.

\section{Applicability}\label{applicability-18}

Borehole points are applicable wherever subsurface information is
collected within Upper Mississippi River System (UMRS) restoration
project areas, including investigations for levees, structures, habitat
features, borrow areas, disposal sites, and other engineered features.
They provide a consistent spatial reference for integrating subsurface
data with restoration planning and design. (UMRR Design Handbook, 2012)

\section{Notes}\label{notes-21}

Restoration features now have a Program\_UID field to distinguish
between NESP and UMRR features. Aliases without program identification
will be created for any field that currently has UMRR in the name, so
they can be used for NESP or UMRR data.

Attribute criteria for each field includes the following:

\begin{itemize}
\tightlist
\item
  Attribute Status -- defines whether the field requires a value or if
  it is optional based on available information.
\item
  Phases -- indicates the earliest time in the project lifecycle values
  for this attribute are available.
\item
  Digitizing requirements -- indicates when the field value should be
  entered in the attribute table.
\end{itemize}

Geometry Note: No Z values are needed for this feature.

\section{BoreholePoint}\label{boreholepoint}

(point) -- A shaft drilled, driven or dug into the earth for the purpose
of subsurface exploration.

\begin{center}\rule{0.5\linewidth}{0.5pt}\end{center}

\section{Feature Attribute
Criteria}\label{feature-attribute-criteria-18}

\#

Attributes

Attribute Status

Phase

Digitizing Requirement

1

Program\_UID

Required

Feasibility

At feature creation

2

PROJECT\_ID

Required

Feasibility

At feature creation

3

Type

Required

Feasibility

At feature creation

4

Subtype

Required

Feasibility

At feature creation

5

RestorFeatStatus

Required

Design

At feature creation

6

sdsFeatureName

Optional

Design

As required and available

7

sdsFeatureDescription

Optional

Design

As required and available

8

RestorFeatVdatum

Optional

Design

As required and available

9

BoringMethod

Optional

Design

As required and available

10

Date

Optional

Construction

As required and available

11

BoreholeDepthFT

Optional

Construction

As required and available

12

GroundElevation

Optional

Construction

As required and available

13

Station

Optional

Construction

As required and available

14

sdsMetaDataID

Optional

Construction

As required and available

15

RestorFeatIDPK

Optional

Construction

As required and available

16

sdsID

Optional

Construction

As required and available

17

mediaIDFK

Optional

Construction

As required and available

\begin{center}\rule{0.5\linewidth}{0.5pt}\end{center}

\section{Parameters}\label{parameters-18}

\subsection{Program\_UID}\label{program_uid-18}

Designate program NESP or UMRR.

\subsection{Project\_ID (alias:
PROJECT\_ID)}\label{project_id-alias-project_id-18}

Foreign key to the unique NESP/UMRR project ID.

\subsection{Type (alias)}\label{type-alias-17}

A broad classification used to group NESP/UMRR restoration features into
categories of similar function.

\begin{quote}
\textbf{Note:} Final coded value table for Type is to be defined (Needs
Table).
\end{quote}

\subsection{SubType (alias: Subtype)}\label{subtype-alias-subtype-18}

A classification used to group NESP/UMRR restoration features into
specific design feature categories.

\begin{quote}
\textbf{Note:} Final coded value table for SubType is to be defined
(Needs Table).
\end{quote}

\subsection{RestorFeatStatus}\label{restorfeatstatus-18}

A description of the status of the individual restoration. It also
includes the party responsible for the individual feature.

Value

Description

USACE Planned

Feature is in a USACE planning stage.

USACE Built

Feature has been built by USACE. If a feature has been planned and built
by USACE, once constructed, it will always be considered USACE built
even if it has been repaired under later project stages.

USACE Enhanced

The function of an existing feature that was not built by USACE has been
enhanced by a feature that was USACE built or repaired.

USACE Repaired

An existing feature that was not built by USACE has been directly
improved or repaired.

Other Planned

A feature that is planned by another organization or government agency.

Other Built

Feature has been built by another organization or government agency and
is an integral part of the UMRR project.

Other Enhanced

The function of an existing feature, that was built by another
organization or government agency, has been enhanced by a feature that
was built or repaired by a non-USACE organization or government agency.

Other Repaired

An existing feature that has been directly improved or repaired by
another organization or government agency.

\subsection{sdsFeatureName}\label{sdsfeaturename-18}

The common name of the feature.

\subsection{sdsFeatureDescription}\label{sdsfeaturedescription-18}

A narrative describing the feature.

\subsection{RestorFeatVdatum}\label{restorfeatvdatum-18}

Vertical datum of elevations recorded in the attribute table or used as
feature z values.

Value

Description

MSL 1912

Mean Sea Level of 1912

NGVD 1929

National Geodetic Vertical Datum of 1929

NAVD 1988

North American Vertical Datum of 1988

Other

Other, note datum in sdsFeatureDescription.

Unknown

Unknown vertical datum.

\subsection{BoringMethod}\label{boringmethod}

Method used to advance the borehole for subsurface exploration.

Value

Description

Hollow-stem Auger

Borehole advanced using hollow-stem augers.

Solid-stem Auger

Borehole advanced using solid-stem augers.

Rotary Wash

Borehole advanced using rotary drilling with fluid circulation.

Rotary Core

Borehole advanced using rotary coring methods.

Direct Push

Borehole created using direct-push/probing equipment.

Hand Auger

Shallow boring advanced manually using hand augers.

Test Pit / Excavation

Subsurface exploration completed by excavation rather than drilling.

Other

Other boring method; describe in sdsFeatureDescription.

Unknown

Boring method not recorded.

\subsection{Date}\label{date}

The date the boring was performed or completed.

\subsection{BoreholeDepthFT}\label{boreholedepthft}

Total depth of the borehole in feet, measured from ground surface to the
bottom of the exploration.

\subsection{GroundElevation}\label{groundelevation}

Ground elevation at the boring site, referenced to the vertical datum
specified in RestorFeatVdatum.

\subsection{Station}\label{station-5}

Stationing of the feature, typically referenced to a project baseline or
alignment (e.g., 10+50).

\subsection{sdsMetaDataID}\label{sdsmetadataid-18}

References attached metadata.

\subsection{RestorFeatIDPK}\label{restorfeatidpk-18}

Primary key used to uniquely identify the restoration feature record.

\subsection{sdsID}\label{sdsid-18}

Identifier used to link the feature to supporting SDS (spatial data
system) records or related datasets.

\subsection{mediaIDFK}\label{mediaidfk-18}

Foreign key used to reference associated media (e.g., boring logs,
photographs, lab reports, as-built drawings) documenting the borehole.

\chapter{Electrical Utility Node}\label{electrical-utility-node}

\section{Spatial \& Geometric
Definition}\label{spatial-geometric-definition-19}

\begin{itemize}
\item
  Geometry Type: Point (Electrical\_Utility\_Node).
\item
  Spatial Delineation: Point locations representing electrical
  infrastructure components powering restoration projects.
\item
  Component Types: Step-down transformers, primary utility metering
  poles, service disconnect panels, automatic transfer switches (ATS),
  backup diesel generators, motor control centers (MCC), and SCADA
  telemetry panels.
\end{itemize}

\subsection{Primary Engineering \& Operational
Objectives}\label{primary-engineering-operational-objectives-1}

\begin{itemize}
\item
  Power Supply \& Conditioning: Deliver and step down commercial
  high-voltage power (480V, 3-phase) to operational voltage for
  high-capacity axial pumps, gate actuators, and lighting.
\item
  Automated System Control: House Programmable Logic Controllers (PLCs)
  and variable-frequency drives (VFDs) governing automated pumping
  regimes based on water level sensors.
\item
  Emergency Redundancy: Provide automatic generator startup within 10
  seconds of commercial utility grid failure to protect facilities
  during flood events.
\end{itemize}

\subsection{Technical \& Electrical Engineering
Specifications}\label{technical-electrical-engineering-specifications}

\begin{itemize}
\item
  Enclosure Standards: NEMA 4X stainless steel or heavy marine-grade
  aluminum enclosures to resist river corridor corrosion and moisture.
\item
  Elevated Platforms: All electrical equipment, MCCs, and generators
  mounted on structural steel platforms elevated minimum 2.0 to 3.0 ft
  above the 100-year Base Flood Elevation (BFE).
\item
  Surge \& Lightning Protection: Industrial metal-oxide varistor (MOV)
  surge protection devices (SPDs) and copper grounding rod grids tied to
  all metallic structures.
\end{itemize}

\subsection{Construction \& Implementation
Methodology}\label{construction-implementation-methodology-18}

\begin{itemize}
\item
  Utility Coordination: Point-of-service coordination with local
  commercial power utilities, extending single-phase or three-phase
  overhead/underground feeders to project boundaries.
\item
  Commissioning: Full load-bank testing of transformers and standby
  generators, megger testing of insulation on all pump motor leads, and
  telemetry testing with central control stations.
\end{itemize}

\subsection{Quality Control, Inspection \& Long-Term
Maintenance}\label{quality-control-inspection-long-term-maintenance-18}

\begin{itemize}
\tightlist
\item
  Infrared Thermography Scans: Annual thermal imaging scans of
  electrical panels under full pumping loads to detect hot spots, loose
  connections, or failing breaker contacts.
\end{itemize}

\section{Spatial Feature Attributes}\label{spatial-feature-attributes-3}

\subsection{Definition}\label{definition-18}

An Electrical Utility Node is a point feature representing a subdivision
of an electrical distribution network, specifically an asset that
participates in the transmission, transformation, measurement, or
control of electricity but is not a line. Examples include transformers,
fuses, generators, meters, and switches.

\subsection{Purpose}\label{purpose-22}

Electrical utility nodes are used to document and manage key components
of the electrical distribution system that control, transform, or
monitor electrical power. Mapping these nodes supports design
coordination, safety and clearance reviews, maintenance planning, outage
management, and protection of electrical infrastructure during
restoration project construction and operations.

\subsection{Applicability}\label{applicability-19}

Electrical utility nodes are applicable wherever electrical distribution
infrastructure intersects, serves, or is affected by Upper Mississippi
River System (UMRS) restoration projects. Typical applications include
documenting power feeds and components at project facilities,
structures, pump stations, and control sites, as well as mapping utility
assets that may constrain or influence design and construction
activities. (UMRR Design Handbook, 2012)

\subsection{Notes}\label{notes-22}

Restoration features now have a Program\_UID field to distinguish
between NESP and UMRR features. Aliases without program identification
will be created for any field that currently has UMRR in the name, so
they can be used for NESP or UMRR data.

Attribute criteria for each field includes the following:

\begin{itemize}
\tightlist
\item
  Attribute Status -- defines whether the field requires a value or if
  it is optional based on available information.
\item
  Phases -- indicates the earliest time in the project lifecycle values
  for this attribute are available.
\item
  Digitizing requirements -- indicates when the field value should be
  entered in the attribute table.
\end{itemize}

Geometry Note: No Z values are needed for this feature.

\subsection{ElectricalUtilityNode}\label{electricalutilitynode}

(point) -- A subdivision of an electrical distribution network,
particularly an asset that participates in the transmission of
electricity but is not a line (e.g., transformer, fuse, generator,
meter, switch).

\begin{center}\rule{0.5\linewidth}{0.5pt}\end{center}

\section{Feature Attribute
Criteria}\label{feature-attribute-criteria-19}

\#

Attribute

Attribute Status

Phase

Digitizing Requirements

1

Program\_UID

Required

Feasibility

At feature creation

2

PROJECT\_ID (Alias)

Required

Feasibility

At feature creation

3

Type (Alias)

Required

Feasibility

At feature creation

4

Subtype (Alias)

Required

Feasibility

At feature creation

5

RestorFeatStatus

Required

Feasibility

At feature creation

6

sdsFeatureName

Optional

Design

As required and available

7

sdsFeatureDescription

Optional

Design

As required and available

8

RestorFeatVdatum

Optional

Design

As required and available

9

electricalSegmentType

Optional

Design

As required and available

10

OperationStatus

Optional

Design

As required and available

11

realPropertySiteUniqueID

Optional

Design

As required and available

12

realPropertyNetworkUniqueID

Optional

Design

As required and available

13

propertyCodeID

Optional

Construction

As required and available

14

sdsMetaDataID

Optional

Construction

As required and available

15

RestorFeatIDPK

Optional

Construction

As required and available

16

sdsID

Optional

Construction

As required and available

\begin{center}\rule{0.5\linewidth}{0.5pt}\end{center}

\section{Parameters}\label{parameters-19}

\subsection{Program\_UID}\label{program_uid-19}

Designate program NESP or UMRR.

\subsection{Project\_ID (alias:
PROJECT\_ID)}\label{project_id-alias-project_id-19}

Foreign key to the unique NESP/UMRR project ID.

\subsection{Type (alias:
Restoration\_Type)}\label{type-alias-restoration_type}

A broad classification used to group NESP/UMRR restoration features into
categories of similar function.

Value

Description

Restoration\_Type

Support feature.

\subsection{SubType (alias: Subtype)}\label{subtype-alias-subtype-19}

A classification used to group NESP/UMRR restoration features into
specific design feature categories.

Value

Description

ElectricalUtility

Electrical utility.

\subsection{RestorFeatStatus}\label{restorfeatstatus-19}

A description of the status of the individual restoration. It also
includes the party responsible for the individual feature.

Value

Description

USACE Planned

Feature is in a USACE planning stage.

USACE Built

Feature has been built by USACE. If a feature has been planned and built
by USACE, once constructed, will always be considered USACE built even
if it has been repaired under later project stages.

USACE Enhanced

The function of an existing feature that was not built by USACE has been
enhanced by a feature that was USACE built or repaired.

USACE Repaired

An existing feature that was not built by USACE has been directly
improved or repaired.

Other Planned

A feature that is planned by another organization or government agency.

Other Built

Feature has been built by another organization or government agency and
is an integral part of the UMRR project.

Other Enhanced

The function of an existing feature, that was built by another
organization or government agency, has been enhanced by a feature that
was built or repaired by a non-USACE organization or government agency.

Other Repaired

An existing feature that has been directly improved or repaired by
another organization or government agency.

\subsection{sdsFeatureName}\label{sdsfeaturename-19}

The common name of the feature.

\subsection{sdsFeatureDescription}\label{sdsfeaturedescription-19}

A narrative describing the feature, including its electrical function
and relationship to nearby structures or facilities.

\subsection{RestorFeatVdatum}\label{restorfeatvdatum-19}

Vertical datum of any elevations recorded in the attribute table (if
used in related contexts).

Value

Description

MSL 1912

Mean Sea Level of 1912.

NGVD 1929

National Geodetic Vertical Datum of 1929.

NAVD 1988

North American Vertical Datum of 1988.

Other

Other; note datum in sdsFeatureDescription.

Unknown

Unknown vertical datum.

\subsection{electricalSegmentType}\label{electricalsegmenttype}

Classification of the electrical utility node within the distribution
network (based on configuration or function).

Value

Description

Transformer

Device used to step voltage up or down in the distribution network.

Switch / Breaker

Switching or circuit-breaking device used to isolate or reroute
electrical circuits.

Fuse / Protective Device

Fuse or other protective device used to interrupt fault currents or
protect equipment.

Meter

Metering device used to measure electrical usage or load.

Generator

Local generation unit (e.g., backup generator) connected to the
distribution network.

Junction / Cabinet

Junction box, cabinet, or similar node where multiple circuits are
connected.

Other

Electrical node type not covered above; describe in
sdsFeatureDescription.

Unknown

Node type has not been determined.

\subsection{OperationStatus}\label{operationstatus}

Operational status of the electrical utility node.

Value

Description

Planned

Node is proposed but not yet installed or energized.

Active

Node is installed and in normal service.

De-energized

Node is installed but not currently energized.

Out of Service

Node is no longer used for regular electrical service but remains in
place.

Removed

Node has been physically removed or abandoned per utility practice.

Unknown

Operational status has not been determined.

\subsection{realPropertySiteUniqueID}\label{realpropertysiteuniqueid}

Unique identifier linking the electrical utility node to a specific real
property site or installation record in the real property inventory,
when applicable.

\subsection{realPropertyNetworkUniqueID}\label{realpropertynetworkuniqueid}

Unique identifier linking the electrical utility node to a network-level
real property or asset management record (e.g., regional or enterprise
utility asset system), when applicable.

\subsection{propertyCodeID}\label{propertycodeid}

Code indicating the real property classification or category of the
asset (e.g., electrical equipment, transformer, switchgear) as defined
by the managing agency's real property coding system.

\subsection{sdsMetaDataID}\label{sdsmetadataid-19}

References attached metadata.

\subsection{RestorFeatIDPK}\label{restorfeatidpk-19}

Primary key used to uniquely identify the restoration feature record.

\subsection{sdsID}\label{sdsid-19}

Identifier used to link the feature to supporting SDS (spatial data
system) records or related datasets.

\subsection{mediaIDFK}\label{mediaidfk-19}

Foreign key used to reference associated media (e.g., photographs,
utility plans, manufacturer cut sheets, as-built drawings) documenting
the electrical utility node.

\chapter{Electrical Utility Segment}\label{electrical-utility-segment}

\section{Spatial \& Geometric
Definition}\label{spatial-geometric-definition-20}

\begin{itemize}
\item
  Geometry Type: Polyline (Electrical\_Utility\_Segment).
\item
  Spatial Delineation: Linear paths of underground, conduit-encased,
  aerial, or submarine electrical power distribution lines linking
  transformers to pumps and control nodes.
\item
  Structural Types: Direct-buried armored cables, concrete-encased duct
  banks, rigid galvanized steel (RGS) conduits, and aerial conductor
  spans.
\end{itemize}

\subsection{Primary Engineering \& Operational
Objectives}\label{primary-engineering-operational-objectives-2}

\begin{itemize}
\item
  Safe Power Distribution: Route electrical power across remote river
  floodplains, levee crests, and wetland environments without exposure
  to environmental or mechanical damage.
\item
  Subsurface Protection: Safeguard high-voltage feeder cables from
  burrowing rodents, flooding, seasonal frost heave, and maintenance
  grading.
\end{itemize}

\subsection{Technical Specifications \& Burial
Depths}\label{technical-specifications-burial-depths}

\begin{itemize}
\item
  Underground Feeders: Copper conductors (THHN/THWN−2 or XHHW−2) sized
  to limit voltage drop to \textless3\% at maximum continuous operating
  current.
\item
  Burial Depth Standards: Minimum burial depth of 36 to 48 in below
  finished grade.
\item
  Conduit Specification: Schedule 80 heavy-wall PVC conduits or rigid
  galvanized steel (RGS) where traversing roadways or embankment
  crossings.
\item
  Detectable Warning Ribbon: High-visibility red magnetic warning tape
  (``Caution: Buried Electric Line Below'') placed 12 in above conduits
  during trench backfill.
\end{itemize}

\subsection{Construction \& Implementation
Methodology}\label{construction-implementation-methodology-19}

\begin{itemize}
\tightlist
\item
  Trenching \& Horizontal Directional Drilling (HDD): Installed via
  open-cut trenching backfilled with compacted sand bedding, or HDD
  under existing levees and navigation channels to avoid disturbing
  structural embankments.
\end{itemize}

\subsection{Quality Control, Inspection \& Long-Term
Maintenance}\label{quality-control-inspection-long-term-maintenance-19}

\begin{itemize}
\item
  As-Built Traceability: All linear conduits and pull boxes mapped using
  sub-meter GPS prior to trench backfill.
\item
  Insulation Resistance Testing: Periodic high-potential (Hi-Pot)
  insulation testing to detect degradation of cable jackets exposed to
  high water tables.
\end{itemize}

\section{Spatial Feature Definition}\label{spatial-feature-definition}

\subsection{Definition}\label{definition-19}

A subdivision of an electrical distribution network, represented as a
linear feature that carries electrical power between structures such as
poles, towers, transformers, or service points.

\subsection{Purpose}\label{purpose-23}

Electrical utility segments are used to document and manage individual
runs of electrical distribution or transmission lines within a project.
Segmenting the electrical network supports design, coordination with
utilities, clearance and safety reviews, maintenance planning, and
impact assessment for restoration features and project construction
activities.

\subsection{Applicability}\label{applicability-20}

Electrical utility segments are applicable wherever overhead or
underground electrical lines intersect, parallel, or serve Upper
Mississippi River System (UMRS) restoration project areas. Typical uses
include mapping power feeds to project facilities, documenting utility
crossings of levees or channels, and coordinating with local utilities
during design and construction. (UMRR Design Handbook, 2012)

\subsection{Notes}\label{notes-23}

Restoration features now have a Program\_UID field to distinguish
between NESP and UMRR features. Aliases without program identification
will be created for any field that currently has UMRR in the name, so
they can be used for NESP or UMRR data.

Attribute criteria for each field includes the following:

\begin{itemize}
\tightlist
\item
  Attribute Status -- defines whether the field requires a value or if
  it is optional based on available information.
\item
  Phases -- indicates the earliest time in the project lifecycle values
  for this attribute are available.
\item
  Digitizing requirements -- indicates when the field value should be
  entered in the attribute table.
\end{itemize}

Geometry Note: No Z values are needed for this feature.

\subsection{ElectricalUtilitySegment}\label{electricalutilitysegment}

(line) -- A subdivision of an electrical distribution network,
particularly a line for the transmission of electricity.

\begin{center}\rule{0.5\linewidth}{0.5pt}\end{center}

\section{Feature Attribute
Criteria}\label{feature-attribute-criteria-20}

\#

Attribute

Attribute Status

Phase

Digitizing Requirements

1

Program\_UID

Required

Feasibility

At feature creation

2

PROJECT\_ID (Alias)

Required

Feasibility

At feature creation

3

Type (Alias)

Required

Feasibility

At feature creation

4

Subtype (Alias)

Required

Feasibility

At feature creation

5

RestorFeatStatus

Required

Feasibility

At feature creation

6

sdsFeatureName

Optional

Design

As required and available

7

sdsFeatureDescription

Optional

Design

As required and available

8

RestorFeatVdatum

Optional

Design

As required and available

9

electricalSegmentType

Optional

Design

As required and available

10

OperationStatus

Optional

Design

As required and available

11

realPropertySiteUniqueID

Optional

Design

As required and available

12

realPropertyNetworkUniqueID

Optional

Design

As required and available

13

propertyCodeID

Optional

Construction

As required and available

14

sdsMetaDataID

Optional

Construction

As required and available

15

RestorFeatIDPK

Optional

Construction

As required and available

16

sdsID

Optional

Construction

As required and available

\begin{center}\rule{0.5\linewidth}{0.5pt}\end{center}

\section{Parameters}\label{parameters-20}

\subsection{Program\_UID}\label{program_uid-20}

Designate program NESP or UMRR.

\subsection{Project\_ID (alias:
PROJECT\_ID)}\label{project_id-alias-project_id-20}

Foreign key to the unique NESP/UMRR project ID.

\subsection{Type (alias:
Restoration\_Type)}\label{type-alias-restoration_type-1}

A broad classification used to group NESP/UMRR restoration features into
categories of similar function.

Value

Description

Restoration\_Type

Support feature

\subsection{SubType (alias: Subtype)}\label{subtype-alias-subtype-20}

A classification used to group NESP/UMRR restoration features into
specific design feature categories.

Value

Description

ElectricalUtility

Electrical utility

\subsection{RestorFeatStatus}\label{restorfeatstatus-20}

A description of the status of the individual restoration. It also
includes the party responsible for the individual feature.

Value

Description

USACE Planned

Feature is in a USACE planning stage.

USACE Built

Feature has been built by USACE. If a feature has been planned and built
by USACE, once constructed, will always be considered USACE built even
if it has been repaired under later project stages.

USACE Enhanced

The function of an existing feature that was not built by USACE has been
enhanced by a feature that was USACE built or repaired.

USACE Repaired

An existing feature that was not built by USACE has been directly
improved or repaired.

Other Planned

A feature that is planned by another organization or government agency.

Other Built

Feature has been built by another organization or government agency and
is an integral part of the UMRR project.

Other Enhanced

The function of an existing feature, that was built by another
organization or government agency, has been enhanced by a feature that
was built or repaired by a non-USACE organization or government agency.

Other Repaired

An existing feature that has been directly improved or repaired by
another organization or government agency.

\subsection{sdsFeatureName}\label{sdsfeaturename-20}

The common name of the feature.

\subsection{sdsFeatureDescription}\label{sdsfeaturedescription-20}

A narrative describing the feature.

\subsection{RestorFeatVdatum}\label{restorfeatvdatum-20}

Vertical datum of elevations recorded in the attribute table or used as
feature z values.

Value

Description

MSL 1912

Mean Sea Level of 1912

NGVD 1929

National Geodetic Vertical Datum of 1929

NAVD 1988

North American Vertical Datum of 1988

Other

Other, note datum in sdsFeatureDescription

Unknown

Unknown vertical datum

\subsection{electricalSegmentType}\label{electricalsegmenttype-1}

Classification of the electrical utility segment based on configuration,
voltage class, or installation type.

Value

Description

Overhead Distribution

Overhead electrical line used for local distribution.

Overhead Transmission

Overhead electrical line used for higher-voltage transmission.

Underground Distribution

Underground electrical line used for local distribution.

Service Drop

Segment providing service connection from main line to a facility or
structure.

Street / Area Lighting Feed

Electrical segment dedicated primarily to lighting circuits.

Other

Electrical segment type not covered above; describe in
sdsFeatureDescription.

Unknown

Segment type has not been determined.

\subsection{OperationStatus}\label{operationstatus-1}

Operational status of the electrical utility segment.

Value

Description

Planned

Segment is proposed but not yet constructed or energized.

Active

Segment is installed and energized for normal service.

De-energized

Segment is installed but not currently energized.

Out of Service

Segment is no longer used for regular electrical service but remains in
place.

Removed

Segment has been physically removed or abandoned in place per utility
practice.

Unknown

Operational status has not been determined.

\subsection{realPropertySiteUniqueID}\label{realpropertysiteuniqueid-1}

Unique identifier linking the electrical utility segment to a specific
real property site or installation record in the real property
inventory, when applicable.

\subsection{realPropertyNetworkUniqueID}\label{realpropertynetworkuniqueid-1}

Unique identifier linking the electrical utility segment to a
network-level real property or asset management record (e.g., a regional
or enterprise real property system), when applicable.

\subsection{propertyCodeID}\label{propertycodeid-1}

Code indicating the real property classification or category of the
asset (e.g., utility line, electrical distribution) as defined by the
managing agency's real property coding system.

\subsection{sdsMetaDataID}\label{sdsmetadataid-20}

References attached metadata.

\subsection{RestorFeatIDPK}\label{restorfeatidpk-20}

Primary key used to uniquely identify the restoration feature record.

\subsection{sdsID}\label{sdsid-20}

Identifier used to link the feature to supporting SDS (spatial data
system) records or related datasets.

\subsection{mediaIDFK}\label{mediaidfk-20}

Foreign key used to reference associated media (e.g., photographs,
utility plans, as-built drawings) documenting the electrical utility
segment.

\chapter{Feature Stationing}\label{feature-stationing}

\section{Spatial \& Geometric
Definition}\label{spatial-geometric-definition-21}

\begin{itemize}
\item
  Geometry Type: Point (Feature\_Stationing\_Point).
\item
  Spatial Delineation: Regular linear control points established along
  the alignment of an embankment centerline, access roadway, channel
  cut, or stone dike.
\item
  Format: Standard civil engineering linear stationing markers (0+00,
  1+00, 5+00, etc.), where each station represents 100.0 linear feet.
\end{itemize}

\subsection{Primary Engineering \& Operational
Objectives}\label{primary-engineering-operational-objectives-3}

\begin{itemize}
\item
  Linear Coordinate Referencing: Establishes the authoritative linear
  measurement system used across all contract plan sheets,
  cross-sections, and inspection reports.
\item
  Asset Tracking \& Work Measurement: Links geotechnical borings,
  culverts, spillway notches, and riprap armor transitions to an
  unambiguous linear coordinate along the structure.
\item
  Construction Progress Verification: Enables contract field inspectors
  to precisely verify contractor pay-application quantities based on
  physical station intervals.
\end{itemize}

\subsection{Typical Specifications \& Geometric
Properties}\label{typical-specifications-geometric-properties-2}

\begin{itemize}
\item
  Station Intervals: Points established every 100 ft on tangents and
  every 50 ft or 25 ft on horizontal curves.
\item
  Attribute Domains: Alignment Name, Station Number (XX+YY), Design
  Crest Grade, Left/Right Slope Offsets, and Elevation.
\end{itemize}

\subsection{Survey \& Layout
Implementation}\label{survey-layout-implementation}

\begin{itemize}
\tightlist
\item
  Staking Standards: Survey field crews install marked wooden hubs, iron
  pins, or aluminum caps during construction, tied to the project
  horizontal baseline.
\end{itemize}

\subsection{Quality Control, Inspection \& Long-Term
Maintenance}\label{quality-control-inspection-long-term-maintenance-20}

\begin{itemize}
\tightlist
\item
  Permanent Reference Hubs: Key stationing monuments (every 1,000 ft)
  cast into concrete markers or stamped brass disks to serve as
  permanent physical recovery points for post-flood inspection teams.
\end{itemize}

\section{Spatial Feature Attributes}\label{spatial-feature-attributes-4}

\subsection{Definition}\label{definition-20}

Feature Stationing: This refers to a system of linear referencing used
to locate features along an alignment, such as an embankment or road. It
is used for design, construction, and maintenance purposes. Reference
points, generally expressed in linear units (typically feet), that allow
identification of specific locations on engineered structures such as
levees, control structures, roads, or walls.

\subsection{Purpose}\label{purpose-24}

Feature stationing is used to tie design documents, construction
records, inspections, and maintenance activities to precise locations
along an engineered structure. By assigning station codes and associated
elevations, project teams can consistently reference, locate, and
document conditions, work items, or features along the structure.

\subsection{Applicability}\label{applicability-21}

Feature stationing is applicable to any linear engineered feature within
Upper Mississippi River System (UMRS) restoration projects where
precise, repeatable location referencing is needed. Common examples
include levees, floodwalls, spillways, roads, channels, and other
structures where plans and profiles are based on stationing. (UMRR
Design Handbook, 2012)

\subsection{Notes}\label{notes-24}

Restoration features now have a Program\_UID field to distinguish
between NESP and UMRR features. Aliases without program identification
will be created for any field that currently has UMRR in the name, so
they can be used for NESP or UMRR data.

Attribute criteria for each field includes the following:

\begin{itemize}
\tightlist
\item
  Attribute Status -- defines whether the field requires a value or if
  it is optional based on available information.
\item
  Phases -- indicates the earliest time in the project lifecycle values
  for this attribute are available.
\item
  Digitizing requirements -- indicates when the field value should be
  entered in the attribute table.
\end{itemize}

Geometry Note: No Z values are needed for this feature.

\section{FeatureStationing}\label{featurestationing}

(point) -- Reference points, generally in feet, that allow
identification of specific locations on engineered structures.

\begin{center}\rule{0.5\linewidth}{0.5pt}\end{center}

\section{Feature Attribute
Criteria}\label{feature-attribute-criteria-21}

\#

Attribute

Attribute Status

Phase

Digitizing Requirements

1

Program\_UID

Required

Feasibility

At feature creation

2

PROJECT\_ID (Alias)

Required

Feasibility

At feature creation

3

Type (Alias)

Required

Feasibility

At feature creation

4

Subtype (Alias)

Required

Feasibility

At feature creation

5

RestorFeatStatus

Required

Feasibility

At feature creation

6

sdsFeatureName

Optional

Design

As required and available

7

sdsFeatureDescription

Optional

Design

As required and available

8

RestorFeatVdatum

Optional

Design

As required and available

9

StationCode

Optional

Design

As required and available

10

StationElevation

Optional

Design

As required and available

11

sdsMetaDataID

Optional

Construction

As required and available

12

RestorFeatIDPK

Optional

Construction

As required and available

13

sdsID

Optional

Construction

As required and available

14

mediaIDFK

Optional

Construction

As required and available

\begin{center}\rule{0.5\linewidth}{0.5pt}\end{center}

\section{Parameters}\label{parameters-21}

\subsection{Program\_UID}\label{program_uid-21}

Designate program NESP or UMRR.

\subsection{Project\_ID (alias:
PROJECT\_ID)}\label{project_id-alias-project_id-21}

Foreign key to the unique NESP/UMRR project ID.

\subsection{Type (alias:
Restoration\_Type)}\label{type-alias-restoration_type-2}

A broad classification used to group NESP/UMRR restoration features into
categories of similar function.

Value

Description

Restoration\_Type

Support feature

\subsection{SubType (alias: Subtype)}\label{subtype-alias-subtype-21}

A classification used to group NESP/UMRR restoration features into
specific design feature categories.

Value

Description

Feature Stationing

Stationing reference points on engineered features.

\subsection{RestorFeatStatus}\label{restorfeatstatus-21}

A description of the status of the individual restoration. It also
includes the party responsible for the individual feature.

Value

Description

USACE Planned

Feature is in a USACE planning stage.

USACE Built

Feature has been built by USACE. If a feature has been planned and built
by USACE, once constructed, it will always be considered USACE built
even if it has been repaired under later project stages.

USACE Enhanced

The function of an existing feature that was not built by USACE has been
enhanced by a feature that was USACE built or repaired.

USACE Repaired

An existing feature that was not built by USACE has been directly
improved or repaired.

Other Planned

A feature that is planned by another organization or government agency.

Other Built

Feature has been built by another organization or government agency and
is an integral part of the UMRR project.

Other Enhanced

The function of an existing feature, that was built by another
organization or government agency, has been enhanced by a feature that
was built or repaired by a non-USACE organization or government agency.

Other Repaired

An existing feature that has been directly improved or repaired by
another organization or government agency.

\subsection{sdsFeatureName}\label{sdsfeaturename-21}

The common name of the feature.

\subsection{sdsFeatureDescription}\label{sdsfeaturedescription-21}

A narrative describing the feature.

\subsection{RestorFeatVdatum}\label{restorfeatvdatum-21}

Vertical datum of elevations recorded in the attribute table or used as
feature z values.

Value

Description

MSL 1912

Mean Sea Level of 1912

NGVD 1929

National Geodetic Vertical Datum of 1929

NAVD 1988

North American Vertical Datum of 1988

Other

Other, note datum in sdsFeatureDescription

Unknown

Unknown vertical datum

\subsection{StationCode}\label{stationcode}

Linear reference identifier for the station point, typically recorded in
engineering station format (e.g., 10+00), representing distance along a
baseline or alignment.

\subsection{StationElevation}\label{stationelevation}

Elevation at the station point, referenced to the specified vertical
datum in RestorFeatVdatum.

\subsection{sdsMetaDataID}\label{sdsmetadataid-21}

References attached metadata.

\subsection{RestorFeatIDPK}\label{restorfeatidpk-21}

Primary key used to uniquely identify the restoration feature record.

\subsection{sdsID}\label{sdsid-21}

Identifier used to link the feature to supporting SDS (spatial data
system) records or related datasets.

\subsection{mediaIDFK}\label{mediaidfk-21}

Foreign key used to reference associated media (e.g., photographs,
drawings, as-built plans) documenting the station location.

\chapter{Pavement Section}\label{pavement-section}

\section{Spatial \& Geometric
Definition}\label{spatial-geometric-definition-22}

\begin{itemize}
\item
  Geometry Type: Polygon (Pavement\_Section\_Footprint).
\item
  Spatial Delineation: Surface areas receiving engineered structural
  wearing courses across pump station staging yards, boat ramps, crest
  access roads, and parking turnarounds.
\item
  Surfacing Typology: Hot-mix asphalt (HMA), cast-in-place reinforced
  concrete pavement, and compacted crushed limestone aggregate
  surfacing.
\end{itemize}

\subsection{Primary Engineering \& Civil
Objectives}\label{primary-engineering-civil-objectives}

\begin{itemize}
\item
  Load-Bearing Surface: Support heavy maintenance vehicle loads (e.g.,
  HS-20 truck loading: crane boom trucks, fuel tankers, and
  tractor-trailers) without rutting underlying subgrades.
\item
  Erosion \& Overtopping Resistance: Prevent surface shear erosion on
  embankment crests and ramp approaches during controlled overtopping
  flood events.
\item
  All-Weather Operational Access: Ensure round-the-clock maintenance
  access to remote water control structures during heavy rain and
  flood-fighting operations.
\end{itemize}

\subsection{Typical Structural Pavement
Cross-Sections}\label{typical-structural-pavement-cross-sections}

\begin{itemize}
\item
  Heavy-Duty Concrete (Pump Yards / Ramps): 8.0 to 10.0 in cast-in-place
  concrete (fc′ ≥4,000 psi) with epoxy-coated rebar mats, underlain by
  6.0 in of crushed aggregate base over non-woven geotextile.
\item
  Crushed Aggregate Surfacing (Levee Roads): 6.0 to 8.0 in dense-graded
  crushed limestone road base (IDOT/WDOT dense-graded aggregate),
  compacted to ≥98\% Standard Proctor density.
\item
  Cross-Slope Drainage: Crowned centerline or uniform cross-slope of
  2.0\% directing surface sheet flow into vegetated side swales.
\end{itemize}

\subsection{Construction \& Implementation
Methodology}\label{construction-implementation-methodology-20}

\begin{itemize}
\item
  Subgrade Stabilization: Proof-rolling native subgrades; subgrade
  stabilization using biaxial geogrids where crossing soft alluvial clay
  strata.
\item
  Joint Sealants \& Finishing: Concrete slabs scored with transverse
  contraction joints every 15 ft; silicone joint sealants applied to
  exclude water and freeze-thaw cracking.
\end{itemize}

\subsection{Quality Control, Inspection \& Long-Term
Maintenance}\label{quality-control-inspection-long-term-maintenance-21}

\begin{itemize}
\tightlist
\item
  Pavement Condition Audits: Annual inspections checking for surface
  rutting, washboarding, joint spalling, or subgrade settlement
  following flood inundation.
\end{itemize}

\section{Spatial Feature Attributes}\label{spatial-feature-attributes-5}

\subsection{Definition}\label{definition-21}

A pavement section is a portion of a pavement branch that differs in
some aspect from other sections such that further segmentation is
required to uniquely identify that section. For example, if the only
difference in a road is the pavement material type -- concrete for part
of the road and asphalt for the remainder of the road -- that road
branch has two sections.

\subsection{Purpose}\label{purpose-25}

Pavement sections are defined to document and manage areas of pavement
that have distinct materials, structural composition, age, or
performance characteristics. Segmenting pavement into sections supports
design, maintenance planning, inspection, rehabilitation, and lifecycle
management of project roadways and paved areas.

\subsection{Applicability}\label{applicability-22}

Pavement sections are applicable wherever paved surfaces occur within
Upper Mississippi River System (UMRS) restoration projects, including
access roads, parking areas, staging areas, and operational surfaces.
They are used when differences in pavement type, thickness, construction
date, or condition require unique identification for engineering or
asset management purposes. (UMRR Design Handbook, 2012)

\subsection{Notes}\label{notes-25}

Restoration features now have a Program\_UID field to distinguish
between NESP and UMRR features. Aliases without program identification
will be created for any field that currently has UMRR in the name, so
they can be used for NESP or UMRR data.

Attribute criteria for each field includes the following:

\begin{itemize}
\tightlist
\item
  Attribute Status -- defines whether the field requires a value or if
  it is optional based on available information.
\item
  Phases -- indicates the earliest time in the project lifecycle values
  for this attribute are available.
\item
  Digitizing requirements -- indicates when the field value should be
  entered in the attribute table.
\end{itemize}

Geometry Note: No Z values are needed for this feature.

\subsection{PavementSection}\label{pavementsection}

(polygon) -- A pavement section is a portion of a pavement branch that
differs in some aspect from other sections such that further
segmentation is required to uniquely identify that section.

\begin{center}\rule{0.5\linewidth}{0.5pt}\end{center}

\section{Feature Attribute
Criteria}\label{feature-attribute-criteria-22}

\#

Attribute

Attribute Status

Phase

Digitizing Requirements

1

Program\_UID

Required

Feasibility

At feature creation

2

PROJECT\_ID (Alias)

Required

Feasibility

At feature creation

3

Type (Alias)

Required

Feasibility

At feature creation

4

Subtype (Alias)

Required

Feasibility

At feature creation

5

RestorFeatStatus

Required

Feasibility

At feature creation

6

sdsFeatureName

Optional

Design

As required and available

7

sdsFeatureDescription

Optional

Design

As required and available

8

RestorFeatVdatum

Optional

Design

As required and available

9

PavementSectionType

Optional

Design

As required and available

10

OperationStatus

Optional

Design

As required and available

11

realPropertySiteUniqueID

Optional

Design

As required and available

12

realPropertyNetworkUniqueID

Optional

Design

As required and available

13

propertyCodeID

Optional

Construction

As required and available

14

sdsMetaDataID

Optional

Construction

As required and available

15

RestorFeatIDPK

Optional

Construction

As required and available

16

sdsID

Optional

Construction

As required and available

\begin{center}\rule{0.5\linewidth}{0.5pt}\end{center}

\section{Parameters}\label{parameters-22}

\subsection{Program\_UID}\label{program_uid-22}

Designate program NESP or UMRR.

\subsection{Project\_ID (alias:
PROJECT\_ID)}\label{project_id-alias-project_id-22}

Foreign key to the unique NESP/UMRR project ID.

\subsection{Type (alias:
Restoration\_Type)}\label{type-alias-restoration_type-3}

A broad classification used to group NESP/UMRR restoration features into
categories of similar function.

Value

Description

Restoration\_Type

Support feature

\subsection{SubType (alias: Subtype)}\label{subtype-alias-subtype-22}

A classification used to group NESP/UMRR restoration features into
specific design feature categories.

Value

Description

Road Utility

Road utility

\subsection{RestorFeatStatus}\label{restorfeatstatus-22}

A description of the status of the individual restoration. It also
includes the party responsible for the individual feature.

Value

Description

USACE Planned

Feature is in a USACE planning stage.

USACE Built

Feature has been built by USACE. If a feature has been planned and built
by USACE, once constructed, it will always be considered USACE built
even if it has been repaired under later project stages.

USACE Enhanced

The function of an existing feature that was not built by USACE has been
enhanced by a feature that was USACE built or repaired.

USACE Repaired

An existing feature that was not built by USACE has been directly
improved or repaired.

Other Planned

A feature that is planned by another organization or government agency.

Other Built

Feature has been built by another organization or government agency and
is an integral part of the UMRR project.

Other Enhanced

The function of an existing feature, that was built by another
organization or government agency, has been enhanced by a feature that
was built or repaired by a non-USACE organization or government agency.

Other Repaired

An existing feature that has been directly improved or repaired by
another organization or government agency.

\subsection{sdsFeatureName}\label{sdsfeaturename-22}

The common name of the feature.

\subsection{sdsFeatureDescription}\label{sdsfeaturedescription-22}

A narrative describing the feature.

\subsection{RestorFeatVdatum}\label{restorfeatvdatum-22}

Vertical datum of elevations recorded in attribute table or used as
feature z values.

Value

Description

MSL 1912

Mean Sea Level of 1912

NGVD 1929

National Geodetic Vertical Datum of 1929

NAVD 1988

North American Vertical Datum of 1988

Other

Other, note datum in sdsFeatureDescription

Unknown

Unknown vertical datum

\subsection{PavementSectionType}\label{pavementsectiontype}

Classification of the pavement section based on its structural or
material characteristics.

Value

Description

Rigid Concrete

Pavement section constructed with portland cement concrete.

Flexible Asphalt

Pavement section constructed with asphalt concrete or bituminous
surface.

Composite

Pavement section consisting of a combination of concrete and asphalt
layers.

Aggregate / Gravel

Pavement section with unbound aggregate or gravel surface.

Stabilized Soil

Pavement section with chemically or mechanically stabilized soil
surface.

Other

Pavement section type not covered above; describe in
sdsFeatureDescription.

Unknown

Pavement section type has not been determined.

\subsection{OperationStatus}\label{operationstatus-2}

Operational status of the pavement section.

Value

Description

Planned

Pavement section is proposed but not yet constructed.

Active

Pavement section is constructed and in regular use.

Limited Use

Pavement section is in place but subject to seasonal, restricted, or
intermittent use.

Out of Service

Pavement section is no longer used for regular traffic but remains in
place.

Removed

Pavement section has been removed or reclaimed.

Unknown

Operational status has not been determined.

\subsection{realPropertySiteUniqueID}\label{realpropertysiteuniqueid-2}

Unique identifier linking the pavement section to a specific real
property site or installation record in the real property inventory,
when applicable.

\subsection{realPropertyNetworkUniqueID}\label{realpropertynetworkuniqueid-2}

Unique identifier linking the pavement section to a network-level real
property or asset management record (e.g., a regional or enterprise real
property system), when applicable.

\subsection{propertyCodeID}\label{propertycodeid-2}

Code indicating the real property classification or category of the
asset (e.g., paved road, parking area, apron) as defined by the managing
agency's real property coding system.

\subsection{sdsMetaDataID}\label{sdsmetadataid-22}

References attached metadata.

\subsection{RestorFeatIDPK}\label{restorfeatidpk-22}

Primary key used to uniquely identify the restoration feature record.

\subsection{sdsID}\label{sdsid-22}

Identifier used to link the feature to supporting SDS (spatial data
system) records or related datasets.

\chapter{Road Segment}\label{road-segment}

\section{Spatial \& Geometric
Definition}\label{spatial-geometric-definition-23}

\begin{itemize}
\item
  Geometry Type: Polyline (Road\_Segment\_Centerline).
\item
  Spatial Delineation: Linear centerlines of service corridors, levee
  crest roads, haul paths, and permanent public access roadways
  servicing restoration sites.
\item
  Data Fields: Road Name, Road Class (Maintenance Service vs.~Public
  Access), Surface Material, Top Width, Design Speed, and Maximum Axle
  Load Rating.
\end{itemize}

\subsection{Primary Engineering \& Operational
Objectives}\label{primary-engineering-operational-objectives-4}

\begin{itemize}
\item
  Civil Network Connectivity: Form the continuous vehicular
  transportation network allowing USACE, USFWS, and State DNR personnel
  to access water control structures, pump stations, and habitat
  monitoring reaches.
\item
  Emergency Flood-Fighting Corridor: Serve as heavy haul routes for rock
  trucks and sandbagging operations during high-water river emergencies.
\end{itemize}

\subsection{Typical Cross-Sections, Dimensions \& Sizing
Criteria}\label{typical-cross-sections-dimensions-sizing-criteria-12}

\begin{itemize}
\item
  Roadway Widths: Crest service roads span 10.0 to 14.0 ft; two-way
  access roads span 20.0 to 24.0 ft.
\item
  Turnout Geometry: Passing turnouts (width20 ft, length 50 to 100 ft)
  established every 1,000 to 1,500 ft along single-lane levee roads.
\item
  Maximum Longitudinal Gradients: Grades held below 6\% (approaches to
  levee crests restricted to \textless8\%).
\end{itemize}

\subsection{Construction \& Implementation
Methodology}\label{construction-implementation-methodology-21}

\begin{itemize}
\item
  Grading \& Compaction: Motor graders establishing crowned profiles;
  heavy vibratory rollers sealing surface aggregates to shed rainfall
  without softening subgrades.
\item
  Culvert Crossings: Corrugated metal pipes (CMP) or dual-wall HDPE
  culverts (18 to 36 in) installed at low points with riprap flared end
  sections to ensure ditch conveyance.
\end{itemize}

\subsection{Quality Control, Inspection \& Long-Term
Maintenance}\label{quality-control-inspection-long-term-maintenance-22}

\begin{itemize}
\tightlist
\item
  Surface Regrading \& Dust Abatement: Annual post-winter grading, spot
  aggregate replenishment, and application of calcium chloride dust
  palliatives near public trail intersections.
\end{itemize}

\section{Spatial Feature Attributes}\label{spatial-feature-attributes-6}

\subsection{Definition}\label{definition-22}

A road segment is a portion of a road with similar characteristics such
as number of lanes, surface type, or functional use, represented as a
continuous linear feature. Road Segment: A road segment represents a
section of a road or access route within a project area.

\subsection{Purpose}\label{purpose-26}

Road segments are used to document and manage distinct portions of
project roads with consistent physical or operational characteristics.
Segmenting roads in this way supports design, construction, maintenance,
access management, and analysis of transportation and operational access
within restoration projects.

\subsection{Applicability}\label{applicability-23}

Road segments are applicable wherever project-related roads or access
routes occur within Upper Mississippi River System (UMRS) restoration
areas. They are typically used to represent project access roads,
levee-top roads, service roads, and internal circulation routes that
provide construction, operations, maintenance, or public access. (UMRR
Design Handbook, 2012)

\subsection{Notes}\label{notes-26}

Restoration features now have a Program\_UID field to distinguish
between NESP and UMRR features. Aliases without program identification
will be created for any field that currently has UMRR in the name, so
they can be used for NESP or UMRR data.

Attribute criteria for each field includes the following:

\begin{itemize}
\tightlist
\item
  Attribute Status -- defines whether the field requires a value or if
  it is optional based on available information.
\item
  Phases -- indicates the earliest time in the project lifecycle values
  for this attribute are available.
\item
  Digitizing requirements -- indicates when the field value should be
  entered in the attribute table.
\end{itemize}

Geometry Note: No Z values are needed for this feature.

\subsection{RoadSegment}\label{roadsegment}

(line) -- A portion of a road with similar characteristics such as the
number of lanes or surface type.

\begin{center}\rule{0.5\linewidth}{0.5pt}\end{center}

\section{Feature Attribute
Criteria}\label{feature-attribute-criteria-23}

\#

Attribute

Attribute Status

Phase

Digitizing Requirements

1

Program\_UID

Required

Feasibility

At feature creation

2

PROJECT\_ID (Alias)

Required

Feasibility

At feature creation

3

Type (Alias)

Required

Feasibility

At feature creation

4

Subtype (Alias)

Required

Feasibility

At feature creation

5

RestorFeatStatus

Required

Feasibility

At feature creation

6

sdsFeatureName

Optional

Design

As required and available

7

sdsFeatureDescription

Optional

Design

As required and available

8

RestorFeatVdatum

Optional

Design

As required and available

9

fieldMeasure

Optional

Design

As required and available

10

isAnchorSection

Optional

Design

As required and available

11

segmentLength

Optional

Design

As required and available

12

segmentLengthUOM

Optional

Design

As required and available

13

roadSegStatus

Optional

Construction

As required and available

14

sdsMetaDataID

Optional

Construction

As required and available

15

RestorFeatIDPK

Optional

Construction

As required and available

16

sdsID

Optional

Construction

As required and available

\begin{center}\rule{0.5\linewidth}{0.5pt}\end{center}

\section{Parameters}\label{parameters-23}

\subsection{Program\_UID}\label{program_uid-23}

Designate program NESP or UMRR.

\subsection{Project\_ID (alias:
PROJECT\_ID)}\label{project_id-alias-project_id-23}

Foreign key to the unique NESP/UMRR project ID.

\subsection{Type (alias:
Restoration\_Type)}\label{type-alias-restoration_type-4}

A broad classification used to group NESP/UMRR restoration features into
categories of similar function.

Value

Description

Restoration\_Type

Support feature

\subsection{SubType (alias: Subtype)}\label{subtype-alias-subtype-23}

A classification used to group NESP/UMRR restoration features into
specific design feature categories.

Value

Description

Road Utility

Road utility

\subsection{RestorFeatStatus}\label{restorfeatstatus-23}

A description of the status of the individual restoration. It also
includes the party responsible for the individual feature.

Value

Description

USACE Planned

Feature is in a USACE planning stage.

USACE Built

Feature has been built by USACE. If a feature has been planned and built
by USACE, once constructed, it will always be considered USACE built
even if it has been repaired under later project stages.

USACE Enhanced

The function of an existing feature that was not built by USACE has been
enhanced by a feature that was USACE built or repaired.

USACE Repaired

An existing feature that was not built by USACE has been directly
improved or repaired.

Other Planned

A feature that is planned by another organization or government agency.

Other Built

Feature has been built by another organization or government agency and
is an integral part of the UMRR project.

Other Enhanced

The function of an existing feature, that was built by another
organization or government agency, has been enhanced by a feature that
was built or repaired by a non-USACE organization or government agency.

Other Repaired

An existing feature that has been directly improved or repaired by
another organization or government agency.

\subsection{sdsFeatureName}\label{sdsfeaturename-23}

The common name of the feature.

\subsection{sdsFeatureDescription}\label{sdsfeaturedescription-23}

A narrative describing the feature.

\subsection{RestorFeatVdatum}\label{restorfeatvdatum-23}

Vertical datum of elevations recorded in attribute table or used as
feature z values.

Value

Description

MSL 1912

Mean Sea Level of 1912

NGVD 1929

National Geodetic Vertical Datum of 1929

NAVD 1988

North American Vertical Datum of 1988

Other

Other, note datum in sdsFeatureDescription

Unknown

Unknown vertical datum

\subsection{fieldMeasure}\label{fieldmeasure}

Linear reference value or measure (e.g., stationing or chainage)
associated with the road segment, used to relate the segment to field
measurements, plans, or other linear referencing systems.

\subsection{isAnchorSection}\label{isanchorsection}

Boolean indicator identifying whether this segment serves as an anchor
or reference section for the road network (e.g., a control/benchmark
segment used for indexing or linear referencing).

Value

Description

Yes

Segment is an anchor/reference section.

No

Segment is not an anchor/reference section.

\subsection{segmentLength}\label{segmentlength}

Numeric length of the road segment, recorded in the units specified by
segmentLengthUOM.

\subsection{segmentLengthUOM}\label{segmentlengthuom}

Unit of measure for segmentLength.

Value

Description

FT

Feet

M

Meters

KM

Kilometers

MI

Miles

Other

Other unit; describe in sdsFeatureDescription.

Unknown

Unit of measure is not known.

\subsection{roadSegStatus}\label{roadsegstatus}

Operational or physical status of the road segment.

Value

Description

Planned

Road segment is proposed or planned but not yet constructed.

Active

Road segment is constructed and in use.

Limited Use

Road segment is in place but subject to seasonal or restricted use.

Decommissioned

Road segment is no longer maintained or intended for regular use.

Removed

Road segment has been physically removed or reclaimed.

Unknown

Status has not been determined.

\subsection{sdsMetaDataID}\label{sdsmetadataid-23}

References attached metadata.

\subsection{RestorFeatIDPK}\label{restorfeatidpk-23}

Primary key used to uniquely identify the restoration feature record.

\subsection{sdsID}\label{sdsid-23}

Identifier used to link the feature to supporting SDS (spatial data
system) records or related datasets.

\chapter{Sign}\label{sign}

\section{Spatial \& Geometric
Definition}\label{spatial-geometric-definition-24}

\begin{itemize}
\item
  Geometry Type: Point (Sign\_Feature\_Location).
\item
  Spatial Delineation: Point location, orientation azimuth, and
  structural support of physical signage deployed across terrestrial
  paths and aquatic river navigation channels.
\item
  Functional Typology: Regulatory water safety markers, USFWS refuge
  sanctuary warnings, project boundary markers, and public educational
  interpretative kiosks.
\end{itemize}

\subsection{Primary Engineering \& Operational
Objectives}\label{primary-engineering-operational-objectives-5}

\begin{itemize}
\item
  Safety \& Regulatory Enforcement: Mark low-head rock sills, submerged
  weirs, and pumping intakes to protect recreational boaters from marine
  hazards.
\item
  Ecological Protection Enforcement: Delineate restricted-access
  wildlife zones, seasonal ``Slow-No Wake'' corridors, and autumn
  migratory waterfowl resting sanctuaries.
\item
  Public Education: Display ecological maps and interpretative
  narratives detailing the engineering and ecological outcomes of the
  UMRR-EMP restoration project.
\end{itemize}

\subsection{Structural \& Material
Specifications}\label{structural-material-specifications}

\begin{itemize}
\item
  Aquatic Piling Signs: Heavy-duty marine signs mounted on treated
  timber pilings or steel structural pipes driven 10 to 15 ft into the
  river bed, equipped with solar-powered flashing amber LED marine
  lanterns.
\item
  Terrestrial Posts: Breakaway galvanized steel U-channel posts or
  treated timber posts anchored in concrete foundations outside levee
  shoulder lines.
\item
  Sign Face Materials: Heavy-gauge (0.080- to 0.125-in) rust-free
  marine-grade aluminum sheeting faced with ASTM D4956 Type IV or Type
  XI High-Intensity Prismatic reflective sheeting.
\end{itemize}

\subsection{Quality Control, Inspection \& Long-Term
Maintenance}\label{quality-control-inspection-long-term-maintenance-23}

\begin{itemize}
\tightlist
\item
  Visual Audits: Biannual inspections checking for bullet damage, solar
  battery degradation, structural tilt, or flood-debris impact
  displacement.
\end{itemize}

\section{Spatial Feature Attributes}\label{spatial-feature-attributes-7}

\subsection{Definition}\label{definition-23}

A sign is a fixed structure used to convey directional, navigational,
regulatory, safety, interpretive, or other project-related information
to the public or operational users. Sign: refers to any type of sign
used at a project site, including regulatory signs, warning signs, and
informational signs.

\subsection{Purpose}\label{purpose-27}

Signs are installed to clearly communicate important information such as
navigation guidance, safety warnings, access restrictions, ownership or
jurisdiction, and project identification. Properly placed signs support
safe use of project areas, protect project features and habitat, and
ensure compliance with management objectives.

\subsection{Applicability}\label{applicability-24}

Signs are used throughout Upper Mississippi River System (UMRS)
restoration projects wherever clear, visible communication is required
on the landscape or waterway. Typical locations include project access
points, navigation channels, control structures, habitat management
areas, and public use facilities. (UMRR Design Handbook, 2012)

\subsection{Notes}\label{notes-27}

Restoration features now have a Program\_UID field to distinguish
between NESP and UMRR features. Aliases without program identification
will be created for any field that currently has UMRR in the name, so
they can be used for NESP or UMRR data.

Attribute criteria for each field includes the following:

\begin{itemize}
\tightlist
\item
  Attribute Status -- defines whether the field requires a value or if
  it is optional based on available information.
\item
  Phases -- indicates the earliest time in the project lifecycle values
  for this attribute are available.
\item
  Digitizing requirements -- indicates when the field value should be
  entered in the attribute table.
\end{itemize}

Geometry Note: No Z values are needed for this feature.

\subsection{Sign}\label{sign-1}

(point) -- A structure that conveys directional, warning, regulatory,
interpretive, or other information.

\begin{center}\rule{0.5\linewidth}{0.5pt}\end{center}

\section{Feature Attribute
Criteria}\label{feature-attribute-criteria-24}

\#

Attribute

Attribute Status

Phase

Digitizing Requirements

1

Program\_UID

Required

Feasibility

At feature creation

2

PROJECT\_ID (Alias)

Required

Feasibility

At feature creation

3

Type (Alias)

Required

Feasibility

At feature creation

4

Subtype (Alias)

Required

Feasibility

At feature creation

5

RestorFeatStatus

Required

Feasibility

At feature creation

6

sdsFeatureName

Optional

Design

As required and available

7

sdsFeatureDescription

Optional

Design

As required and available

8

RestorFeatVdatum

Optional

Design

As required and available

9

signText

Optional

Design

As required and available

10

signType

Optional

Design

As required and available

11

realPropertySiteUniqueID

Optional

Design

As required and available

12

realPropertyNetworkUniqueID

Optional

Design

As required and available

13

propertyCodeID

Optional

Construction

As required and available

14

sdsMetaDataID

Optional

Construction

As required and available

15

RestorFeatIDPK

Optional

Construction

As required and available

16

sdsID

Optional

Construction

As required and available

17

mediaIDFK

Optional

Construction

As required and available

\begin{center}\rule{0.5\linewidth}{0.5pt}\end{center}

\section{Parameters}\label{parameters-24}

\subsection{Program\_UID}\label{program_uid-24}

Designate program NESP or UMRR.

\subsection{Project\_ID (alias:
PROJECT\_ID)}\label{project_id-alias-project_id-24}

Foreign key to the unique NESP/UMRR project ID.

\subsection{Type (alias)}\label{type-alias-18}

A broad classification used to group NESP/UMRR restoration features into
categories of similar function.

Value

Description

Restoration\_Type

Support feature

\subsection{SubType (alias: Subtype)}\label{subtype-alias-subtype-24}

A classification used to group NESP/UMRR restoration features into
specific design feature categories.

Value

Description

Sign

Sign

\subsection{RestorFeatStatus}\label{restorfeatstatus-24}

A description of the status of the individual restoration. It also
includes the party responsible for the individual feature.

Value

Description

USACE Planned

Feature is in a USACE planning stage.

USACE Built

Feature has been built by USACE. If a feature has been planned and built
by USACE, once constructed, it will always be considered USACE built
even if it has been repaired under later project stages.

USACE Enhanced

The function of an existing feature that was not built by USACE has been
enhanced by a feature that was USACE built or repaired.

USACE Repaired

An existing feature that was not built by USACE has been directly
improved or repaired.

Other Planned

A feature that is planned by another organization or government agency.

Other Built

Feature has been built by another organization or government agency and
is an integral part of the UMRR project.

Other Enhanced

The function of an existing feature, that was built by another
organization or government agency, has been enhanced by a feature that
was built or repaired by a non-USACE organization or government agency.

Other Repaired

An existing feature that has been directly improved or repaired by
another organization or government agency.

\subsection{sdsFeatureName}\label{sdsfeaturename-24}

The common name of the feature.

\subsection{sdsFeatureDescription}\label{sdsfeaturedescription-24}

A narrative describing the feature.

\subsection{RestorFeatVdatum}\label{restorfeatvdatum-24}

Vertical datum of elevations recorded in attribute table or used as
feature z values.

Value

Description

MSL 1912

Mean Sea Level of 1912

NGVD 1929

National Geodetic Vertical Datum of 1929

NAVD 1988

North American Vertical Datum of 1988

Other

Other, note datum in sdsFeatureDescription

Unknown

Unknown vertical datum

\subsection{signText}\label{signtext}

Full text displayed on the sign face, recorded as it appears in the
field (including key abbreviations or messages needed for navigation,
safety, or management).

\subsection{signType}\label{signtype}

General classification of the sign based on its primary function.

Value

Description

Navigational

Sign providing channel alignment, hazard, regulatory, or other
navigation-related information to boaters.

Warning / Safety

Sign warning users of hazards, restricted areas, or other safety
concerns (e.g., ``Danger -- Keep Out'').

Regulatory

Sign stating rules, restrictions, or required actions (e.g., no wake, no
trespassing, speed limits).

Informational / Interpretive

Sign providing general information, education, or interpretation about
the project, resources, or facilities.

Identification

Sign identifying project features, structures, ownership, or facility
names.

Other

Sign type not captured by the above categories; describe in
sdsFeatureDescription.

Unknown

Sign type has not been determined.

\subsection{realPropertySiteUniqueID}\label{realpropertysiteuniqueid-3}

Unique identifier linking the sign to a specific real property site or
installation record in the real property inventory, when applicable.

\subsection{realPropertyNetworkUniqueID}\label{realpropertynetworkuniqueid-3}

Unique identifier linking the sign to a network-level real property or
asset management record (e.g., a regional or enterprise real property
system), when applicable.

\subsection{propertyCodeID}\label{propertycodeid-3}

Code indicating the real property classification or category of the
asset (e.g., sign, marker, or similar facility type) as defined by the
managing agency's real property coding system.

\subsection{sdsMetaDataID}\label{sdsmetadataid-24}

References attached metadata.

\subsection{RestorFeatIDPK}\label{restorfeatidpk-24}

Primary key used to uniquely identify the restoration feature record.

\subsection{sdsID}\label{sdsid-24}

Identifier used to link the feature to supporting SDS (spatial data
system) records or related datasets.

\subsection{mediaIDFK}\label{mediaidfk-22}

Foreign key used to reference associated media (e.g., photographs,
drawings, as-built plans) documenting the sign.

\chapter{Water Utility Node}\label{water-utility-node}

\section{Spatial \& Geometric
Definition}\label{spatial-geometric-definition-25}

\begin{itemize}
\item
  Geometry Type: Point (Water\_Utility\_Node).
\item
  Spatial Delineation: Geographic point locations representing discrete
  hydraulic fixtures, fittings, and mechanical control nodes within
  project water delivery systems.
\item
  Component Typology: Gate valves, butterfly valves, check valves,
  vacuum/air-relief valves, cleanouts, hydrants, flow meters, and
  manifold tee-junctions.
\end{itemize}

\subsection{Primary Engineering \& Hydraulic
Objectives}\label{primary-engineering-hydraulic-objectives}

\begin{itemize}
\item
  Pipeline Flow Isolation: Enable operators to isolate, balance, or
  throttle flow rates across individual pipe legs within complex
  multi-cell wetland delivery systems.
\item
  Surge \& Vacuum Protection: Air/vacuum relief valves release entrapped
  air pockets during pipeline filling and introduce air during rapid
  draining to prevent catastrophic vacuum pipeline collapse.
\item
  Volume Flow Metering: Ultrasonic or magnetic flow meters recording
  pumped water volumes for water-budget accounting and regulatory
  reporting.
\end{itemize}

\subsection{Technical \& Mechanical
Specifications}\label{technical-mechanical-specifications}

\begin{itemize}
\item
  Valving Specifications: Resilient-seated wedge gate valves or
  high-performance rubber-lined butterfly valves rated for minimum 150
  psi working pressure.
\item
  Access Vaults: Underground nodes housed inside precast reinforced
  concrete utility vaults equipped with aluminum diamond-plate hatch
  doors (H−20 traffic-rated), interior drain sumps, and access ladders.
\item
  Corrosion Protection: All valves and fittings epoxy-coated per AWWA
  C550, with stainless steel assembly bolts and sacrificial zinc anode
  bags in corrosive alluvial soils.
\end{itemize}

\subsection{Quality Control, Inspection \& Long-Term
Maintenance}\label{quality-control-inspection-long-term-maintenance-24}

\begin{itemize}
\tightlist
\item
  Valve Exercising Programs: Semi-annual operational cycling of all
  buried and vault-housed valves through their complete open/close
  travel to prevent sediment binding.
\end{itemize}

\section{Spatial Feature Attributes}\label{spatial-feature-attributes-8}

\subsection{Definition}\label{definition-24}

Water Utility Node: Represents a point in a water distribution system,
such as a valve, hydrant, or meter. A subdivision of a water
distribution network, particularly an asset that participates in the
transmission of water but that is not a pipeline. For example, a pump,
valve or fitting.

\subsection{Purpose}\label{purpose-28}

Spatial Delineation: Geographic point locations representing discrete
hydraulic fixtures, fittings, and mechanical control nodes within
project water delivery systems.

\subsection{Applicability}\label{applicability-25}

Primary Engineering \& Hydraulic Objectives • Pipeline Flow Isolation:
Enable operators to isolate, balance, or throttle flow rates across
individual pipe legs within complex multi-cell wetland delivery systems.
• Surge \& Vacuum Protection: Air/vacuum relief valves release entrapped
air pockets during pipeline filling and introduce air during rapid
draining to prevent catastrophic vacuum pipeline collapse. • Volume Flow
Metering: Ultrasonic or magnetic flow meters recording pumped water
volumes for water-budget accounting and regulatory reporting.

\subsection{Notes}\label{notes-28}

Restoration features now have a Program\_UID field to distinguish
between NESP and UMRR features. Aliases without program identification
will be created for any field that currently has UMRR in the name, so
they can be used for NESP or UMRR data.

Attribute criteria for each field includes the following:

\begin{itemize}
\tightlist
\item
  Attribute Status -- defines whether the field requires a value or if
  it is optional based on available information.
\item
  Phases -- indicates the earliest time in the project lifecycle values
  for this attribute as available.
\item
  Digitizing requirements -- indicates when the field value should be
  entered in the attribute table.
\end{itemize}

Geometry Note: No Z values are needed for this feature.

\subsection{WaterUtilityNode}\label{waterutilitynode}

(point) -- A subdivision of a water distribution network, particularly
an asset that participates in the transmission of water but that is not
a pipeline. For example, a pump, valve or fitting.

\section{Feature Attribute
Criteria}\label{feature-attribute-criteria-25}

\#

Attribute

Attribute Status

Phase

Digitizing Requirements

1

Program\_UID

Required

Feasibility

At feature creation

2

PROJECT\_ID

Required

Feasibility

At feature creation

3

Type

Required

Feasibility

At feature creation

4

Subtype

Required

Feasibility

At feature creation

5

RestorFeatStatus

Required

Feasibility

At feature creation

6

sdsFeatureName

Optional

Design

As required and available

7

sdsFeatureDescription

Optional

Design

As required and available

8

RestorFeatVdatum

Optional

Design

As required and available

9

waterNetworkSubtype

Optional

Design

As required and available

10

OperationStatus

Optional

Design

As required and available

11

realPropertySiteUniqueID

Optional

Design

As required and available

12

realPropertyNetworkUniqueID

Optional

Design

As required and available

13

realPropertyUniqueID

Optional

Construction

As required and available

14

propertyCodeID

Optional

Construction

As required and available

15

waterNodeType

Optional

Construction

As required and available

16

sdsMetaDataID

Optional

Construction

As required and available

17

RestorFeatIDPK

Optional

Construction

As required and available

18

sdsID

Optional

Construction

As required and available

19

mediaIDFK

Optional

Construction

As required and available

\section{Parameters}\label{parameters-25}

\subsection{Program\_UID}\label{program_uid-25}

Designate program NESP or UMRR.

\subsection{Project\_ID (alias:
PROJECT\_ID)}\label{project_id-alias-project_id-25}

Foreign key to the unique NESP/UMRR project ID.

\subsection{Type (alias)}\label{type-alias-19}

A broad classification used to group NESP/UMRR restoration features into
categories of similar function.

Value

Description

Restoration\_Type

Support feature

\subsection{SubType (alias: Subtype)}\label{subtype-alias-subtype-25}

A classification used to group NESP/UMRR restoration features into
specific design feature categories.

Value

Description

Water Utility

Water utility

\subsection{RestorFeatStatus}\label{restorfeatstatus-25}

A description of the status of the individual restoration. It also
includes the party responsible for the individual feature.

Value

Description

USACE Planned

Feature is in a USACE planning stage.

USACE Built

Feature has been built by USACE. If a feature has been planned and built
by USACE, once constructed, will always be considered USACE built even
if it has been repaired under later project stages.

USACE Enhanced

The function of an existing feature that was not built by USACE has been
enhanced by a feature that was USACE built or repaired.

USACE Repaired

An existing feature that was not built by USACE has been directly
improved or repaired.

Other Planned

A feature that is planned by another organization or government agency.

Other Built

Feature has been built by another organization or government agency and
is an integral part of the UMRR project.

Other Enhanced

The function of an existing feature, that was built by another
organization or government agency, has been enhanced by a feature that
was built or repaired by a non-USACE organization or government agency.

Other Repaired

An existing feature that has been directly improved or repaired by
another organization or government agency.

\subsection{sdsFeatureName}\label{sdsfeaturename-25}

The common name of the feature.

\subsection{sdsFeatureDescription}\label{sdsfeaturedescription-25}

A narrative describing the feature.

\subsection{RestorFeatVdatum}\label{restorfeatvdatum-25}

Vertical datum of elevations recorded in attribute table or used as
feature z values.

Value

Description

MSL 1912

Mean Sea Level of 1912

NGVD 1929

National Geodetic Vertical Datum of 1929

NAVD 1988

North American Vertical Datum of 1988

Other

Other, note datum in sdsFeatureDescription

Unknown

Unknown vertical datum

\subsection{waterNetworkSubtype}\label{waternetworksubtype}

(Need text)

\subsection{OperationStatus}\label{operationstatus-3}

(Need text)

\subsection{realPropertySiteUniqueID}\label{realpropertysiteuniqueid-4}

(Need text)

\subsection{realPropertyNetworkUniqueID}\label{realpropertynetworkuniqueid-4}

(Need text)

\subsection{realPropertyUniqueID}\label{realpropertyuniqueid}

(Need text)

\subsection{propertyCodeID}\label{propertycodeid-4}

(Need text)

\subsection{waterNodeType}\label{waternodetype}

(Need text)

\subsection{sdsMetaDataID}\label{sdsmetadataid-25}

References attached metadata.

\subsection{RestorFeatIDPK}\label{restorfeatidpk-25}

References attached media.

\subsection{sdsID}\label{sdsid-25}

References attached media.

\subsection{mediaIDFK}\label{mediaidfk-23}

References attached media.

\chapter{Water Utility Segment}\label{water-utility-segment}

\section{Spatial \& Geometric
Definition}\label{spatial-geometric-definition-26}

\begin{itemize}
\item
  Geometry Type: Polyline (Water\_Utility\_Segment).
\item
  Spatial Delineation: Linear paths of pressurized supply pipes, gravity
  discharge lines, siphon tubes, and dewatering conduits linking pump
  stations, wells, and wetland cells.
\item
  Material Typology: High-density polyethylene (HDPE) DR-11/DR-17 pipe,
  ductile iron pipe (DIP), and reinforced concrete cylinder pipe (RCCP).
\end{itemize}

\subsection{Primary Engineering \& Hydraulic
Objectives}\label{primary-engineering-hydraulic-objectives-1}

\begin{itemize}
\item
  Bulk Water Conveyance: Transport large volumes of water (5 to 50+ cfs)
  between river sources, wells, and isolated interior moist soil units
  across long distances.
\item
  Pressure Containment: Withstand sustained pumping operating heads,
  transient hydraulic water hammer surges, and external earth/vehicle
  loadings.
\item
  Embankment Integrity: Provide closed, continuous conveyance across or
  through earthen levees without inducing structural piping or seepage
  along pipe outer walls.
\end{itemize}

\subsection{Technical Specifications \& Hydraulic
Sizing}\label{technical-specifications-hydraulic-sizing}

\begin{itemize}
\item
  Pipe Diameters: Sized between 12 and 48 in inside diameter based on
  fluid mechanics equations (Q=V⋅A), restricting internal velocities to
  4.0 to 8.0 ft/s to minimize dynamic head loss.
\item
  Pipe Jointing: HDPE pipelines joined using continuous thermal
  butt-fusion welding per ASTM F2620, producing zero-leakage, fully
  restrained joint strings.
\item
  Anti-Seepage Measures: Where pipes penetrate through earthen levees,
  concrete cutoff collars, bentonite seepage rings, and impervious clay
  backfill are placed along the barrel.
\item
  Siphon Configurations: For over-the-levee pipelines, pipe crowns cross
  above the levee crest and feature vacuum-breaker valves at the high
  point to terminate siphon action.
\end{itemize}

\subsection{Quality Control, Inspection \& Long-Term
Maintenance}\label{quality-control-inspection-long-term-maintenance-25}

\begin{itemize}
\item
  Hydrostatic Pressure Testing: Pressure testing conducted per AWWA
  standards (e.g., testing at 1.5× working design pressure for 2 hours)
  prior to trench backfill.
\item
  Internal CCTV Video Inspections: Closed-circuit robotic crawler
  cameras run through pipe barrels at 5-year intervals to verify the
  absence of joint displacement, sediment plugs, or internal scaling.
\end{itemize}

\section{Spatial Feature Attributes}\label{spatial-feature-attributes-9}

\subsection{Definition}\label{definition-25}

Water Utility Segment: Represents a linear segment of a water
distribution system, such as a pipeline or ditch. A subdivision of a
wastewater network, particularly a pipeline for the transport of
stormwater, sewage or industrial waste between the source, holding
facilities, and/or treatment facilities.

\subsection{Purpose}\label{purpose-29}

Spatial Delineation: Linear paths of pressurized supply pipes, gravity
discharge lines, siphon tubes, and dewatering conduits linking pump
stations, wells, and wetland cells.

\subsection{Applicability}\label{applicability-26}

Primary Engineering \& Hydraulic Objectives • Bulk Water Conveyance:
Transport large volumes of water () between river sources, wells, and
isolated interior moist soil units across long distances. • Pressure
Containment: Withstand sustained pumping operating heads, transient
hydraulic water hammer surges, and external earth/vehicle loadings. •
Embankment Integrity: Provide closed, continuous conveyance across or
through earthen levees without inducing structural piping or seepage
along pipe outer walls.

\subsection{Notes}\label{notes-29}

Restoration features now have a Program\_UID field to distinguish
between NESP and UMRR features. Aliases without program identification
will be created for any field that currently has UMRR in the name, so
they can be used for NESP or UMRR data.

Attribute criteria for each field includes the following:

\begin{itemize}
\tightlist
\item
  Attribute Status -- defines whether the field requires a value or if
  it is optional based on available information.
\item
  Phases -- indicates the earliest time in the project lifecycle values
  for this attribute as available.
\item
  Digitizing requirements -- indicates when the field value should be
  entered in the attribute table. Geometry Note: No Z values are needed
  for this feature.
\end{itemize}

\subsection{WaterUtilitySegment}\label{waterutilitysegment}

(line) -- A subdivision of a wastewater network, particularly a pipeline
for the transport of stormwater, sewage or industrial waste between the
source, holding facilities, and/or treatment facilities.

\section{Feature Attribute
Criteria}\label{feature-attribute-criteria-26}

\#

Attribute

Attribute Status

Phase

Digitizing Requirements

1

Program\_UID

Required

Feasibility

At feature creation

2

PROJECT\_ID

Required

Feasibility

At feature creation

3

Type

Required

Feasibility

At feature creation

4

Subtype

Required

Feasibility

At feature creation

5

RestorFeatStatus

Required

Feasibility

At feature creation

6

sdsFeatureName

Optional

Design

As required and available

7

sdsFeatureDescription

Optional

Design

As required and available

8

RestorFeatVdatum

Optional

Design

As required and available

9

waterNetworkSubtype

Optional

Design

As required and available

10

OperationStatus

Optional

Design

As required and available

11

realPropertySiteUniqueID

Optional

Design

As required and available

12

realPropertyNetworkUniqueID

Optional

Design

As required and available

13

realPropertyUniqueID

Optional

Construction

As required and available

14

propertyCodeID

Optional

Construction

As required and available

15

waterSegmentType

Optional

Construction

As required and available

16

sdsMetaDataID

Optional

Construction

As required and available

17

RestorFeatIDPK

Optional

Construction

As required and available

18

sdsID

Optional

Construction

As required and available

19

mediaIDFK

Optional

Construction

As required and available

\section{Parameters}\label{parameters-26}

\subsection{Program\_UID}\label{program_uid-26}

Designate program NESP or UMRR.

\subsection{Project\_ID (alias:
PROJECT\_ID)}\label{project_id-alias-project_id-26}

Foreign key to the unique NESP/UMRR project ID.

\subsection{Type (alias)}\label{type-alias-20}

A broad classification used to group NESP/UMRR restoration features into
categories of similar function.

Value

Description

Restoration\_Type

Support feature

\subsection{SubType (alias: Subtype)}\label{subtype-alias-subtype-26}

A classification used to group NESP/UMRR restoration features into
specific design feature categories.

Value

Description

Water Utility

Water utility

\subsection{RestorFeatStatus}\label{restorfeatstatus-26}

A description of the status of the individual restoration. It also
includes the party responsible for the individual feature.

Value

Description

USACE Planned

Feature is in a USACE planning stage.

USACE Built

Feature has been built by USACE. If a feature has been planned and built
by USACE, once constructed, it will always be considered USACE built
even if it has been repaired under later project stages.

USACE Enhanced

The function of an existing feature that was not built by USACE has been
enhanced by a feature that was USACE built or repaired.

USACE Repaired

An existing feature that was not built by USACE has been directly
improved or repaired.

Other Planned

A feature that is planned by another organization or government agency.

Other Built

Feature has been built by another organization or government agency and
is an integral part of the UMRR project.

Other Enhanced

The function of an existing feature, that was built by another
organization or government agency, has been enhanced by a feature that
was built or repaired by a non-USACE organization or government agency.

Other Repaired

An existing feature that has been directly improved or repaired by
another organization or government agency.

\subsection{sdsFeatureName}\label{sdsfeaturename-26}

The common name of the feature.

\subsection{sdsFeatureDescription}\label{sdsfeaturedescription-26}

A narrative describing the feature.

\subsection{RestorFeatVdatum}\label{restorfeatvdatum-26}

Vertical datum of elevations recorded in attribute table or used as
feature z values.

Value

Description

MSL 1912

Mean Sea Level of 1912

NGVD 1929

National Geodetic Vertical Datum of 1929

NAVD 1988

North American Vertical Datum of 1988

Other

Other, note datum in sdsFeatureDescription

Unknown

Unknown vertical datum

\subsection{waterNetworkSubtype}\label{waternetworksubtype-1}

Subtype describing the specific role or classification of the
water-related network feature within the restoration project (e.g.,
distribution line, transmission main, service connection). Used to
further categorize water utility network components.

\subsection{OperationStatus}\label{operationstatus-4}

Indicates the current operational state of the feature within the water
network (e.g., Active, Inactive, Under Construction, Out of Service).
Used to distinguish features that are in-service from those that are
planned, decommissioned, or temporarily unavailable.

\subsection{realPropertySiteUniqueID}\label{realpropertysiteuniqueid-5}

Unique identifier for the specific real property site (parcel or
facility) associated with the feature. Used to link the feature to
site-level real property records.

\subsection{realPropertyNetworkUniqueID}\label{realpropertynetworkuniqueid-5}

Unique identifier for the real property network or system (e.g., a
defined utility or infrastructure network) that the feature belongs to.
Used to associate individual features with a larger network asset
record.

\subsection{realPropertyUniqueID}\label{realpropertyuniqueid-1}

Unique identifier for the individual real property asset represented by
the feature. Used to cross-reference the feature with the authoritative
real property inventory or asset management system.

\subsection{propertyCodeID}\label{propertycodeid-5}

Code or identifier representing the property classification or type
(e.g., building, utility, easement, parcel). Used for grouping and
reporting on property types consistent with real property management
standards.

\subsection{waterSegmentType}\label{watersegmenttype}

Classification describing the type of water segment represented by the
feature (e.g., main channel, side channel, backwater area, intake,
outfall). Used to distinguish functional waterbody segments within the
project area.

\subsection{sdsMetaDataID}\label{sdsmetadataid-26}

References attached metadata.

\subsection{RestorFeatIDPK}\label{restorfeatidpk-26}

References attached media.

\subsection{sdsID}\label{sdsid-26}

References attached media.

\subsection{mediaIDFK}\label{mediaidfk-24}

References attached media.

\part{Restoration Feature Crosswalk}

\chapter{Restoration Features Crosswalk to Navigation and Ecosystem
Restoration Program
Features}\label{restoration-features-crosswalk-to-navigation-and-ecosystem-restoration-program-features}

\section{Introduction}\label{introduction-1}

The following provides restoration features alignment to the Navigation
and Ecosystem Restoration Program (NESP) measures.

\section{NESP Restoration Measures}\label{nesp-restoration-measures}

\subsection{Island Building}\label{island-building}

Includes the following Restoration Features:

Island (line and polygon) Dredging Notches Riverine Flow Structures
(line and polygon) Wetlands

\subsection{Floodplain Restoration}\label{floodplain-restoration}

Includes the following Restoration Features:

Dredging Embankments (line and polygon) Moist Soil Management Units
Flora Planting and Seeding Wells

\subsection{Water Level Management}\label{water-level-management}

Includes the following Restoration Features:

Dredging Embankments (line and polygon) Water Control Structures Fish
Structures Pump Station

\subsection{Backwater Restoration}\label{backwater-restoration}

Includes the following Restoration Features:

Dredging Riverine Flow Structures (line and polygon) Fishing Structures
Notches Revetments

\subsection{Side Channel Restoration}\label{side-channel-restoration}

Includes the following Restoration Features:

Dredging Wetlands Flora Planting and Seeding Dredging Revetments

\subsection{Wing Dam / Dike
Alterations}\label{wing-dam-dike-alterations}

Includes the following Restoration Features:

Revetments Dredging Riverine Flow Structures (line and polygon)

\subsection{Island Protection}\label{island-protection}

Includes the following Restoration Features:

Dredging Revetments Riverine Flow Structures (line and polygon)

\subsection{Shoreline Protection}\label{shoreline-protection}

Includes the following Restoration Features:

Island (line and polygon) Dredging Notches Riverine Flow Structures
(line and polygon) Wetlands Flora Planting and Seeding

\subsection{Topographic Diversity}\label{topographic-diversity}

Includes the following Restoration Features:

Island (line and polygon) Dredging Notches Riverine Flow Structures
(line and polygon) Wetlands

\subsection{Dam Point Control}\label{dam-point-control}

Includes the following Restoration Features:

????

\subsection{Dam Embankment Lowering}\label{dam-embankment-lowering}

Includes the following Restoration Features:

Dredging Notches Riverine Flow Structures (line and polygon)

\subsection{Reduced Water Level
Fluctuation}\label{reduced-water-level-fluctuation}

Includes the following Restoration Features:

Pump Station Dredging Notches Riverine Flow Structures (line and
polygon) Wetlands

\subsection{Forest Management}\label{forest-management}

Includes the following Restoration Features:

Flora Planting and Seeding Riverine Flow Structures (line and polygon)
Wetlands

\chapter{Restoration Features Crosswalk to Environmental Design Handbook
Chapters}\label{restoration-features-crosswalk-to-environmental-design-handbook-chapters}

Environmental Design Handbook, 2012 documents the array of restoration
tools and lessons learned to aid in the design of future Habitat
Rehabilitation and Enhancement Projects (HREP). These restoration tools
include shoreline protection, island creation, water level management,
backwater dredging, secondary channel modifications and river training
structures, aeration, and floodplain and tributary restoration. These
tools are captured to a GIS database as Restoration Features for
mapping, analysis for feature impact, and future planning.

The following provides restoration features alignment to the chapters of
the design handbook.

\section{Chapter 4}\label{chapter-4}

\subsection{Shoreline and Riverbank
Protection}\label{shoreline-and-riverbank-protection}

Includes the following Restoration Features:

Revetments Riverine Flow Structures (line and polygon) Embankments (line
and polygon)

\section{Chapter 5}\label{chapter-5}

\subsection{Water Level Management}\label{water-level-management-1}

Includes the following Restoration Features:

Dredging Fish Structures Embankments (line and polygon) Moist Soil
Management Units Pump Stations Wells

\section{Chapter 6}\label{chapter-6}

\subsection{Dredging}\label{dredging}

Includes the following Restoration Features:

Dredging Embankments (line and polygon) Water Control Structures

\section{Chapter 7}\label{chapter-7}

\subsection{River Training Structures and Secondary Channel
Modification}\label{river-training-structures-and-secondary-channel-modification}

Includes the following Restoration Features:

Riverine Flow Structures (line and polygon) Fishing Structures Notches
Revetments

\section{Chapter 8}\label{chapter-8}

\subsection{Floodplain Restoration}\label{floodplain-restoration-1}

Includes the following Restoration Features:

Wetlands Flora Planting and Seeding Dredging

\section{Chapter 9}\label{chapter-9}

\subsection{Island Design}\label{island-design}

Includes the following Restoration Features:

Island (line and polygon) Dredging Notches Riverine Flow Structures
(line and polygon) Wetlands

\bookmarksetup{startatroot}

\chapter*{Resources}\label{resources}
\addcontentsline{toc}{chapter}{Resources}

\markboth{Resources}{Resources}

\bookmarksetup{startatroot}

\chapter*{References}\label{references}
\addcontentsline{toc}{chapter}{References}

\markboth{References}{References}

\protect\phantomsection\label{refs}




\end{document}
